\documentclass[11pt]{article}

%================================================
% Packages
%================================================
\usepackage[margin=1in]{geometry}
\usepackage{setspace}
\usepackage{amsmath,amssymb,amsthm,mathtools,bm,mathrsfs}
\usepackage{enumitem}
\usepackage{microtype}
\usepackage{graphicx}
\usepackage{natbib}
\usepackage{array,booktabs,longtable,threeparttable,makecell}
\usepackage{siunitx}
\usepackage{placeins}
\usepackage{xcolor}
\usepackage[colorlinks=true,linkcolor=blue,citecolor=blue,urlcolor=blue]{hyperref}
\usepackage[nameinlink,capitalise]{cleveref}

\onehalfspacing
\emergencystretch=6em
\numberwithin{equation}{section}

%================================================
% Theorem environments
%================================================
\newtheorem{assumption}{Assumption}
\newtheorem{theorem}{Theorem}
\newtheorem{proposition}{Proposition}
\newtheorem{lemma}{Lemma}
\newtheorem{corollary}{Corollary}
\theoremstyle{remark}
\newtheorem{remark}{Remark}
\crefname{assumption}{Assumption}{Assumptions}
\Crefname{assumption}{Assumption}{Assumptions}
\crefname{theorem}{Theorem}{Theorems}
\Crefname{theorem}{Theorem}{Theorems}
\crefname{proposition}{Proposition}{Propositions}
\Crefname{proposition}{Proposition}{Propositions}
\crefname{lemma}{Lemma}{Lemmas}
\Crefname{lemma}{Lemma}{Lemmas}
\crefname{corollary}{Corollary}{Corollaries}
\Crefname{corollary}{Corollary}{Corollaries}
\crefname{remark}{Remark}{Remarks}
\Crefname{remark}{Remark}{Remarks}
\crefname{section}{Section}{Sections}
\Crefname{section}{Section}{Sections}
\crefname{table}{Table}{Tables}
\Crefname{table}{Table}{Tables}
\crefname{figure}{Figure}{Figures}
\Crefname{figure}{Figure}{Figures}

%================================================
% Notation
%================================================
\newcommand{\R}{\mathbb R}
\newcommand{\Ezero}{\mathbb E_0}
\newcommand{\Pzero}{\mathbb P_0}
\newcommand{\rank}{\operatorname{rank}}
\newcommand{\op}{\operatorname{op}}
\newcommand{\argmin}{\operatorname*{arg\,min}}
\newcommand{\ip}[2]{\left\langle #1,#2\right\rangle}
\newcommand{\be}{\bm e}
\newcommand{\bq}{\bm q}
\newcommand{\bu}{\bm u}
\newcommand{\br}{\bm r}
\newcommand{\bs}{\bm s}
\newcommand{\bA}{\bm A}
\newcommand{\bB}{\bm B}
\newcommand{\bC}{\bm C}
\newcommand{\bD}{\bm D}
\newcommand{\bH}{\bm H}
\newcommand{\bL}{\bm L}
\newcommand{\bM}{\bm M}
\newcommand{\bQ}{\bm Q}
\newcommand{\bR}{\bm R}
\newcommand{\bU}{\bm U}
\newcommand{\bV}{\bm V}
\newcommand{\bDelta}{\bm\Delta}
\newcommand{\bTheta}{\bm\Theta}
\newcommand{\bSigma}{\bm\Sigma}
\newcommand{\bPhi}{\bm\Phi}
\newcommand{\calC}{\mathcal C}
\newcommand{\calD}{\mathcal D}
\newcommand{\calM}{\mathcal M}
\newcommand{\calP}{\mathcal P}
\newcommand{\calR}{\mathcal R}
\newcommand{\calT}{\mathcal T}
\newcommand{\doilink}[1]{\href{https://doi.org/#1}{doi:\ #1}}

%================================================
% Placeholders for tables and figures
%================================================
\newcommand{\placeholdertable}[3]{%
\begin{table}[!htbp]
\centering
\caption{#1}
\label{#2}
\begin{tabular}{p{0.88\textwidth}}
\toprule
\textit{Placeholder.} #3\\
\bottomrule
\end{tabular}
\end{table}}

\newcommand{\placeholderfigure}[3]{%
\begin{figure}[!htbp]
\centering
\fbox{\begin{minipage}[c][2.2in][c]{0.86\textwidth}\centering #3\end{minipage}}
\caption{#1}
\label{#2}
\end{figure}}

% Full-width notes for Monte Carlo tables
\newcommand{\mctablenote}[1]{%
  \par\vspace{3pt}
  \noindent\begin{minipage}{\linewidth}
  \footnotesize
  \raggedright
  \textit{Notes:} #1
  \end{minipage}
}

%================================================
\begin{document}

\title{Debiased Inference for Dynamic Panels with
Low-Rank \\ Coefficient Heterogeneity}
\author{}
\date{}
\maketitle

\begin{abstract}
We study dynamic panel models in which the coefficients on lagged outcomes and observed covariates vary across units and time. Each autoregressive coefficient matrix, each slope coefficient matrix, and the interactive-effects matrix has its own fixed low rank; these matrices need not share ranks, loadings, or factors. Rank selection uses one joint least-squares pilot whose rank bound for each coefficient matrix is one above its largest reportable rank. The pilot singular values are normalized by observable regressor scales, and each matrix rank is selected by a deterministic ridge ratio with ridge $a_{NT}=1/\log(NT)$. The selected rank vector is then imposed in one unpenalized joint fixed-rank least-squares fit. All reported coefficients and inferential objects are constructed from this final fit, not from the pilot. A target-specific Riesz representer accounts for estimation error in every coefficient matrix. Individual coefficients and suitable fixed-time means and contrasts use the full-panel plug-in estimator, whereas a fixed two-way split-panel correction removes the leading loading- and factor-estimation biases for averages over many units and periods. The probability theory is conditional on common shocks and allows dynamic feedback through lagged outcomes, predetermined covariates, conditional heteroskedasticity, spatial dependence, separate coefficient-matrix ranks, and different growth rates for $N$ and $T$. A spatial HAC estimator consistently estimates the conditional variance. Rank selection is performed once on the full panel; four fixed-rank split-panel fits are added only when a broad target requires bias correction.
\end{abstract}

\noindent\textbf{Keywords:} dynamic panel data; heterogeneous coefficients; low-rank matrices; interactive effects; full-panel inference; rank selection; spatial dependence; near-epoch dependence.

\clearpage
\tableofcontents
\clearpage

%================================================
\section{Introduction}
\label{sec:introduction}

Dynamic panel studies often ask whether persistence or covariate effects differ across units, over time, or between economically meaningful groups. A common-coefficient model cannot answer these questions, while unrestricted unit--time coefficients cannot be estimated from a single panel. We use separate low-rank restrictions to make the complete coefficient matrices estimable without imposing known groups, smooth time paths, or a common latent structure across regressors.

The model contains one coefficient matrix for each lagged outcome and each observed covariate. A further low-rank interactive-effects matrix absorbs additive latent conditional-mean heterogeneity. Every matrix has its own rank and singular spaces. This separation is useful when persistence, covariate effects, and common movements need not be generated by the same factors.

Inference is difficult for three reasons. First, all coefficient matrices enter one outcome equation jointly, so identification must prevent changes in different matrices from offsetting one another. Second, a current disturbance enters future observations through lagged outcomes, while shocks to nearby units may be correlated. Third, estimating loading and factor rows from the same panel creates biases of order $T^{-1}$ and $N^{-1}$. These biases are negligible for an individual coefficient but can be first order for an average over many units and periods.

The estimator has two stages. First, one joint spectral pilot estimates the singular spectrum of each coefficient matrix. Observable regressor scales put these spectra in comparable fitted-value units, and a deterministic ridge-ratio rule selects each matrix rank separately. Second, the selected rank vector is imposed in one unpenalized joint fixed-rank least-squares fit. The pilot is used for rank determination only; it is never the reported coefficient estimator. For a parameter of interest, a Riesz representer solves the least-squares Gram equation on the estimated fixed-rank tangent space. This inverse-Gram adjustment translates estimation errors in all autoregressive, slope, and interactive-effect matrices into first-order uncertainty for the target. At an exact interior least-squares solution, the residual-weighted Riesz sum is zero, so the Riesz weights determine the influence expansion and standard error rather than an additional point-estimate correction.

Individual coefficients, fixed unit--time lag sums, and suitable fixed-time means and contrasts use the full-panel plug-in estimator because their $N^{-1}+T^{-1}$ estimation bias is smaller than the standard error. For averages whose weights extend over many units and periods, the standard error can be of order $(NT)^{-1/2}$, so the same bias is not negligible. A fixed two-way split-panel correction removes the leading $T^{-1}$ loading bias and $N^{-1}$ factor bias. It uses two time-split fits and two unit-split fits, with the selected full-panel ranks held fixed.

The stochastic assumptions are stated conditionally on common shocks. The common shocks may be strong, persistent, and relevant for covariates, conditional variances, and spatial dependence. Conditional on them, the remaining unit-specific innovations are weakly dependent over space and time. Predetermined covariates are approximable by recent innovations from nearby units, while dynamic stability controls how current disturbances affect distant future outcomes.

The main contribution is a joint inferential framework for several separately ranked coefficient matrices in a dynamic panel. The method combines coefficient-matrix rank selection from one spectral pilot, unpenalized joint least-squares estimation at the selected ranks, explicit alignment and expansion of the estimated singular spaces, and a target-specific Riesz equation that accounts for estimation spillovers across all coefficient matrices. The rank-consistency result uses the same strong-signal and pilot-recovery rates as the fixed-rank analysis and allows $N$ and $T$ to grow at different rates. The resulting theory provides full-panel inference without random folds for individual coefficients and suitable fixed-time means and contrasts, a fixed split-panel correction for broad two-way averages, and conditional spatial inference under common shocks. We assume exact fixed ranks and consistent selection of the separate rank vector; approximate-low-rank inference, growing-rank asymptotics, semiparametric efficiency, and inference robust to rank overestimation are not claimed.

The remainder is organized as follows. \Cref{sec:literature} relates the paper to the literature. \Cref{sec:model} presents the model and parameters of interest. \Cref{sec:estimation} defines estimation, rank selection, and variance estimation. \Cref{sec:assumptions} states the assumptions. \Cref{sec:theory} gives the asymptotic results. \Cref{sec:implementation} describes computation. \Cref{sec:MC} and \Cref{sec:application} present the Monte Carlo and empirical designs. The appendix follows the theorem order.

%================================================
\section{Related Literature}
\label{sec:literature}

The paper relates to interactive-effects panels, dynamic panels, heterogeneous coefficients, low-rank estimation, rank determination, and inference under spatial and temporal dependence.

Large-panel factor methods and interactive-effects regressions use latent factors to absorb unobserved common heterogeneity \citep{bai2002,bai2003,pesaran2006,moonweidner2015}. We retain a low-rank interactive-effects matrix but allow every autoregressive and slope coefficient matrix to vary over units and time. Each coefficient matrix has its own rank and singular spaces; no common loading or factor space is imposed across matrices.

The classical dynamic-panel literature uses instruments or transformations to address short panels with common coefficients \citep{andersonhsiao1981,andersonhsiao1982,arellanobond1991}. Here both panel dimensions are large, the autoregressive coefficients are heterogeneous, and least squares relies on conditional sequential exogeneity. In dynamic interactive-effects models, predetermined regressors and heteroskedastic or correlated idiosyncratic errors generate incidental-parameter biases \citep{moonweidner2017}. The loading-side and factor-side biases below have the same economic sources, but they arise for target functionals of several separately ranked coefficient matrices under conditional spatial dependence. If lagged outcomes remain endogenous after conditioning on current covariates, past panel information, and common shocks, the least-squares criterion would need to be replaced by an instrumented one.

Random-coefficient, mean-group, and grouped-pattern methods impose different structures on coefficient heterogeneity \citep{swamy1970,pesaransmith1995,bonhommemanresa2015}. Low rank instead represents variation through latent unit and time dimensions. It supports individual coefficients, cross-sectional and time averages, group contrasts, and lag-sum persistence targets without a known group assignment.

The closest low-rank inferential benchmark is \citet{chernozhukov2023}. They study linear functionals of one matrix that is well approximated by a spiked low-rank structure and combine regularized initialization, unpenalized least-squares refinement, and explicit rotation analysis. Under their model, the procedure attains a semiparametric efficiency bound. The present paper shares the emphasis on unpenalized refinement and singular-space alignment, but the estimator and proof must be joint across several coefficient matrices that multiply different stochastic regressors, including a lagged dependent variable. Conditional spatial dependence also changes the score and variance arguments.

The target-specific Riesz equation represents the derivative of a selected parameter under the joint least-squares Gram form and therefore accounts simultaneously for estimation errors in every coefficient matrix. This inverse-Gram adjustment is related in purpose to the rotation and least-squares refinements in \citet{chernozhukov2023}, but the present paper establishes validity rather than semiparametric efficiency. Efficiency under conditional heteroskedasticity and spatial dependence would require covariance-optimal weighting of the joint Gram equation.

\citet{choi2024} obtain low-rank completion inference without implemented sample splitting by combining regularized initialization with loading and factor refinements. Their proof uses an auxiliary leave-one-out estimator that is not computed. Our full-panel proof follows the same broad device, but it deletes a space--time neighborhood because a current disturbance reappears in future lagged outcomes and may remain correlated with nearby units. The four implemented split-panel fits have a different purpose: they estimate and remove the leading loading- and factor-estimation biases for broad targets.

Recent rank-robust inference provides a useful contrast. \citet{choi2026rank} allow a prespecified rank no smaller than the truth and show that overestimating the low-rank space can itself create nonnegligible bias. Our tangent spaces and joint Riesz equation are rank specific, so the present theory instead requires consistent selection of every coefficient-matrix rank. Extending rank-robust inference would require controlling the effects of several overestimated and interacting singular spaces throughout the dynamic equation.

Adjacent-eigenvalue ratios are a standard device for estimating factor dimension; see \citet{ahnhorenstein2013}. The closest recent motivation for our selector is \citet{puetal2025}, who use a ridge-type eigenvalue ratio in a reduced-rank spatio-temporal model. \citet{barigozzietal2026} develop an iterative eigenvalue-ratio procedure for tensor factor dimensions under weak cross-sectional and temporal dependence. Their spectral objects, dependence conditions, and dimension-growth restrictions differ from ours. We use only the ridge-ratio principle and prove consistency from this paper's joint pilot rate, strong coefficient-matrix signals, and rectangular-panel growth condition. In particular, we do not import an aspect-ratio restriction.

The comparison also clarifies the scope of the low-rank assumptions. \citet{chernozhukov2023} and \citet{choi2024} permit approximate low rank or slowly growing rank in their settings. We impose exact fixed ranks because approximation errors in an autoregressive coefficient matrix propagate through future outcomes, while approximation errors from all matrices enter the same Gram and Riesz equations. Approximate-low-rank inference and growing-rank theory are left for future work.

Related work on matrix completion and causal panels includes \citet{athey2021matrix} and \citet{bai_ng2021matrix}. Those settings do not directly cover a dynamic equation with several stochastic-regressor coefficient matrices, predetermined covariates, and conditional spatial dependence. The conditional weak-dependence framework follows \citet{jenishprucha2012}, and spatial HAC variance estimation follows \citet{conley1999gmm}.

%================================================
\section{Model and Parameters of Interest}
\label{sec:model}

\subsection{Model specification}

We observe a complete panel for units $i=1,\ldots,N$ and periods $t=1,\ldots,T$. The model is
\begin{equation}
\label{eq:model}
y_{it}
=
\sum_{\ell=1}^{P}a_{0,it,\ell}y_{i,t-\ell}
+
\sum_{k=1}^{K}x_{it,k}\beta_{0,it,k}
+
h_{0,it}
+
u_{it}.
\end{equation}
Here $y_{i,t-\ell}$ is the $\ell$th lag of the dependent variable, $x_{it,k}$ is the $k$th observed covariate, and $u_{it}$ is the remaining unit-specific innovation. The term $h_{0,it}$ is an unobserved interactive fixed effect that captures latent common conditional-mean variation with heterogeneous unit loadings and is treated as a nuisance parameter. The autoregressive coefficients $a_{0,it,\ell}$, for $\ell=1,\ldots,P$, and the slope coefficients $\beta_{0,it,k}$, for $k=1,\ldots,K$, may vary across both units and time. The initial values $y_{i,0},\ldots,y_{i,1-P}$ are observed. Throughout, $P\ge1$ and $K$ are fixed and known.

Let $\mathscr C_t$ be the sigma-field generated by the common shocks at time $t$, and let
$
\mathscr C=\bigvee_{s=1}^{T}\mathscr C_s
$
collect the common shocks over the sample. These shocks need not be observed by the researcher, and conditioning on $\mathscr C$ is used only in the probability theory. The common shocks may affect the observed covariates, the heterogeneous coefficient paths, and the distribution of the unit-specific innovations. Their effects on the conditional mean of $y_{it}$ are represented by the complete fitted value $\mathcal Y_{it}(\bTheta_0)$: through the observed covariates, through the autoregressive and slope coefficient matrices when those coefficients vary with aggregate conditions, and through $\bH_0$ for additive latent common mean effects that are not mediated by the observed regressors. Thus, $\bH_0$ does not represent every effect of the common shocks. It represents the remaining additive common conditional-mean variation after the lagged outcomes, observed covariates, and heterogeneous coefficients have been included. The common shocks may additionally affect conditional variances and the strength of spatial dependence.

For each lag $\ell$, collect the autoregressive coefficients in the $N\times T$ matrix
$
\bA_0^{(\ell)}=(a_{0,it,\ell})_{i\le N,t\le T}.
$
For each covariate $k$, collect the slope coefficients in
$
\bB_0^{(k)}=(\beta_{0,it,k})_{i\le N,t\le T},
$
and write $\bH_0=(h_{0,it})_{i\le N,t\le T}$ for the interactive fixed-effect matrix. The maintained dimension-reduction restriction is that every one of these matrices has fixed low rank as $N,T\to\infty$:
$$
r_{A,\ell}=\rank\{\bA_0^{(\ell)}\},
\qquad
r_{B,k}=\rank\{\bB_0^{(k)}\},
\qquad
r_H=\rank(\bH_0).
$$
The rank vector is $\br_0=(r_{A,1},\ldots,r_{A,P},r_{B,1},\ldots,r_{B,K},r_H)$, and any rank may equal zero. No common rank, factor, or loading structure is imposed across coefficient matrices.

The low-rank restriction allows flexible heterogeneity while reducing the effective number of parameters. For example, if $\bA_0^{(1)}$ has rank one, then $a_{0,it,1}=\lambda_i f_t$ for unit loading $\lambda_i$ and time factor $f_t$. A larger rank permits several such terms. A rank-$r$ $N\times T$ matrix has $r(N+T-r)$ locally free parameters rather than $NT$. Thus, each coefficient matrix may vary over both dimensions while retaining an effective dimension of order $N+T$.

The model contains no unrestricted unit--time intercept outside $\bH_0$, because such a matrix would not generally be identified separately from $\bH_0$.

\subsection{Matrix notation}

Write
$$
\bTheta
=
\bigl(
\bA^{(1)},\ldots,\bA^{(P)},
\bB^{(1)},\ldots,\bB^{(K)},
\bH
\bigr)
\in
\left(\R^{N\times T}\right)^{P+K+1}.
$$
For two such collections, define
$$
\ip{\bQ}{\bDelta}
=
\sum_{\ell=1}^{P}\operatorname{tr}\{(\bQ_A^{(\ell)})'\bDelta_A^{(\ell)}\}
+
\sum_{k=1}^{K}\operatorname{tr}\{(\bQ_B^{(k)})'\bDelta_B^{(k)}\}
+
\operatorname{tr}(\bQ_H'\bDelta_H),
$$
and $\|\bQ\|^2=\ip{\bQ}{\bQ}$.

For each observation, define the scalar fitted-value map
\begin{equation}
\label{eq:fitted_value}
\mathcal Y_{it}:
\left(\R^{N\times T}\right)^{P+K+1}
\longrightarrow
\R,
\qquad
\mathcal Y_{it}(\bTheta)
=
\sum_{\ell=1}^{P}A_{it}^{(\ell)}y_{i,t-\ell}
+
\sum_{k=1}^{K}B_{it}^{(k)}x_{it,k}
+
H_{it}.
\end{equation}
Thus, $\mathcal Y_{it}(\bTheta)$ is one scalar, while
$
\mathcal Y(\bTheta)
=
\{\mathcal Y_{it}(\bTheta)\}_{i\le N,t\le T}
\in\R^{N\times T}
$
collects the fitted values for the full panel. Conditional on the observed regressors and lagged outcomes, $\mathcal Y_{it}(\bTheta)$ is linear in $\bTheta$, and the model can be written as $y_{it}=\mathcal Y_{it}(\bTheta_0)+u_{it}$.

\subsection{Parameters of interest}

We write the main formulas for one generic scalar parameter of interest. This avoids carrying an index through every equation; joint inference for any fixed finite collection is stated separately in the theorems. The parameter is a deterministic linear summary of the autoregressive and slope coefficient matrices:
\begin{equation}
\label{eq:target}
\phi(\bTheta_0)=\ip{\bD}{\bTheta_0}.
\end{equation}
The direction assigned to $\bH_0$ is zero. Every matrix in $\bD$ is $N\times T$, and matrices not involved in the parameter are set equal to zero.

Let $\be_i\in\R^N$ and $\be_t\in\R^T$ denote coordinate vectors. For a unit group $G_N\subseteq\{1,\ldots,N\}$ and a time set $R_T\subseteq\{1,\ldots,T\}$, define
$$
\bm a_{G_N}
=
\frac{1}{|G_N|}
\sum_{i\in G_N}\be_i,
\qquad
\bm b_{R_T}
=
\frac{1}{|R_T|}
\sum_{t\in R_T}\be_t.
$$
The directions for the main parameters are as follows.

\begin{enumerate}[label=(\roman*),leftmargin=2em]
\item For the individual autoregressive coefficient $a_{0,i_0t_0,\ell}$, set
$
\bD_A^{(\ell)}=\be_{i_0}\be_{t_0}',
$
and set every other target matrix to zero. For $\beta_{0,i_0t_0,k}$, use the same direction in $\bD_B^{(k)}$.

\item For the mean of a coefficient matrix $\bM_0$ over $G_N\times R_T$, set
$
\bD_M=\bm a_{G_N}\bm b_{R_T}'.
$
Then $\ip{\bD}{\bTheta_0}=|G_N|^{-1}|R_T|^{-1}\sum_{i\in G_N}\sum_{t\in R_T}M_{0,it}$.

\item For a cross-sectional mean at a fixed time $t_0$, set
$
\bD_M=\bm a_{G_N}\be_{t_0}'.
$
The full cross-sectional mean is obtained by taking $G_N=\{1,\ldots,N\}$.

\item For two unit groups $G_{1,N}$ and $G_{2,N}$, the time-$t_0$ contrast uses
$
\bD_M
=
(\bm a_{G_{1,N}}-\bm a_{G_{2,N}})\be_{t_0}'.
$
A contrast averaged over $R_T$ replaces $\be_{t_0}$ by $\bm b_{R_T}$.

\item For the unit--time lag-sum persistence index $\sum_{\ell=1}^{P}a_{0,i_0t_0,\ell}$, set
$
\bD_A^{(\ell)}=\be_{i_0}\be_{t_0}'
$
 for every $\ell=1,\ldots,P.
$
Group, time, and group--time averages of the lag sum use the same averaging direction in every autoregressive coefficient matrix.
\end{enumerate}

Empirical work often focuses on economically meaningful groups. For example, the mean of the $k$th slope coefficient for group $G_N$ at time $t$ is
$
\phi(\bTheta_0)
=
\frac{1}{|G_N|}
\sum_{i\in G_N}\beta_{0,it,k},
$
and the contrast between groups $G_{1,N}$ and $G_{2,N}$ is
$$
\phi(\bTheta_0)
=
\frac{1}{|G_{1,N}|}\sum_{i\in G_{1,N}}\beta_{0,it,k}
-
\frac{1}{|G_{2,N}|}\sum_{i\in G_{2,N}}\beta_{0,it,k}.
$$
Autoregressive coefficient means and contrasts are defined in the same way.

A target whose nonzero weights cover many units and many periods is called a \emph{two-way average}. This expression refers to the support of the target weights, not to the rank of any coefficient matrix. The full-panel mean is the leading example. The lag-sum target is a linear summary of the autoregressive coefficients; it is not an impulse response or a long-run multiplier.


%================================================
\section{Full-Panel Estimation, Rank Selection, and Variance Estimation}
\label{sec:estimation}

The procedure has four steps. First, estimate all coefficient matrices jointly under supplied or selected ranks. Second, use the estimated singular spaces to construct a Riesz representer for each parameter of interest. Third, use the full-panel plug-in estimate for local parameters and a fixed split-panel correction for averages over many units and time periods. Fourth, estimate the variance from the residuals and Riesz weights, allowing for spatial dependence. Theorems~\ref{thm:recovery}--\ref{thm:rank_consistency} provide the formal justification for these steps.

\subsection{Full-panel fixed-rank estimation}

A rank vector records the rank assigned to every coefficient matrix:
$$
\br
=
(r_{A,1},\ldots,r_{A,P},r_{B,1},\ldots,r_{B,K},r_H).
$$
For a supplied rank vector $\br$, let $\calM(\br)$ contain all collections $\bTheta$ for which $\rank\{\bA^{(\ell)}\}=r_{A,\ell}$, $\rank\{\bB^{(k)}\}=r_{B,k}$, and $\rank(\bH)=r_H$. Rank zero means that the corresponding matrix is identically zero.

Let $B<\infty$ be a fixed coefficient bound, and assume that the true matrices lie strictly inside this bound. Let $\|\bTheta\|_{\max}$ denote the largest absolute entry among all coefficient matrices. For supplied ranks, define
\begin{equation}
\label{eq:fixed_rank_estimator}
\widehat\bTheta(\br)
\in
\argmin_{\substack{\bTheta\in\calM(\br)\\ \|\bTheta\|_{\max}\le B}}
\frac1{2NT}\sum_{i=1}^{N}\sum_{t=1}^{T}
\{y_{it}-\mathcal Y_{it}(\bTheta)\}^2.
\end{equation}
Fix a supplied rank vector $\br$ and write $\widehat\bTheta=\widehat\bTheta(\br)$. Set
$$
\widehat u_{it}=y_{it}-\mathcal Y_{it}(\widehat\bTheta),
\qquad
\widehat\phi=\ip{\bD}{\widehat\bTheta}.
$$
Thus, the point estimate is obtained by inserting the estimated coefficient matrices into the parameter of interest. \Cref{thm:recovery} proves consistency for supplied ranks. The implementation reports numerical stationarity and agreement across starting values.

\subsection{Tangent spaces and Riesz representers}
\label{subsec:riesz}

If a positive-rank matrix has compact singular-value decomposition $\bM=\bU\bSigma\bV'$, its fixed-rank tangent space is
$$
\calT_M
=
\{\bU\bR'+\bm L\bV':\ \bR\in\R^{T\times r_M},\ \bm L\in\R^{N\times r_M}\}.
$$
A rank-zero matrix has tangent space $\{0\}$. Let $\calT_0$ and $\widehat\calT$ collect the tangent spaces at $\bTheta_0$ and $\widehat\bTheta$, and let $\calP_0$ and $\widehat\calP$ denote their orthogonal projections.

For a parameter $\phi(\bTheta)=\ip{\bD}{\bTheta}$, the map $\bDelta\mapsto\ip{\bD}{\bDelta}$ gives its first-order change along a tangent direction. Conditional on the common shocks $\mathscr C$, the least-squares Gram form on $\calT_0$ is
$$
(\bQ,\bDelta)
\longmapsto
\Ezero\!\left[
\sum_{i,t}\mathcal Y_{it}(\bQ)\mathcal Y_{it}(\bDelta)
\middle|\mathscr C
\right].
$$
When this Gram form is nonsingular, the finite-dimensional Riesz representation theorem gives a unique $\bq_0\in\calT_0$ satisfying
\begin{equation}
\label{eq:population_target_weight}
\Ezero\!\left[
\sum_{i,t}\mathcal Y_{it}(\bDelta)\mathcal Y_{it}(\bq_0)
\middle|\mathscr C
\right]
=
\ip{\bD}{\bDelta}
\quad
\text{for every }\bDelta\in\calT_0.
\end{equation}
The zero subscript indicates the population Riesz representer associated with the true coefficient matrices. It is the matrix analogue of a parameter derivative multiplied by the inverse of a population least-squares Gram matrix.

The empirical Riesz representer $\widehat\bq\in\widehat\calT$ solves
\begin{equation}
\label{eq:empirical_target_weight}
\sum_{i,t}\mathcal Y_{it}(\bDelta)\mathcal Y_{it}(\widehat\bq)
=
\ip{\bD}{\bDelta}
\quad
\text{for every }\bDelta\in\widehat\calT.
\end{equation}
Thus, $\bq_0$ and $\widehat\bq$ are the population and sample versions of the same object. The zero subscript denotes the population weight, while a hat denotes its sample estimate:
$$
\Psi_{0,it}=\mathcal Y_{it}(\bq_0),
\qquad
\widehat\Psi_{it}=\mathcal Y_{it}(\widehat\bq).
$$
Define the conditional variance of the population first-order term by
\begin{equation}
\label{eq:population_score_variance}
s_{NT}^2
=
\Ezero\!\left[
\left(
\sum_{i,t}\Psi_{0,it}u_{it}
\right)^2
\middle|\mathscr C
\right].
\end{equation}
\Cref{thm:recovery} shows that the empirical Riesz representer consistently estimates its population counterpart, and \Cref{thm:target_expansion} shows how the population Riesz weights enter the first-order distribution of $\widehat\phi$.

A tangent direction is a feasible infinitesimal change in the fitted coefficient matrices. Formally, $\bDelta\in\widehat\calT$ means that there is a differentiable curve $a\mapsto\bTheta(a)$ within the fitted fixed-rank set such that $\bTheta(0)=\widehat\bTheta$ and $d\bTheta(a)/da|_{a=0}=\bDelta$. If the coefficient bound is inactive and $\widehat\bTheta$ is a first-order stationary point, the least-squares objective cannot decrease to first order along any such curve. Since $\mathcal Y_{it}$ is linear in the coefficient matrices,
$$
0
=
\left.\frac{d}{da}
\frac1{2NT}\sum_{i,t}
\{y_{it}-\mathcal Y_{it}(\bTheta(a))\}^2
\right|_{a=0}
=
-\frac1{NT}\sum_{i,t}\widehat u_{it}\mathcal Y_{it}(\bDelta).
$$
Taking $\bDelta=\widehat\bq$ gives
\begin{equation}
\label{eq:weighted_residual_zero}
\sum_{i,t}\widehat\Psi_{it}\widehat u_{it}=0.
\end{equation}
This is the constrained least-squares normal equation, not an additional debiasing assumption. It explains why the Riesz-weighted residual is not added as a separate same-panel correction. The formal argument and the proof that the coefficient bound is inactive with probability approaching one are given in \Cref{app:recovery}.

\subsection{Point estimates and the split-panel correction}
\label{subsec:split_correction}

Use the full-panel plug-in $\widehat\phi$ for a regular individual coefficient and for a fixed unit--time lag sum. For a mean over $G_N\times R_T$, use the plug-in when
$$
\sqrt{|G_N||R_T|}
\left(\frac1N+\frac1T\right)
\longrightarrow0.
$$
For a contrast between disjoint groups $G_{1,N}$ and $G_{2,N}$ averaged over $R_T$, use the plug-in when
$$
\sqrt{
\frac{|R_T|}{|G_{1,N}|^{-1}+|G_{2,N}|^{-1}}
}
\left(\frac1N+\frac1T\right)
\longrightarrow0.
$$
\Cref{prop:target_score,thm:target_expansion} show that these conditions make the $T^{-1}$ loading bias and the $N^{-1}$ factor bias smaller than the corresponding standard error. The conditions hold automatically for a regular individual coefficient and a fixed unit--time lag sum.

For an average whose weights extend over many units and many time periods, the standard error can be as small as order $(NT)^{-1/2}$, so biases of order $T^{-1}$ and $N^{-1}$ may affect first-order inference. Draw a balanced random partition $\mathcal I_{T,1},\mathcal I_{T,2}$ of $\{1,\ldots,T\}$ and an independent balanced random partition $\mathcal I_{N,1},\mathcal I_{N,2}$ of $\{1,\ldots,N\}$. The unit partition is stratified by every unit group entering the parameter of interest. Let $T_h=|\mathcal I_{T,h}|$ and $N_g=|\mathcal I_{N,g}|$.

For a general two-way target, define $\bD^{T,h}$ by retaining in $\bD$ only the columns whose time indices belong to $\mathcal I_{T,h}$, and define $\bD^{N,g}$ by retaining only the rows whose unit indices belong to $\mathcal I_{N,g}$. All other entries are set to zero. Then
$$
\bD^{T,1}+\bD^{T,2}=\bD,
\qquad
\bD^{N,1}+\bD^{N,2}=\bD.
$$
Thus, this construction preserves any mean or contrast defined in \Cref{sec:model}, including a target supported on proper subsets $G_N$ and $R_T$. For a full mean, it is equivalent to averaging the two normalized split means with weights $T_h/T$ or $N_g/N$.

For $h=1,2$, $\widehat\bTheta_{T,h}$ is the fixed-rank estimator using all units and only times in $\mathcal I_{T,h}$; every selected time retains its observed lagged outcomes even when a lag belongs to the other time partition. For $g=1,2$, $\widehat\bTheta_{N,g}$ uses all times and only units in $\mathcal I_{N,g}$. Define the split target contributions by
$$
\widehat\phi_{T,h}
=
\ip{\bD^{T,h}}{\widehat\bTheta_{T,h}},
\qquad
\widehat\phi_{N,g}
=
\ip{\bD^{N,g}}{\widehat\bTheta_{N,g}},
$$
and set
$$
\overline\phi_T
=
\widehat\phi_{T,1}+\widehat\phi_{T,2},
\qquad
\overline\phi_N
=
\widehat\phi_{N,1}+\widehat\phi_{N,2}.
$$
The full-panel selected ranks are held fixed in all four fits. These splits are used only to estimate the leading incidental-parameter biases; they are not validation samples, and no observations are buffered.

Define the corrected estimator
\begin{equation}
\label{eq:half_panel_correction}
\widehat\phi^{\mathrm{corr}}
=
3\widehat\phi-\overline\phi_T-\overline\phi_N.
\end{equation}
The superscript $\mathrm{corr}$ denotes the split-panel bias-corrected estimate.

\Cref{thm:target_expansion} establishes that the full-panel bias has the form $B_T/T+B_N/N$. \Cref{prop:split_preservation} proves three facts used here: the restricted split targets reassemble the original parameter exactly; the random halves preserve the relevant ranks and Gram matrices; and the split bias constants equal the corresponding portions of the full-panel constants up to $o_p(1)$. Consequently, the aggregate time-split estimate has leading bias $2B_T/T+B_N/N$, while the aggregate unit-split estimate has leading bias $B_T/T+2B_N/N$. Hence
$$
3\left(\frac{B_T}{T}+\frac{B_N}{N}\right)
-
\left(\frac{2B_T}{T}+\frac{B_N}{N}\right)
-
\left(\frac{B_T}{T}+\frac{2B_N}{N}\right)
=0.
$$
Thus, \eqref{eq:half_panel_correction} removes both leading bias terms. \Cref{thm:twoway} proves the resulting asymptotic normality and the validity of its standard error.

\subsection{Variance estimation}
\label{subsec:variance}

For each observation $(i,t)$, $\Psi_{0,it}^{T}$ and $\Psi_{0,it}^{N}$ denote the population Riesz weights from the time split containing $t$ and the unit split containing $i$. Their sample versions are $\widehat\Psi_{it}^{T}$ and $\widehat\Psi_{it}^{N}$, with residuals $\widehat u_{it}^{T}$ and $\widehat u_{it}^{N}$. The conditional variance of the corrected first-order score is
\begin{equation}
\label{eq:corrected_population_variance}
(s_{NT}^{\mathrm{corr}})^2
=
\operatorname{Var}_0\!\left(
\sum_{i,t}u_{it}
\left[
3\Psi_{0,it}
-\Psi_{0,it}^{T}
-\Psi_{0,it}^{N}
\right]
\middle|\mathscr C
\right).
\end{equation}
This single variance includes every covariance among the full-panel and four split-panel estimates, which are based on overlapping observations. \Cref{prop:split_preservation} also shows that $(s_{NT}^{\mathrm{corr}})^2/s_{NT}^2\to_p1$, so the correction does not collapse the first-order score.

For an uncorrected parameter, conditional cross-sectional independence gives
$$
\widehat s_{\mathrm{diag}}^2
=
\sum_{i,t}
(\widehat\Psi_{it}\widehat u_{it})^2.
$$
Although \eqref{eq:weighted_residual_zero} makes the sum of the residual-weighted Riesz terms zero at an exact interior solution, the individual terms are not zero. Their squares and cross-products estimate the variance, just as observation-level residual--regressor products do in the usual heteroskedasticity-robust regression variance formula.

When unit-specific shocks may be spatially dependent, let $\bm W_N=(w_{ik,N})$ be a prespecified positive-semidefinite spatial weight matrix. It is symmetric, $w_{ii,N}=1$, $|w_{ik,N}|\le1$, and $w_{ik,N}=0$ whenever $d_N(i,k)>h_N$, where $h_N=\lceil c_h\log(NT)\rceil$ for a fixed sufficiently large $c_h$. The spatial HAC estimator is
\begin{equation}
\label{eq:spatial_variance}
\widehat s_{\mathrm{sp}}^2
=
\sum_{t=1}^{T}\sum_{i=1}^{N}\sum_{k=1}^{N}
w_{ik,N}
\widehat\Psi_{it}\widehat u_{it}
\widehat\Psi_{kt}\widehat u_{kt}.
\end{equation}
The matrix $\bm W_N$ retains nearby cross-unit covariances and excludes sufficiently distant pairs. If units belong to known conditionally independent clusters, set $w_{ik,N}=1$ within a cluster and zero across clusters. The diagonal estimator is the special case $\bm W_N=\bm I_N$. A general distance-based calculation uses a positive-semidefinite construction so that the reported variance is nonnegative.

For the corrected estimator, use the same spatial quadratic form with the full- and split-panel residual terms combined inside each observation:
\begin{align}
\widehat s_{\mathrm{corr}}^2
={}&
\sum_{t=1}^{T}\sum_{i=1}^{N}\sum_{k=1}^{N}
w_{ik,N}
\Big[
3\widehat\Psi_{it}\widehat u_{it}
-\widehat\Psi_{it}^{T}\widehat u_{it}^{T}
-\widehat\Psi_{it}^{N}\widehat u_{it}^{N}
\Big]
\nonumber\\
&\qquad\times
\Big[
3\widehat\Psi_{kt}\widehat u_{kt}
-\widehat\Psi_{kt}^{T}\widehat u_{kt}^{T}
-\widehat\Psi_{kt}^{N}\widehat u_{kt}^{N}
\Big].
\label{eq:corrected_spatial_variance}
\end{align}
This construction automatically includes the variances of the full- and split-panel estimates and every covariance among them. \Cref{thm:twoway} proves consistency of \eqref{eq:spatial_variance} and \eqref{eq:corrected_spatial_variance} for the corresponding population variances.

No temporal HAC terms are added because conditional sequential exogeneity makes the population time-level Riesz sums martingale differences. For $s<t$,
$$
\Ezero\!\left[
\left\{\sum_i\Psi_{0,it}u_{it}\right\}
\left\{\sum_i\Psi_{0,is}u_{is}\right\}
\middle|\mathscr C
\right]
=0.
$$
Thus, cross-time covariance does not enter the first-order variance, while same-time spatial covariance remains in \eqref{eq:spatial_variance} and \eqref{eq:corrected_spatial_variance}.

\subsection{Rank selection}
\label{subsec:rank_selection}

Let $\br=(r_{A,1},\ldots,r_{A,P},r_{B,1},\ldots,r_{B,K},r_H)$ collect the ranks assigned to the coefficient matrices. The reportable rank caps define
\begin{equation}
\label{eq:rank_reportable_set}
\calR_{\max}
=
\prod_{\ell=1}^{P}\{0,\ldots,\bar r_{A,\ell}\}
\times
\prod_{k=1}^{K}\{0,\ldots,\bar r_{B,k}\}
\times
\{0,\ldots,\bar r_H\}.
\end{equation}
Ranks are selected separately; no common rank or common singular space is imposed.

Write the normalized least-squares loss as $\mathcal L_{NT}(\bTheta)=(2NT)^{-1}\sum_{i,t}\{y_{it}-\mathcal Y_{it}(\bTheta)\}^2$. To evaluate the ridge ratio at a reportable cap, the pilot must reveal one additional singular value. We therefore compute one joint spectral pilot
\begin{equation}
\label{eq:rank_spectral_pilot}
\widehat\bTheta^{\mathrm{pil}}
\in
\argmin_{\substack{\|\bTheta\|_{\max}\le B\\
\rank(\bM)\le \bar r_M+1\ \text{for every coefficient matrix }\bM}}
\mathcal L_{NT}(\bTheta).
\end{equation}
The extra pilot rank is never reportable. The pilot is used only to estimate singular spectra and is not a reported coefficient estimator.
Write $\calD_{\max}^{+}$ for the admissible differences $\bTheta-\bTheta_0$ generated by this cap+1 pilot class and the maintained coefficient box.

Let
\begin{equation}
\label{eq:rank_scale_weights}
w_{A,\ell}=\left(\frac1{NT}\sum_{i,t}y_{i,t-\ell}^2\right)^{1/2},
\qquad
w_{B,k}=\left(\frac1{NT}\sum_{i,t}x_{it,k}^2\right)^{1/2},
\qquad
w_H=1.
\end{equation}
The reference scale $w_{A,1}$ is bounded above and away from zero in probability under the maintained conditions. For each coefficient matrix $M$ and $j\ge1$, define the normalized squared singular values
\begin{equation}
\label{eq:rank_scaled_spectrum}
\widehat\lambda_{M,j}
=
\left(\frac{w_M}{w_{A,1}}\right)^2
\frac{\sigma_j(\widehat\bM^{\mathrm{pil}})^2}{NT}.
\end{equation}
This normalization expresses each coefficient spectrum in relative fitted-value units. Under a fixed nonzero change of measurement units, with the coefficient parameterization and its box transformed consistently, the regressor scale and the corresponding coefficient matrix move inversely; the normalized spectrum is therefore unchanged. All empirical and Monte Carlo calculations use prespecified canonical units.

Set $a_{NT}=1/\log(NT)$ and $\widehat\lambda_{M,0}=1$. Define
\begin{equation}
\label{eq:rank_ratios}
R_{M,0}=\frac{\widehat\lambda_{M,1}+a_{NT}}{1+a_{NT}},
\qquad
R_{M,j}=\frac{\widehat\lambda_{M,j+1}+a_{NT}}
{\widehat\lambda_{M,j}+a_{NT}},
\quad j=1,\ldots,\bar r_M.
\end{equation}
The selected rank of matrix $M$ is
\begin{equation}
\label{eq:rank_selector}
\widehat r_M=\min\argmin_{0\le j\le\bar r_M}R_{M,j},
\qquad
\widehat\br=(\widehat r_{A,1},\ldots,\widehat r_{A,P},
\widehat r_{B,1},\ldots,\widehat r_{B,K},\widehat r_H).
\end{equation}
Thus rank zero is part of the same minimization rule and exact ties are resolved in favor of the smaller rank. At $j=\bar r_M$, the numerator uses the genuine $(\bar r_M+1)$st pilot singular value rather than a zero created by truncation.

Finally, compute one unpenalized constrained fixed-rank estimator $\widehat\bTheta(\widehat\br)$ from \eqref{eq:fixed_rank_estimator}. Every reported coefficient, target estimate, tangent space, Riesz representer, split correction, residual, standard error, and confidence interval is constructed from this final estimator. Rank selection is performed once on the full panel and is not repeated in the split panels.

\section{Assumptions}
\label{sec:assumptions}

Let $\mathscr F_{t-1}=\sigma(y_{is},x_{is}:i\le N,s\le t-1)$, including the observed presample outcomes. The common-shock sigma-field $\mathscr C$ was defined in \Cref{sec:model}. Conditional on $\mathscr C$, the true coefficient matrices and deterministic parameter weights are treated as fixed. The true rank vector $\br_0$ is nonrandom and does not depend on the realization of $\mathscr C$. A \emph{regular common-shock realization} is a realization for which the conditional moment, dependence, covariance, signal, and Gram bounds stated below hold with their deterministic envelopes; the assumptions require this set of realizations to have probability one. A \emph{regular individual coefficient} is an entry target satisfying the leverage and target-variance noncancellation requirements in \Cref{a:target}. All conditional statements below hold for almost every regular common-shock realization. Fix $\eta>0$. Throughout this section, $c$ and $C$ denote fixed positive bounded constants that do not depend on $N$, $T$, $i$, $t$, or the realization of the common shocks; their values may differ from line to line. Other constants introduced below are also fixed unless stated otherwise.

\begin{assumption}[Dynamic stability]
\label{a:stab}
Let $\bC_{0,it}$ be the companion matrix formed from $a_{0,it,1},\ldots,a_{0,it,P}$. For some $\rho\in(0,1)$,
\begin{equation}
\label{eq:stability}
\sup_{i,h,t+h\le T}
\|\bC_{0,i,t+h}\cdots\bC_{0,i,t+1}\|_{\op}
\le C\rho^h.
\end{equation}
The observed initial outcome states satisfy $\sup_{i,s=1-P,\ldots,0}\Ezero(|y_{is}|^{8+\eta}\mid\mathscr C)\le C$ almost surely.
\end{assumption}
This assumption permits time-varying autoregressive coefficients but rules out explosive propagation. The effect of a current disturbance through future lagged outcomes decreases geometrically with the time horizon, while the initial observations have the same conditional moment order used elsewhere in the paper.

\begin{assumption}[Conditional sequential exogeneity]
\label{a:exog}
For every $i$ and $t$,
\begin{equation}
\label{eq:conditional_exogeneity}
\Ezero[u_{it}\mid x_{it},\mathscr F_{t-1},\mathscr C]=0.
\end{equation}
\end{assumption}
This is the conditional sequential-exogeneity formulation used in dynamic interactive-effects and dynamic spatial-panel models; see \citet{moonweidner2017} and \citet{kuersteinerprucha2020}. It defines $u_{it}$ as the new unit-specific innovation after the current covariates, past panel information, and common shocks are held fixed. Future unit-level covariates are not in the conditioning set and may respond to past outcome disturbances.

\begin{assumption}[Spatial sampling geometry]
\label{a:geometry}
The unit locations lie in $\R^{d_s}$ for fixed $d_s$ and are uniformly separated. If
$$
v_N(r)
=
\max_{i\le N}
\left|\{j\le N:d_N(i,j)\le r\}\right|,
$$
then $v_N(r)\le C(1+r)^{d_s}$ uniformly in $N$ and $r\ge0$.
\end{assumption}
This assumption limits how quickly the number of spatial neighbors can grow with distance. It allows irregular locations while making sums of local spatial covariances manageable.

Let $\varepsilon_{it}^x$ collect the unit-specific innovations that generate the covariates and set $\varepsilon_{it}=((\varepsilon_{it}^x)',u_{it})'$. For a set $\mathcal A$ of unit--time indices, write
$
\mathscr G(\mathcal A)
=
\sigma\{\varepsilon_{it}:(i,t)\in\mathcal A\}.
$
For two index sets $\mathcal A$ and $\mathcal B$, define
$$
\operatorname{dist}(\mathcal A,\mathcal B)
=
\min_{\substack{(i,t)\in\mathcal A\\(j,s)\in\mathcal B}}
\{d_N(i,j)+|t-s|\}.
$$
Conditional on $\mathscr C$, the strong-mixing, or $\alpha$-mixing, coefficient at separation $a$ is
$$
\begin{aligned}
\alpha_{\mathscr C}(a)
=
\sup_{\operatorname{dist}(\mathcal A,\mathcal B)\ge a}
\sup_{\substack{E\in\mathscr G(\mathcal A)\\F\in\mathscr G(\mathcal B)}}
\big|
\Pzero(E\cap F\mid\mathscr C)
-\Pzero(E\mid\mathscr C)\Pzero(F\mid\mathscr C)
\big|.
\end{aligned}
$$
Let $\bar\alpha(a)$ be a deterministic, nonincreasing sequence, independent of $N,T$ and the realization of $\mathscr C$, that uniformly bounds $\alpha_{\mathscr C}(a)$. For $h,r\ge0$, let
$$
\mathscr E_{it}(h,r)
=
\sigma\Big(
\varepsilon^x_{js}:
t-h\le s\le t,\ d_N(i,j)\le r;
u_{js}:
t-h\le s\le t-1,\ d_N(i,j)\le r
\Big)
$$

\begin{assumption}[Conditional alpha mixing and near-epoch dependence]
\label{a:ned}
Almost surely, $\alpha_{\mathscr C}(a)\le\bar\alpha(a)\le Ce^{-ca}$. Moreover, for every covariate,
\begin{equation}
\label{eq:conditional_ned}
\sup_{i,t,k}
\left\{
\Ezero\!\left[
\left|
x_{it,k}
-
\Ezero[x_{it,k}\mid\mathscr E_{it}(h,r)\vee\mathscr C]
\right|^{8+\eta}
\middle|\mathscr C
\right]
\right\}^{1/(8+\eta)}
\le
C\{\psi_T(h)+\psi_S(r)\},
\end{equation}
where $\psi_T(h)\le Ce^{-ch}$ and $\psi_S(r)\le Ce^{-cr}$. Current covariate innovations are included in $\mathscr E_{it}(h,r)$, while the current outcome innovation is excluded.
\end{assumption}
Conditional on the common shocks, distant unit-specific innovations become nearly independent, while each covariate can be approximated by recent innovations from nearby units. The deterministic envelope $\bar\alpha(a)$ permits the realized conditional dependence to vary with the common shocks while imposing one uniform decay rate. The near-epoch-dependence formulation follows the spatial random-field framework of \citet{jenishprucha2012}, adapted here to the combined unit--time index and to conditioning on common shocks. Excluding current $u_{it}$ from the approximation preserves the predetermined timing of the covariates.

\begin{assumption}[Conditional moments and weak cross-sectional covariance]
\label{a:moments}
Let $\bm u_t=(u_{1t},\ldots,u_{Nt})'$, let $\bm x_t$ collect the current covariates of all units, and define
$
\bm\Sigma_{u,t}
=
\Ezero[\bm u_t\bm u_t'\mid\bm x_t,\mathscr F_{t-1},\mathscr C].
$

\begin{enumerate}
\item[(i)] For the fixed $\eta>0$, uniformly in $i,t,k$
$$
\Ezero[|u_{it}|^{8+\eta}\mid x_{it},\mathscr F_{t-1},\mathscr C]\le C,
\qquad
\Ezero[|x_{it,k}|^{8+\eta}\mid\mathscr C]\le C.
$$

\item[(ii)] Uniformly in $t$,
$$
c\le (\bm\Sigma_{u,t})_{ii}\le C
\quad\text{for every }i, \quad
\max_{i\le N}\sum_{j=1}^{N}|(\bm\Sigma_{u,t})_{ij}|\le C.
$$
\end{enumerate}
\end{assumption}
The moment bounds control the products and quadratic forms used in estimation and variance calculation. The row-sum bound permits conditional spatial dependence but prevents the variance of a cross-sectional sum from growing because each unit is strongly related to too many other units. It implies $\|\bm\Sigma_{u,t}\|_{\op}\le C$. We do not require the smallest eigenvalue of the entire $N\times N$ covariance matrix to be uniformly positive. Instead, Assumption~\ref{a:target} imposes nondegeneracy only in the Riesz directions used for inference. For comparison, \citet{baich2021} impose the stronger requirement that each relevant covariance block have a smallest eigenvalue bounded away from zero and bounded absolute row sums.

\begin{assumption}[Separate low-rank signals and incoherence]
\label{a:signal}
For every positive-rank matrix
$$
\bM_0
\in
\{\bA_0^{(1)},\ldots,\bA_0^{(P)},
\bB_0^{(1)},\ldots,\bB_0^{(K)},\bH_0\},
$$
let $\bM_0=\bU_M\bSigma_M\bV_M'$ be its compact singular-value decomposition. Uniformly over these matrices,
$$
c\sqrt{NT}
\le
\sigma_{r_M}(\bM_0)
\le
\sigma_1(\bM_0)
\le
C\sqrt{NT},
$$
and
$$
\max_i\|\be_i'\bU_M\|_2^2\le \frac{\mu r_M}{N},
\qquad
\max_t\|\be_t'\bV_M\|_2^2\le \frac{\mu r_M}{T}.
$$
Every true coefficient entry is bounded in absolute value by $B-c_B$ for some fixed interior margin $c_B\in(0,B)$. No singular-value condition is imposed on a rank-zero matrix.
\end{assumption}
The singular-value bounds make every nonzero low-rank matrix statistically visible, while incoherence spreads its variation across units and time periods. Thus, no single row or column carries the entire low-rank signal. The margin $c_B$ need not be small; it only places the true coefficients a fixed positive distance from the boundary of the estimation set. Together with max-norm consistency, this makes the coefficient bound asymptotically inactive and permits the ordinary fixed-rank first-order conditions to be used.

\begin{assumption}[Conditional prediction identification]
\label{a:identification}
The lower bound
\begin{equation}
\label{eq:conditional_identification}
\Ezero\!\left[
\sum_{i,t}\mathcal Y_{it}(\bDelta)^2
\middle|\mathscr C
\right]
\ge c\|\bDelta\|^2
\end{equation}
holds in each of the following cases:
\begin{enumerate}[label=(\roman*)]
\item \emph{Global identification.}
For every nonzero admissible difference
$\bDelta=\bTheta-\bTheta_0$, where $\bTheta$ has a rank
vector in $\calR_{\max}$ and satisfies
$\|\bTheta\|_{\max}\le B$.

\item \emph{Local identification.}
For every nonzero tangent direction
$\bDelta\in\calT_0$.
\end{enumerate}
\end{assumption}
Part (i) ensures that two distinct admissible low-rank coefficient
collections cannot generate the same conditional fitted values.
Part (ii) is the local identification condition used for the joint
least-squares Gram form and the Riesz equation. A simple sufficient setting is one in which the lagged outcomes, covariates, and constant have conditionally nonsingular second moments after their low-rank directions are taken into account. Identification can fail, for example, if a covariate is exactly equal to a lagged outcome and the associated coefficient matrices can offset one another, or if a constant covariate duplicates the unrestricted role of $\bH_0$. The corresponding upper bound follows from the fixed numbers of lags and covariates and the moment bounds in \Cref{a:moments}.

\begin{assumption}[Pilot-only prediction identification]
\label{a:pilot_identification}
For some fixed, possibly smaller constant $c_+>0$, the global identification condition in Assumption~\ref{a:identification}(i) additionally holds on the cap+1 pilot difference class $\calD_{\max}^{+}$:
\[
\Ezero\!\left[
\sum_{i,t}\mathcal Y_{it}(\bDelta)^2
\middle|\mathscr C
\right]
\ge
c_+\|\bDelta\|^2
\]
for every nonzero $\bDelta\in\calD_{\max}^{+}$.
\end{assumption}
This is a localized strengthening of the identification domain for data-driven rank selection. Supplied-rank estimation and inference continue to use the original reportable-rank domain in Assumption~\ref{a:identification}(i), while the local tangent-space condition in Assumption~\ref{a:identification}(ii) is not enlarged. The extension is needed only to estimate one singular value beyond each reportable cap. No temporal, spatial, moment, signal, target, or panel-growth assumption is strengthened.

For the next assumption, write every positive-rank matrix as
$$
M_{0,it}=\bm\lambda_{M,0,i}'\bm f_{M,0,t},
\qquad
\frac1N\sum_{i=1}^{N}\bm\lambda_{M,0,i}\bm\lambda_{M,0,i}'=\bm I_{r_M}.
$$
Let $r_+=\sum_{\ell=1}^{P}r_{A,\ell}+\sum_{k=1}^{K}r_{B,k}+r_H$, omitting zero-rank matrices. Set $z_{it,A^{(\ell)}}=y_{i,t-\ell}$, $z_{it,B^{(k)}}=x_{it,k}$, and $z_{it,H}=1$. Stack the $r_+$ loading-regression terms $z_{it,M}\bm f_{M,0,t}$ in $\bm g_{it}^{\lambda}\in\R^{r_+}$, and stack the factor-regression terms $z_{it,M}\bm\lambda_{M,0,i}$ in $\bm g_{it}^{f}\in\R^{r_+}$.

\begin{assumption}[Loading and factor Gram matrices]
\label{a:gram}
Let $\bm I$ below denote the $r_+\times r_+$ identity matrix. Uniformly over $i$ and $t$,
$$
c\bm I
\preceq
\Ezero\!\left[
\frac1T\sum_{s=1}^{T}
\bm g_{is}^{\lambda}(\bm g_{is}^{\lambda})'
\middle|\mathscr C
\right]
\preceq
C\bm I,
$$
and
$$
c\bm I
\preceq
\Ezero\!\left[
\frac1N\sum_{j=1}^{N}
\bm g_{jt}^{f}(\bm g_{jt}^{f})'
\middle|\mathscr C
\right]
\preceq
C\bm I.
$$
\end{assumption}
These are ordinary least-squares second-moment matrices, not autoregressive or slope coefficient matrices. The assumption requires enough time-series variation to estimate each unit's loading rows and enough cross-sectional variation to estimate each time factor row. A close recent analogue is Assumption~5 of \citet{wangsu2022}, which imposes uniform eigenvalue bounds on the row- and column-wise regression matrices in a low-rank panel model. \citet{pengetal2025} impose a related positive-definiteness condition on regressor moments after accounting for interactive effects.

\begin{assumption}[Regularity of the parameters of interest]
\label{a:target}
If a unit or time averaging weight has growing support, its largest absolute value converges to zero. For a mean or contrast,
$$
\|\calP_0\bD\|\ge c\|\bD\|.
$$
For an individual entry of a positive-rank matrix $\bM_0=\bU_M\bSigma_M\bV_M'$, the selected unit and time satisfy
$$
\|\be_i'\bU_M\|_2^2\ge c/N,
\qquad
\|\be_t'\bV_M\|_2^2\ge c/T.
$$
Let $\bm\psi_t=(\Psi_{0,1t},\ldots,\Psi_{0,Nt})'$. For every parameter of interest,
$$
\sum_{t=1}^{T}
\Ezero[
\bm\psi_t'\bm\Sigma_{u,t}\bm\psi_t
\mid\mathscr C]
\ge
c\,
\Ezero\!\left[
\sum_{t=1}^{T}\|\bm\psi_t\|_2^2
\middle|\mathscr C
\right].
$$
For every fixed finite vector of parameters of interest, the limiting conditional correlation matrix of their variance-normalized first-order Riesz terms is positive definite.
\end{assumption}
The first restriction prevents one unit or one time from dominating a growing average. The projection and leverage restrictions ensure that the chosen parameter changes along statistically identified low-rank directions. The displayed variance restriction rules out near cancellation only for the Riesz-weighted combinations used in the reported inference; unlike a global lower eigenvalue bound on $\bm\Sigma_{u,t}$, it does not restrict irrelevant cross-sectional directions. The final restriction rules out redundant linear combinations in joint inference.

Let
\begin{equation}
\label{eq:tail_scale}
b_{NT}=(NT)^{1/(8+\eta)}\log(NT).
\end{equation}
The quantity $b_{NT}$ is used only in the theoretical rates and reflects the finite $8+\eta$ moments; it is not selected from the data. The same first-stage rate also controls the spectral-pilot error used by the rank-selection theorem, while the deterministic ridge $a_{NT}=1/\log(NT)$ is specified in \Cref{subsec:rank_selection}.

\begin{assumption}[Panel growth]
\label{a:growth}
As $N,T\to\infty$,
\begin{equation}
\label{eq:panel_growth}
(NT)^{4/(8+\eta)}
\{\log(NT)\}^{d_s+6}
\left(\frac1N+\frac1T\right)
\longrightarrow0.
\end{equation}
\end{assumption}
This condition allows $N$ and $T$ to grow at different rates but requires both dimensions to be large enough relative to the moment and spatial-complexity terms. It delivers the recovery rate used by the full-panel estimator and the spectral pilot. It also implies $\zeta_{NT}\log(NT)\to0$, which separates the singular values beyond the true rank from the deterministic ridge $a_{NT}=1/\log(NT)$.

The assumptions above contain the maintained restrictions on the data-generating process, parameters of interest, and panel growth. The only rank-selection-specific strengthening is Assumption~\ref{a:pilot_identification}, which enlarges the prediction-identification domain for the spectral pilot. The appendix derives the uniform low-rank score and Gram bounds, the proof-only local-deletion rates, the preservation of ranks and Gram matrices by the balanced random splits, feasible spatial variance estimation, and rank consistency. Short numerical requirements for approximate optimization are stated with the results that use them.


%================================================
\section{Asymptotic Theory}
\label{sec:theory}

All probability statements in this section are conditional on $\mathscr C$. The deterministic bounds in the assumptions make the conclusions uniform over regular realizations of the common shocks, so conditional validity implies unconditional validity. Let $\Phi(z)$ denote the cumulative distribution function of a standard normal random variable, and define the first-stage rate
$$
\zeta_{NT}
=
b_{NT}^2\{\log(NT)\}^{d_s+2}
\left(\frac1N+\frac1T\right).
$$
The logarithmic power reflects the fixed-dimensional unit--time blocking and the low-rank covering argument in \Cref{lem:weighted_concentration,prop:uniform_concentration}. The panel-growth assumption is stronger than what is needed for $\zeta_{NT}\to0$ and continues to imply all rank-selection and split-panel rates below.
Proofs are collected in the mathematical appendix: \Cref{app:population} proves \Cref{thm:population_clt}; \Cref{app:recovery,app:weights} prove \Cref{thm:recovery}; \Cref{app:bias} proves \Cref{thm:target_expansion,thm:twoway}; \Cref{app:variance} proves the variance statements in \Cref{thm:twoway}; and \Cref{app:rank} proves \Cref{thm:rank_consistency}.

\begin{theorem}[Conditional normality of the population Riesz term]
\label{thm:population_clt}
Under \Cref{a:stab,a:exog,a:geometry,a:ned,a:moments,a:signal,a:identification,a:gram,a:target,a:growth},
$$
\frac{\sum_{i,t}\Psi_{0,it}u_{it}}{s_{NT}}
\xrightarrow{d}
N(0,1)
\qquad
\text{conditionally on }\mathscr C.
$$
Equivalently,
\begin{equation}
\label{eq:conditional_clt}
\sup_{z\in\R}
\left|
\Pzero\!\left[
\frac{\sum_{i,t}\Psi_{0,it}u_{it}}{s_{NT}}\le z
\middle|\mathscr C
\right]
-\Phi(z)
\right|
\longrightarrow_p0.
\end{equation}
The conclusion holds jointly for every fixed finite vector of parameters of interest after conditional covariance normalization.
\end{theorem}

\Cref{a:identification} is a population restriction: it says that distinct admissible coefficient matrices imply distinct conditional fitted values. The next theorem is a sample result. It uses identification together with the moment and dependence assumptions to show that the least-squares estimator, its singular spaces, and the empirical Riesz representer converge to their population counterparts.

\begin{theorem}[Consistency of the full-panel estimator and Riesz representer]
\label{thm:recovery}
Suppose the true ranks are supplied and \eqref{eq:fixed_rank_estimator} is solved exactly. Under \Cref{a:stab,a:exog,a:geometry,a:ned,a:moments,a:signal,a:identification,a:gram,a:growth},
$$
\frac1{NT}\|\widehat\bTheta(\br_0)-\bTheta_0\|^2
+
\frac1{NT}\sum_{i,t}
\mathcal Y_{it}\{\widehat\bTheta(\br_0)-\bTheta_0\}^2
=
O_p(\zeta_{NT}),
$$
$$
\|\widehat\bTheta(\br_0)-\bTheta_0\|_{\max}=o_p(1),
\qquad
\|\widehat\calP-\calP_0\|_{\op}=O_p(\zeta_{NT}^{1/2}).
$$
If \Cref{a:target} also holds for $\phi$, then
$$
\frac{\|\widehat\bq-\bq_0\|}{\|\calP_0\bD\|}=o_p(1),
\qquad
\sum_{i,t}(\widehat\Psi_{it}-\Psi_{0,it})^2
=o_p(\|\calP_0\bD\|^2).
$$
\end{theorem}

The first line establishes consistency of all separately ranked coefficient matrices in Frobenius and prediction norms. The second line gives entrywise and tangent-space recovery, and the last line justifies replacing the population Riesz weights by feasible weights. Root-$T$ loading-row and root-$N$ factor-row expansions used in the bias proof are given in \Cref{app:recovery}. The same conclusions hold for numerical solutions satisfying the objective, first-order, and Riesz-equation tolerances reported in \Cref{sec:implementation}.

\begin{theorem}[Full-panel expansion and asymptotic normality]
\label{thm:target_expansion}
Suppose the true ranks are supplied. Under \Cref{a:stab,a:exog,a:geometry,a:ned,a:moments,a:signal,a:identification,a:gram,a:target,a:growth}, there are bounded $\mathscr C$-measurable terms $B_T$ and $B_N$ such that
\begin{equation}
\label{eq:target_expansion}
\widehat\phi-\phi(\bTheta_0)
=
\sum_{i,t}\Psi_{0,it}u_{it}
+
\frac{B_T}{T}
+
\frac{B_N}{N}
+
r_{NT},
\end{equation}
where $r_{NT}=o_p(s_{NT})+o_p(N^{-1}+T^{-1})$. The $T^{-1}$ term is generated by loading estimation and predetermined dynamic regressors; the $N^{-1}$ term is generated by factor estimation and contemporaneous cross-sectional dependence.

For a regular individual coefficient or a fixed unit--time lag sum,
$$
\frac{\widehat\phi-\phi(\bTheta_0)}{s_{NT}}
\xrightarrow{d}
N(0,1)
\qquad
\text{conditionally on }\mathscr C.
$$
The same conclusion holds for a mean over $G_N\times R_T$ when
$$
\sqrt{|G_N||R_T|}
\left(\frac1N+\frac1T\right)
\longrightarrow0,
$$
and for a contrast between disjoint groups $G_{1,N}$ and $G_{2,N}$ averaged over $R_T$ when
$$
\sqrt{
\frac{|R_T|}{|G_{1,N}|^{-1}+|G_{2,N}|^{-1}}
}
\left(\frac1N+\frac1T\right)
\longrightarrow0.
$$
\end{theorem}

\Cref{prop:target_score} proves the variance rates used in this statement. In particular, the support-growth restriction is automatic for a regular individual coefficient and a fixed unit--time lag sum. A full-panel mean generally fails the displayed mean condition and is therefore covered by the split-panel correction in \Cref{thm:twoway}.

\begin{theorem}[Split-panel correction and feasible asymptotic normality]
\label{thm:twoway}
Suppose the true ranks are supplied. Use the independent balanced time and unit partitions defined in \Cref{subsec:split_correction}, drawn independently of the panel, and hold the full-panel ranks fixed in all four split-panel fits. The conclusions hold conditionally on $\mathscr C$ for almost every partition sequence. Under \Cref{a:stab,a:exog,a:geometry,a:ned,a:moments,a:signal,a:identification,a:gram,a:target,a:growth},
$$
\widehat\phi^{\mathrm{corr}}-\phi(\bTheta_0)
=
\sum_{i,t}u_{it}
\left[
3\Psi_{0,it}
-\Psi_{0,it}^{T}
-\Psi_{0,it}^{N}
\right]
+o_p(s_{NT}^{\mathrm{corr}}).
$$
Moreover,
$$
\frac{(s_{NT}^{\mathrm{corr}})^2}{s_{NT}^2}
\to_p1,
\qquad
\frac{\widehat s_{\mathrm{corr}}^2}{(s_{NT}^{\mathrm{corr}})^2}
\to_p1,
\qquad
\frac{\widehat\phi^{\mathrm{corr}}-\phi(\bTheta_0)}
{\widehat s_{\mathrm{corr}}}
\xrightarrow{d}
N(0,1)
\quad
\text{conditionally on }\mathscr C.
$$
The analogous feasible normality result holds for an uncorrected parameter with \eqref{eq:spatial_variance}.
\end{theorem}

\begin{theorem}[Ridge-ratio rank consistency and selected-rank inference]
\label{thm:rank_consistency}
Suppose \Cref{a:stab,a:exog,a:geometry,a:ned,a:moments,a:signal,a:identification,a:pilot_identification,a:gram,a:growth} hold and $\br_0\in\calR_{\max}$. If a spectral pilot satisfies
\begin{equation}
\label{eq:pilot_operator_rate}
\max_M\|\widehat\bM^{\mathrm{pil}}-\bM_0\|_{\op}
=O_p\{\sqrt{NT\zeta_{NT}}\},
\end{equation}
then the selector in \eqref{eq:rank_selector}, with $a_{NT}=1/\log(NT)$, obeys
\begin{equation}
\label{eq:rank_consistency}
\Pzero(\widehat\br=\br_0\mid\mathscr C)\longrightarrow_p1.
\end{equation}
Consequently, \Cref{thm:recovery,thm:target_expansion,thm:twoway} remain valid when the supplied ranks are replaced by the selected ranks. Rank selection is performed only on the full panel.
\end{theorem}

\begin{corollary}[Rate of the spectral pilot]
\label{cor:rank_pilot_rate}
Suppose \Cref{a:stab,a:exog,a:geometry,a:ned,a:moments,a:signal,a:identification,a:pilot_identification,a:growth} hold, $\br_0\in\calR_{\max}$, and the pilot in \eqref{eq:rank_spectral_pilot} has normalized objective gap
\begin{equation}
\label{eq:rank_pilot_objective_gap}
\mathcal L_{NT}(\widehat\bTheta^{\mathrm{pil}})
-
\inf_{\substack{\|\bTheta\|_{\max}\le B\\
\rank(\bM)\le\bar r_M+1\ \forall M}}
\mathcal L_{NT}(\bTheta)
=o_p(\zeta_{NT}).
\end{equation}
Then $\|\widehat\bTheta^{\mathrm{pil}}-\bTheta_0\|=O_p\{\sqrt{NT\zeta_{NT}}\}$ and the operator-rate condition \eqref{eq:pilot_operator_rate} holds. Feasibility, agreement across starts, objective stability, stationarity, and boundary activity are observable numerical diagnostics, but they do not certify the unknown global infimum in \eqref{eq:rank_pilot_objective_gap}.
\end{corollary}

\begin{remark}[Rectangular panels and rank zero]
\label{rem:ridge_rectangular}
The rank result requires neither $N=T$ nor a limiting aspect ratio. Under \Cref{a:growth}, $a_{NT}\to0$ and $\zeta_{NT}/a_{NT}=\zeta_{NT}\log(NT)\to0$ for every admitted rectangular sequence. The deterministic anchor in \eqref{eq:rank_ratios} includes rank-zero coefficient matrices without an additional test or cutoff.
\end{remark}

\section{Implementation}
\label{sec:implementation}

The implemented procedure follows the estimators in \Cref{sec:estimation}. It uses no validation folds and no forward or spatial buffers. The local deletions used in the proofs are never computed.

\subsection{Practical implementation}

A practical implementation proceeds as follows.
\begin{enumerate}[leftmargin=2em]
\item \textit{Preprocessing and caps.} Store the scientific units, compute the scale weights in \eqref{eq:rank_scale_weights}, specify each reportable rank cap $\bar r_M$, set the pilot rank bound to $\bar r_M+1$, and retain the coefficient box $B$.
\item \textit{Spectral pilot.} Solve \eqref{eq:rank_spectral_pilot} jointly for all coefficient matrices. Record all start objectives, objective stability, first-order residuals, feasibility, boundary activity, numerical ranks, and runtime.
\item \textit{Rank selection.} Extract the singular values through $\bar r_M+1$, form \eqref{eq:rank_scaled_spectrum} and \eqref{eq:rank_ratios}, and select each coefficient-matrix rank by \eqref{eq:rank_selector}.
\item \textit{Final coefficient estimate.} Compute one literal unpenalized constrained fixed-rank estimator $\widehat\bTheta(\widehat\br)$. The spectral pilot is not substituted for this estimator.
\item \textit{Inference.} From the final estimator, construct tangent projections, empirical Riesz representers, plug-in targets, the four fixed-rank split fits required for broad targets, the two-way correction, residuals, spatial variances, and intervals.
\end{enumerate}

If the spectral pilot is infeasible, nonfinite, fails the prespecified objective or stationarity criteria, or remains unstable across the maintained starts, record \texttt{rank\_selection\_numerically\_unresolved}. There is no fallback rank. In particular, the procedure does not substitute a pilot rank, an alternative spectral rule, or the supplied truth. No primary selected-rank point estimate or inference is reported for that replication. This rule is distinct from a boundary-active final post-refit, whose point-estimate and interiority treatment remains as stated below.

For a supplied rank vector, each positive-rank coefficient matrix may be written as $M=L_MF_M'$, while a rank-zero matrix is set equal to zero. Conditional on factor rows, the least-squares criterion separates into loading regressions across units; conditional on loading rows, it separates into factor regressions across time. Warm-started alternating least squares or fixed-rank manifold methods may therefore be used. The optimization remains joint because every update uses residuals formed from all coefficient matrices. Multiple starts, objective agreement, stationarity, rank preservation, and box feasibility are reported for both the spectral pilot and the final fit. These diagnostics do not replace the theoretical global objective-gap condition.

After estimating the coefficient matrices, the empirical Riesz equation is solved by matrix--vector products rather than by forming a dense tangent-space basis. The same linear operator and preconditioner can be reused across parameters; only the target direction changes. For the spatial variance calculation, neighbor lists are precomputed for pairs within the distance cutoff, giving cost proportional to $TNv_N(h_N)$ instead of $TN^2$.

\subsection{Tuning choices}
\label{subsec:tuning}

The rank selector has three prespecified ingredients: the reportable caps $\bar r_M$, pilot rank bounds $\bar r_M+1$, and the deterministic ridge $a_{NT}=1/\log(NT)$. It has no information-criterion multiplier, ridge multiplier, singular-value cutoff, rank-zero test, or tuning sensitivity chosen from the reported simulations. Numerical optimization tolerances are fixed in advance and reported with the pilot and final-fit diagnostics.

The coefficient bound $B$ is fixed ex ante from a deterministic support bound for the maintained coefficient matrices. The margin $c_B$ in \Cref{a:signal} places the truth a fixed distance inside the estimation set. A boundary-active constrained solution remains a valid point estimate of the finite-sample constrained problem, but primary theorem-based inference is suppressed when the asymptotic argument requires the bound to be inactive.

For general spatial dependence, the spatial weight matrix and distance cutoff $h_N=\lceil c_h\log(NT)\rceil$ are prespecified. Conditional cross-sectional independence uses $W_N=I_N$, while known clusters use within-cluster weights. The implementation stores the random balanced split assignments and all numerical diagnostics.

\section{Monte Carlo Experiments}
\label{sec:MC}

We consider four data-generating processes arranged in increasing order of difficulty. DGP~1 introduces unit-level heteroskedasticity. DGP~2 adds spatial dependence. DGP~3 adds predetermined covariates. DGP~4 retains the spatial and predetermined structure of DGP~3 and adds systematic group heterogeneity in the autoregressive and slope coefficient matrices.

\subsection{Common dynamic design}

The baseline model has one lag and one observed covariate:
\begin{equation}
\label{eq:sim_dgp}
y_{it}
=
a_{it}y_{i,t-1}
+
\beta_{it}x_{it}
+
H_{0,it}
+
u_{it},
\qquad
i=1,\ldots,N.
\end{equation}
The process is generated for $t=-49,\ldots,T$, with
$$
y_{i,-50}=0,
\qquad
x_{i,-50}=0,
\qquad
g_{a,-50}=g_{b,-50}=g_{h,-50}=f_{x,-50}=0,
$$
and the first 50 periods are discarded. To keep the coefficient matrices uniformly bounded and compatible with the maintained incoherence restrictions, the loading perturbations and common-factor innovations below are independent standardized uniform variables on $[-\sqrt 3,\sqrt 3]$. Gaussian draws are retained for idiosyncratic disturbances.

For $m\in\{a,b\}$,
$$
g_{m,t}
=
\rho_g g_{m,t-1}
+
(1-\rho_g^2)^{1/2}v_{m,t},
\qquad
f_{m,t}=\mu_{f,m}+\kappa_{f,m}g_{m,t},
$$
where $v_{m,t}\stackrel{\mathrm{iid}}{\sim}U[-\sqrt 3,\sqrt 3]$ and $\rho_g=0.5$. For DGPs~1--4, set
$$
\lambda_{a,i}=1+0.1z_{a,i},
\qquad
\lambda_{b,i}=1+0.4z_{b,i},
$$
where $z_{a,i}$ and $z_{b,i}$ are independent standardized uniform variables. Use $\mu_{f,a}=0.5$, $\kappa_{f,a}=0.1$, $\mu_{f,b}=0.6$, and $\kappa_{f,b}=0.15$. Since the bounded AR(1) factor recursion above implies $|g_{b,t}|\le 3$, this choice gives
$$
0.15
\le
f_{b,t}
=
0.6+0.15g_{b,t}
\le
1.05,
$$
so the slope-factor time direction is uniformly separated from zero. The raw autoregressive and slope matrices are
$$
a_{it}^{\mathrm{raw}}=\lambda_{a,i}f_{a,t},
\qquad
\beta_{it}=\lambda_{b,i}f_{b,t}.
$$
Dynamic stability is imposed by
$$
a_{it}=c_a a_{it}^{\mathrm{raw}},
\qquad
c_a=\min\left\{1,\frac{0.85}{\max_{i,t}|a_{it}^{\mathrm{raw}}|}\right\}.
$$

Let $\lambda_{h,i}$ and the innovations in $g_{h,t}$ be independent standardized uniform variables, and generate
$$
g_{h,t}=\rho_g g_{h,t-1}+(1-\rho_g^2)^{1/2}v_{h,t},
\qquad
H_{it}^{\mathrm{raw}}=\lambda_{h,i}g_{h,t}.
$$
For each DGP and panel size, define
$$
H_{0,it}=c_\xi c_H H_{it}^{\mathrm{raw}},
\qquad
u_{it}=c_\xi\widetilde u_{it}.
$$
Thus, \eqref{eq:sim_dgp} exactly matches the maintained model: $\bH_0=(H_{0,it})$ is the true interactive-effect matrix and $u_{it}$ is the unit-specific disturbance. In the baseline design, $\bA_0$, $\bB_0$, and $\bH_0$ all have rank one.

\paragraph{DGP 1: independent heteroskedastic disturbances.}
Let $\widetilde u_{it}=\sigma_i\varepsilon_{it}$, where $\varepsilon_{it}\stackrel{\mathrm{iid}}{\sim}N(0,1)$ and $\sigma_i^2\stackrel{\mathrm{iid}}{\sim}U(0.5,1.5)$.

\paragraph{DGP 2: spatially dependent heteroskedastic disturbances.}
Units lie on a one-dimensional lattice with $d_N(i,j)=|i-j|$. Draw $\varepsilon_{it}\stackrel{\mathrm{iid}}{\sim}N(0,1)$ and $\sigma_i^2\stackrel{\mathrm{iid}}{\sim}U(0.5,1.5)$. For every time $t$, generate
$$
z_{1t}=\varepsilon_{1t},
\qquad
z_{it}=\rho_s z_{i-1,t}+(1-\rho_s^2)^{1/2}\varepsilon_{it},
\qquad
\rho_s=0.5,
$$
and set $\widetilde u_{it}=\sigma_i z_{it}$. Conditional on the unit variances,
$$
\operatorname{Cov}(\widetilde u_{it},\widetilde u_{jt})
=
\sigma_i\sigma_j\rho_s^{|i-j|}.
$$
The recursive construction has linear cost in $N$ and avoids a dense covariance square root.

For DGPs~1 and~2, generate
\begin{equation}
\label{eq:sim_x_exogenous}
x_{it}
=
\rho_xx_{i,t-1}
+
\delta_x\lambda_{x,i}f_{x,t}
+
(1-\rho_x^2)^{1/2}e_{it},
\end{equation}
where $\rho_x=0.5$, $\delta_x=0.5$, $\lambda_{x,i}\stackrel{\mathrm{iid}}{\sim}U[-\sqrt 3,\sqrt 3]$, $\sigma_{e,i}^2\stackrel{\mathrm{iid}}{\sim}U(0.5,1.5)$, and $e_{it}\sim N(0,\sigma_{e,i}^2)$. The common covariate factor satisfies
$$
f_{x,t}=\rho_{fx}f_{x,t-1}+(1-\rho_{fx}^2)^{1/2}v_{x,t},
\qquad
\rho_{fx}=0.5,
$$
where $v_{x,t}\stackrel{\mathrm{iid}}{\sim}U[-\sqrt 3,\sqrt 3]$.

\paragraph{DGP 3: predetermined covariates.}
DGP~3 retains the spatial disturbance from DGP~2 and generates
\begin{equation}
\label{eq:sim_x_predetermined}
x_{it}
=
\rho_xx_{i,t-1}
+
\delta_x\lambda_{x,i}f_{x,t}
+
\eta_x\widetilde u_{i,t-1}
+
(1-\rho_x^2)^{1/2}e_{it},
\end{equation}
where $\eta_x=0.3$ and the other terms follow \eqref{eq:sim_x_exogenous}. Thus, $x_{it}$ responds to the lagged disturbance but not to the current disturbance, matching the maintained predeterminedness restriction.

\paragraph{DGP 4: predetermined covariates with group heterogeneity.}
DGP~4 retains the disturbance and covariate processes from DGP~3 and changes only the coefficient loadings. Divide the units into two equal groups,
$$
G_{1,N}=\{1,\ldots,N/2\},
\qquad
G_{2,N}=\{N/2+1,\ldots,N\}.
$$
For $m\in\{a,b\}$, set
$$
\lambda_{m,i}
=
\mu_{\lambda,m,1}\mathbf 1\{i\in G_{1,N}\}
+
\mu_{\lambda,m,2}\mathbf 1\{i\in G_{2,N}\}
+
\sigma_{\lambda,m}z_{m,i},
$$
where $z_{m,i}\stackrel{\mathrm{iid}}{\sim}U[-\sqrt 3,\sqrt 3]$. The baseline values are
$$
\mu_{\lambda,a,1}=0.9,
\quad
\mu_{\lambda,a,2}=1.1,
\quad
\sigma_{\lambda,a}=0.08,
$$
and
$$
\mu_{\lambda,b,1}=0.8,
\quad
\mu_{\lambda,b,2}=1.2,
\quad
\sigma_{\lambda,b}=0.25.
$$
Construct $a_{it}^{\mathrm{raw}}=\lambda_{a,i}f_{a,t}$, apply the same common stability scaling $c_a$, and set $\beta_{it}=\lambda_{b,i}f_{b,t}$. Because the loading vector remains multiplied by one time factor, all three coefficient matrices remain rank one. DGP~4 therefore adds systematic group heterogeneity without changing the rank vector or the spatial and predetermined structure inherited from DGP~3.

For the Monte Carlo verification of the conditional assumptions, let $\mathscr D_N$ denote the sigma-field generated by all time-invariant unit-level design draws, including the coefficient loadings $\{\lambda_{a,i},\lambda_{b,i},\lambda_{h,i}\}_{i\le N}$, the covariate loadings $\{\lambda_{x,i}\}_{i\le N}$, the heteroskedastic scales $\{\sigma_i,\sigma_{e,i}\}_{i\le N}$, and, in DGP~4, the group-specific loading perturbations. We use the enlarged conditioning field
$$
\mathscr C^{\mathrm{MC}}
=
\sigma\{g_{a,t},g_{b,t},g_{h,t},f_{x,t}:t=-49,\ldots,T\}
\vee
\mathscr D_N.
$$
Conditional on $\mathscr C^{\mathrm{MC}}$, the true coefficient matrices and the time-invariant heteroskedastic scales are fixed, while the time-varying idiosyncratic Gaussian innovations remain random. The current disturbances have conditional mean zero, and DGPs~2--4 retain the stated spatial weak-dependence structure. In verifying the assumptions for the Monte Carlo designs, $\mathscr C^{\mathrm{MC}}$ plays the role of the conditioning sigma-field $\mathscr C$ used in the theory.

\subsection{Calibration, targets, and reported quantities}

A separate ex-ante calibration stage is completed before any reported Monte Carlo replication is generated. First, set $\pi_H=0.3$ and choose $c_H$ from the population moments
$$
c_H^2
=
\frac{\pi_H}{1-\pi_H}
\frac{\operatorname{Var}(\widetilde u_{it})}
{\operatorname{Var}(H_{it}^{\mathrm{raw}})}.
$$
The variances in this expression are population quantities implied by the primitive DGP distributions rather than realized sample variances. Thus, $c_H$ is a deterministic design-specific calibration constant fixed before the reported replications are generated.

Second, for every DGP and panel size, select $c_\xi$ in a separate calibration-only experiment with common random numbers so that the average pooled coefficient of determination is $0.65$, where the fitted value includes the lagged-outcome term, the observed-covariate term, and $H_{0,it}$. The calibration draws are independent of and never reused in the reported Monte Carlo replications. Once calibration is complete, all values of $c_H$ and $c_\xi$ are frozen and stored, and the reported replications read these constants without recalibration.

The resulting coefficient supports are checked analytically before the reported simulations. A single coefficient bound $B$ is then fixed for all maintained designs so that every true coefficient matrix lies at least a fixed margin $c_B>0$ inside the estimation set. The baseline implementation uses $B=10$ and $c_B=1$ after verification of the final calibrated coefficient envelope.

DGP~4 additionally reports the realized true group means and $G_{2,N}$-minus-$G_{1,N}$ differences for both $a_{it}$ and $\beta_{it}$, at the target time and after averaging over time. These truths are computed after applying the common autoregressive stability scaling $c_a$. The loading means and perturbation scales stated above are fixed ex ante and are not adjusted across panel sizes or after inspecting any Monte Carlo result. In particular, the autoregressive loading means remain $0.9$ and $1.1$, and the slope loading means remain $0.8$ and $1.2$. The realized post-stability group differences are reported as features of the generated DGP rather than targeted by an additional calibration step.

The balanced-panel sequence uses $N=T\in\{50,100,200,400\}$. Rectangular panels are included in the appendix to assess unequal growth of $N$ and $T$. The production design uses 1,000 replications per cell when computationally feasible; any reduced count for the largest cells is fixed before results are examined and reported together with Monte Carlo standard errors.

DGP~1 uses the diagonal variance estimator. DGPs~2--4 use the positive-semidefinite Bartlett weights
$$
w_{ik,N}
=
\left(1-\frac{|i-k|}{h_N+1}\right)_+,
\qquad
h_N=\left\lceil c_{\mathrm{sp}}\log(NT)\right\rceil,
$$
with baseline $c_{\mathrm{sp}}=1$ and prespecified sensitivity values in the appendix.

Let $i_0=\lfloor N/4\rfloor$ and $t_0=\lfloor T/2\rfloor$. The main parameters of interest are:
\begin{enumerate}[label=(\roman*),leftmargin=2em]
\item the individual entries $a_{i_0t_0}$ and $\beta_{i_0t_0}$;
\item the cross-sectional means of $a_{it_0}$ and $\beta_{it_0}$;
\item the group-specific fixed-time means for DGPs~1--4 and, for DGP~4, the $G_{2,N}$-minus-$G_{1,N}$ fixed-time contrasts; the analogous fixed-time contrasts for DGPs~1--3 are reported only as weak-target diagnostics outside the uniform projection condition in \Cref{a:target};
\item the full-panel means of $\bA_0$ and $\bB_0$;
\item the group-specific time-averaged means and the $G_{2,N}$-minus-$G_{1,N}$ time-averaged contrasts.
\end{enumerate}
The first three parameter classes use the full-panel plug-in estimator whenever the target regularity conditions in \Cref{a:target} hold. The last two use the two-way split-panel correction. Group-specific means are theorem-covered targets in all four DGPs. The fixed-time $G_{2,N}$-minus-$G_{1,N}$ contrast is a headline theorem-covered target only in DGP~4, where the loading design supplies systematic group separation. The corresponding fixed-time contrasts in DGPs~1--3 are retained, if reported, only as weak-target stress diagnostics and are not used as evidence for the uniform inference theorem. DGP~4 provides systematically nonzero group means and contrasts and therefore evaluates both confidence-interval coverage and power against zero.

For every parameter and panel cell, report the exact true value, mean estimate, bias, RMSE, Monte Carlo standard deviation, mean estimated standard error, the ratio of mean estimated standard error to Monte Carlo standard deviation, empirical size of the two-sided $5\%$ test centered at the truth, power against zero when the truth is nonzero, $95\%$ coverage, Monte Carlo standard errors for size and coverage, and the number and reasons for failed replications.

The main inference designs have true rank vector $(1,1,1)$. A separate rank-selection experiment applies each of the four disturbance and covariate designs to the true rank vectors $(1,1,1)$, $(2,1,1)$, and $(1,0,2)$. Each nonzero rank-two coefficient matrix is constructed from two independent bounded loading--factor terms, with the same bounded primitive distributions used in the baseline design, followed by a deterministic common rescaling that preserves the maintained coefficient envelope and, for autoregressive matrices, dynamic stability. Rank-zero matrices are set identically equal to zero. For every replication, rank selection uses the spectral pilot, the normalized coefficient-matrix spectra, the deterministic ridge $a_{NT}=1/\log(NT)$, and one final selected-rank fit; no information criterion or rank-selection sensitivity multiplier is used.

For the $(1,0,2)$ design, the slope matrix is identically zero. Because the common scaling $c_\xi$ does not identify the pooled coefficient of determination in this design, we fix $c_\xi=1$ and report the induced pooled $R^2$ rather than imposing the baseline $0.65$ target. The exact loading--factor normalization and deterministic rescaling used for the rank-two matrices are reported in \Cref{app:simulation}.

For the Monte Carlo rank selector, the reportable rank caps are
$$
(\bar r_A,\bar r_B,\bar r_H)=(3,3,3),
$$
while the pilot rank bounds are $(4,4,4)$. The reportable caps are fixed before the reported replications are generated and contain every maintained true rank vector. The extra pilot rank is used only to compute the singular value immediately beyond each reportable cap.


\paragraph{Rank-selection reporting.}
For each DGP, rank configuration, and panel size, report attempted replications, numerically acceptable spectral pilots, exact recovery, underselection only, overselection only, mixed rank errors, unresolved selections, selections at a reportable cap, and the complete selected-rank distribution. The appendix reports coefficient-matrix-specific recovery, normalized pilot singular values, every ridge ratio, and the difference between the smallest and second-smallest ratios. Balanced panels are complemented by rectangular panels to assess unequal growth of $N$ and $T$.

\paragraph{Monte Carlo reporting.}
The main fixed-rank results are reported separately for DGPs~1--4 in \Cref{tab:mc_dgp1,tab:mc_dgp2,tab:mc_dgp3,tab:mc_dgp4}; rank selection is reported separately in \Cref{tab:mc_rank}. Every attempted target record remains in the raw records. Let $R_{\mathrm{attempted}}$, $R_{\mathrm{point}}$, and $R_{\mathrm{inference}}$ denote the numbers attempted, yielding a valid finite point estimate, and yielding valid primary inference, respectively, with $R_{\mathrm{inference}}\le R_{\mathrm{point}}\le R_{\mathrm{attempted}}$. For replication $r$, write $e_r=\widehat\phi_r-\phi_r$ for the replication-specific estimation error. The reported bias and RMSE are the mean and root mean square of $e_r$ over point-valid replications. Because the exact simulated truth may vary across replications, ``MC SD'' is the standard deviation of $e_r$, rather than of the raw estimates $\widehat\phi_r$. ``Mean SE'' is the average estimated standard error over inference-valid replications, and ``SE/SD'' divides this average by the standard deviation of $e_r$ computed over the same inference-valid sample. ``Cov.'' is empirical 95\% coverage, while ``Rej.'' is the rejection probability of the two-sided test of $H_0:\phi=0$. ``Point/Inf. ret.'' reports the fractions of attempted target records yielding a valid point estimate and valid primary inference. Numerical failures and theorem-driven inference suppressions are assigned mutually exclusive primary statuses and reported separately rather than silently deleted; no denominator is determined by an undocumented missing-value filter. Finite extreme estimates are never trimmed or winsorized. As a tail diagnostic, we flag inference-valid standardized errors satisfying $|e_r/\widehat s_r|>8$, retain them in every headline calculation, and list the corresponding replication identifiers in the appendix.

\input{tables/mc/mc_dgp1}
\input{tables/mc/mc_dgp2}
\input{tables/mc/mc_dgp3}
\input{tables/mc/mc_dgp4}
\placeholdertable{Blockwise ridge-ratio rank recovery}{tab:mc_rank}{Revision-10 shell with columns: DGP; $N$; $T$; true rank vector; attempted replications; numerically acceptable spectral pilots; exact recovery; under only; over only; mixed; numerically unresolved; and modal selected rank. Exact means $\widehat\br=\br_0$. Under only requires every coordinate weakly below truth and at least one strictly below; over only is defined symmetrically; mixed requires at least one underselected and one overselected coordinate. Detailed block-level cap selections are reserved for the appendix. No numerical entries are populated in RR3.1.}

\input{tables/mc/mc_retention}
\input{tables/mc/mc_optimization}





\paragraph{Graphical summaries.}
The production figures are prespecified before the final replications are examined. Figure~\ref{fig:mc_bias_rmse} shows convergence of bias and RMSE; Figure~\ref{fig:mc_coverage} shows coverage and standard-error accuracy; Figure~\ref{fig:mc_rank_recovery} shows ridge-ratio rank-recovery frequencies and ratio separation; Figure~\ref{fig:mc_power} reports DGP~4 power; and Figure~\ref{fig:mc_retention} reports retention and numerical-failure rates. All Monte Carlo uncertainty bands use the relevant denominator for the statistic being plotted.

\input{figures/mc/mc_bias_rmse}

\input{figures/mc/mc_coverage}

\input{figures/mc/mc_rank_recovery}

\input{figures/mc/mc_power}

\input{figures/mc/mc_retention}


\FloatBarrier

%================================================
\section{Empirical Application}
\label{sec:application}

The empirical section contains two applications---regional unemployment and house-price momentum---and will recompute all estimates with the full-panel estimator.

\subsection{Regional unemployment}

The planned application uses monthly metropolitan unemployment rates, deseasonalized metro by metro. The working data build contains 315 metropolitan areas over 240 months in the main specification. The interactive-effects matrix absorbs the common cycle, while the autoregressive coefficient matrices describe remaining local persistence. Predetermined controls include lagged metropolitan population growth and real-GDP growth.

The main targets will be the full-panel persistence mean, high- and low-unemployment group means, their contrast, selected individual entries, and time paths of cross-sectional mean persistence. A time-invariant high/low split will be defined before estimation. Full-period means and contrasts will use the two-way correction; fixed-time targets and individual entries will use the full-panel plug-in when the target-rate condition holds.

\placeholdertable{Idiosyncratic persistence in metropolitan unemployment}{tab:emp_unemployment}{Placeholder for selected ranks, full and group means, the high-minus-low contrast, bias-corrected estimates, spatial standard errors, confidence intervals, and target-specific diagnostics.}

\subsection{House-price momentum}

The planned housing application stacks the top and bottom house-price tiers for each metropolitan area as separate units. The working data build contains 610 metro-tier units over 315 months. The outcome is monthly log price growth. The main targets are the mean first- and second-lag coefficients, lag-sum persistence, top- and bottom-tier means, their contrast, selected individual entries, and time paths of mean momentum. Predetermined controls include lagged building-permit growth, population growth, and real-GDP growth.

\placeholdertable{House-price dynamics and price-tier contrasts}{tab:emp_housing}{Placeholder for selected ranks, lag means, lag-sum persistence, top and bottom tier means, the top-minus-bottom contrast, bias-corrected estimates, spatial standard errors, and confidence intervals.}

\subsection{Spatial inference and computation}

When metropolitan coordinates are available, the primary dependence-robust calculation will use great-circle distance and several prespecified spatial distance cutoffs. The diagonal calculation will be reported as a benchmark. A known-cluster calculation may be reported when a credible cluster partition is available. Period clustering is not treated as a substitute for the spatial theorem unless additional conditions are verified.

\placeholdertable{Geographic spatial HAC sensitivity}{tab:emp_spatial}{Placeholder for diagonal and spatial HAC standard errors at several geographic distance cutoffs for the main unemployment and housing targets.}

\placeholdertable{Rank, split-panel, and optimization diagnostics}{tab:emp_diagnostics}{Placeholder for spectral pilot singular values, normalized block spectra, ridge-ratio margins, selected ranks, full- and split-panel Gram-matrix eigenvalues, objective tolerances, Riesz iterations, coefficient-envelope checks, and spatial-weight diagnostics.}

\placeholderfigure{Estimated coefficient paths and selected entrywise estimates}{fig:emp_paths}{Placeholder for time paths of cross-sectional mean autoregressive coefficients and selected entrywise confidence intervals.}

\placeholderfigure{Geographic distribution of selected coefficient summaries}{fig:emp_maps}{Placeholder for metropolitan maps, with uncertainty or reliability indicators reported separately from point estimates.}

\FloatBarrier

%================================================
\section{Conclusion}
\label{sec:conclusion}

This paper develops full-panel estimation and inference for dynamic panels with several separately ranked coefficient matrices. One spectral pilot and deterministic ridge ratios select each matrix rank, while all reported coefficients are obtained from a separate unpenalized joint fixed-rank estimator. The analysis aligns and expands the estimated singular spaces and uses a target-specific Riesz representer to translate estimation error in the autoregressive, slope, and interactive-effect matrices into first-order uncertainty. The framework allows predetermined covariates, conditioning on common shocks, dynamic feedback through lagged outcomes, and spatially dependent unit-specific innovations.

The inferential procedure depends on the target. Individual coefficients, fixed unit--time lag sums, and suitable fixed-time means and contrasts use one full-panel fit because the leading loading- and factor-estimation biases are smaller than their standard errors. Averages whose weights extend over many units and periods receive a fixed two-way split-panel correction that removes both leading bias terms. Rank selection is performed once on the full panel. Strong singular values and the pilot operator rate separate the true ranks from the smaller singular values beyond the truth, while the deterministic ridge $1/\log(NT)$ stabilizes the ratios without a data-selected penalty.

The theory assumes exact fixed ranks and consistent selection of each coefficient-matrix rank. Rank selection requires a localized enlargement of the prediction-identification domain for the spectral pilot, but it does not strengthen the temporal or spatial dependence assumptions. The confidence intervals are not presented as approximate-low-rank or rank-misspecification-robust procedures. Growing ranks, rank-robust inference, and semiparametric efficiency under conditional spatial dependence are left for future work.

The Monte Carlo and empirical sections are designed to evaluate finite-sample bias correction, spatial inference, rank recovery, numerical reliability, and computation. Their numerical tables and figures will be populated after implementation is complete.

%================================================
\begin{singlespace}
\begin{thebibliography}{99}

\bibitem[Ahn and Horenstein(2013)]{ahnhorenstein2013}
Ahn, S.~C., and Horenstein, A.~R. (2013).
Eigenvalue ratio test for the number of factors.
\emph{Econometrica}, 81(3), 1203--1227.
\doilink{10.3982/ECTA8968}.

\bibitem[Anderson and Hsiao(1981)]{andersonhsiao1981}
Anderson, T.~W., and Hsiao, C. (1981).
Estimation of dynamic models with error components.
\emph{Journal of the American Statistical Association}, 76(375), 598--606.
\doilink{10.1080/01621459.1981.10477691}.

\bibitem[Anderson and Hsiao(1982)]{andersonhsiao1982}
Anderson, T.~W., and Hsiao, C. (1982).
Formulation and estimation of dynamic models using panel data.
\emph{Journal of Econometrics}, 18(1), 47--82.
\doilink{10.1016/0304-4076(82)90095-1}.

\bibitem[Arellano and Bond(1991)]{arellanobond1991}
Arellano, M., and Bond, S. (1991).
Some tests of specification for panel data: Monte Carlo evidence and an application to employment equations.
\emph{Review of Economic Studies}, 58(2), 277--297.
\doilink{10.2307/2297968}.

\bibitem[Athey et~al.(2021)]{athey2021matrix}
Athey, S., Bayati, M., Doudchenko, N., Imbens, G.~W., and Khosravi, K. (2021).
Matrix completion methods for causal panel data models.
\emph{Journal of the American Statistical Association}, 116(536), 1716--1730.
\doilink{10.1080/01621459.2021.1891924}.

\bibitem[Bai(2003)]{bai2003}
Bai, J. (2003).
Inferential theory for factor models of large dimensions.
\emph{Econometrica}, 71(1), 135--171.
\doilink{10.1111/1468-0262.00392}.

\bibitem[Bai and Ng(2002)]{bai2002}
Bai, J., and Ng, S. (2002).
Determining the number of factors in approximate factor models.
\emph{Econometrica}, 70(1), 191--221.
\doilink{10.1111/1468-0262.00273}.

\bibitem[Bai and Ng(2021)]{bai_ng2021matrix}
Bai, J., and Ng, S. (2021).
Matrix completion, counterfactuals, and factor analysis of missing data.
\emph{Journal of the American Statistical Association}, 116(536), 1746--1763.
\doilink{10.1080/01621459.2021.1967163}.


\bibitem[Bai et~al.(2021)]{baich2021}
Bai, J., Choi, S.~H., and Liao, Y. (2021).
Feasible generalized least squares for panel data with cross-sectional and serial correlations.
\emph{Empirical Economics}, 60(1), 309--326.
\doilink{10.1007/s00181-020-01977-2}.

\bibitem[Bonhomme and Manresa(2015)]{bonhommemanresa2015}
Bonhomme, S., and Manresa, E. (2015).
Grouped patterns of heterogeneity in panel data.
\emph{Econometrica}, 83(3), 1147--1184.
\doilink{10.3982/ECTA11319}.

\bibitem[Bradley(1983)]{bradley1983}
Bradley, R.~C. (1983).
Approximation theorems for strongly mixing random variables.
\emph{Michigan Mathematical Journal}, 30(1), 69--81.
\doilink{10.1307/mmj/1029002789}.

\bibitem[Chernozhukov et~al.(2023)]{chernozhukov2023}
Chernozhukov, V., Hansen, C., Liao, Y., and Zhu, Y. (2023).
Inference for low-rank models.
\emph{Annals of Statistics}, 51(3), 1309--1330.
\doilink{10.1214/23-AOS2293}.

\bibitem[Choi et~al.(2024)]{choi2024}
Choi, J., Kwon, H., and Liao, Y. (2024).
Inference for low-rank completion without sample splitting with application to treatment effect estimation.
\emph{Journal of Econometrics}, 240(1), Article 105682.
\doilink{10.1016/j.jeconom.2024.105682}.

\bibitem[Choi et~al.(2026)]{choi2026rank}
Choi, J., Kwon, H., and Liao, Y. (2026).
Inference for low-rank models without estimating the rank.
\emph{Journal of the American Statistical Association}, 121(553), 655--666.
\doilink{10.1080/01621459.2025.2538272}.

\bibitem[Barigozzi et~al.(2026)]{barigozzietal2026}
Barigozzi, M., He, Y., Li, L., and Trapani, L. (2026).
Statistical inference for large-dimensional tensor factor model by iterative projections.
\emph{Journal of Multivariate Analysis}, 214, Article 105616.
\doilink{10.1016/j.jmva.2026.105616}.

\bibitem[Pu et~al.(2025)]{puetal2025}
Pu, D., Fang, K., Lan, W., Yu, J., and Zhang, Q. (2025).
Reduced rank spatio-temporal models.
\emph{Journal of Business \& Economic Statistics}, 43(1), 98--109.
\doilink{10.1080/07350015.2024.2326142}.

\bibitem[Conley(1999)]{conley1999gmm}
Conley, T.~G. (1999).
GMM estimation with cross sectional dependence.
\emph{Journal of Econometrics}, 92(1), 1--45.
\doilink{10.1016/S0304-4076(98)00084-0}.



\bibitem[Jenish and Prucha(2012)]{jenishprucha2012}
Jenish, N., and Prucha, I.~R. (2012).
On spatial processes and asymptotic inference under near-epoch dependence.
\emph{Journal of Econometrics}, 170(1), 178--190.
\doilink{10.1016/j.jeconom.2012.05.022}.

\bibitem[Moon and Weidner(2015)]{moonweidner2015}
Moon, H.~R., and Weidner, M. (2015).
Linear regression for panel with unknown number of factors as interactive fixed effects.
\emph{Econometrica}, 83(4), 1543--1579.
\doilink{10.3982/ECTA9382}.



\bibitem[Moon and Weidner(2017)]{moonweidner2017}
Moon, H.~R., and Weidner, M. (2017).
Dynamic linear panel regression models with interactive fixed effects.
\emph{Econometric Theory}, 33(1), 158--195.
\doilink{10.1017/S0266466615000328}.

\bibitem[Kuersteiner and Prucha(2020)]{kuersteinerprucha2020}
Kuersteiner, G.~M., and Prucha, I.~R. (2020).
Dynamic spatial panel models: Networks, common shocks, and sequential exogeneity.
\emph{Econometrica}, 88(5), 2109--2146.
\doilink{10.3982/ECTA13660}.

\bibitem[Peng et~al.(2025)]{pengetal2025}
Peng, B., Su, L., Westerlund, J., and Yang, Y. (2025).
Interactive effects panel data models with general factors and regressors.
\emph{Econometric Theory}, 41(2), 472--488.
\doilink{10.1017/S0266466623000270}.

\bibitem[Pesaran(2006)]{pesaran2006}
Pesaran, M.~H. (2006).
Estimation and inference in large heterogeneous panels with a multifactor error structure.
\emph{Econometrica}, 74(4), 967--1012.
\doilink{10.1111/j.1468-0262.2006.00692.x}.

\bibitem[Pesaran and Smith(1995)]{pesaransmith1995}
Pesaran, M.~H., and Smith, R. (1995).
Estimating long-run relationships from dynamic heterogeneous panels.
\emph{Journal of Econometrics}, 68(1), 79--113.
\doilink{10.1016/0304-4076(94)01644-F}.

\bibitem[Swamy(1970)]{swamy1970}
Swamy, P.~A.~V.~B. (1970).
Efficient inference in a random coefficient regression model.
\emph{Econometrica}, 38(2), 311--323.
\doilink{10.2307/1913012}.

\bibitem[Tropp(2012)]{tropp2012}
Tropp, J.~A. (2012).
User-friendly tail bounds for sums of random matrices.
\emph{Foundations of Computational Mathematics}, 12(4), 389--434.
\doilink{10.1007/s10208-011-9099-z}.

\bibitem[Wang et~al.(2022)]{wangsu2022}
Wang, Y., Su, L., and Zhang, Y. (2022).
Low-rank panel quantile regression: Estimation and inference.
\emph{arXiv preprint} \doilink{10.48550/arXiv.2210.11062}.

\bibitem[Wedin(1972)]{wedin1972}
Wedin, P.-\AA. (1972).
Perturbation bounds in connection with singular value decomposition.
\emph{BIT Numerical Mathematics}, 12, 99--111.
\doilink{10.1007/BF01932678}.

\bibitem[Yuan and Spindler(2025)]{yuanspindler2025}
Yuan, Z., and Spindler, M. (2025).
Bernstein-type inequalities and nonparametric estimation under near-epoch dependence.
\emph{Journal of Econometrics}, 251, Article 106054.
\doilink{10.1016/j.jeconom.2025.106054}.

\end{thebibliography}
\end{singlespace}

%================================================
\clearpage
\appendix


\section{Proofs and Supporting Results}
\label{app:proofs}

All proofs are carried out under the conditional law given $\mathscr C$. Conditional $O_p$ and $o_p$ statements are uniform over the regular common-shock realizations covered by the deterministic constants in the assumptions. Integrating a conditional conclusion over $\mathscr C$ gives its unconditional counterpart.

\begin{lemma}[Conditional use of random-field limit results]
\label{lem:conditional_limit_transfer}
Fix a regular realization $\mathscr C=c$ and let $\Pzero^c$ denote a regular conditional probability given $\mathscr C=c$.
\begin{enumerate}[label=(\roman*),leftmargin=2em]
\item If a cited theorem is an asymptotic weak-limit result and its assumptions hold under $\Pzero^c$, then it may be applied under that conditional law for almost every regular $c$. If the resulting conditional distribution functions converge pointwise to a continuous limit, Pólya's theorem gives uniform convergence in the argument for that fixed $c$, and bounded convergence over $c$ yields convergence in probability of the conditional distribution distance and the corresponding unconditional weak limit.
\item If a cited theorem is a finite-sample probability inequality whose constants depend only on deterministic moment, sampling-geometry, NED, and mixing envelopes, then the same inequality holds under $\Pzero^c$ with constants uniform over regular $c$. Integrating the conditional probability bound gives the corresponding unconditional bound.
\end{enumerate}
\end{lemma}

\begin{proof}
For almost every $c$, the regular conditional law $\Pzero^c$ is an ordinary probability measure. Once $c$ is fixed, the arrays of unit-specific innovations, localized covariates, and scores satisfy the same deterministic geometry, moment, NED, and mixing inequalities as in the maintained assumptions. The random-field central limit theorem of \citet{jenishprucha2012} is therefore applied separately under each such $\Pzero^c$; no uniform Berry--Esseen bound is invoked. Because the Gaussian limit is continuous, conditional weak convergence for almost every $c$ implies uniform convergence of the corresponding conditional distribution functions for that fixed $c$, and the distance is bounded by one. Bounded convergence over the distribution of $\mathscr C$ then gives convergence in probability of the conditional distance and the unconditional conclusion.

For finite-sample concentration inequalities the transfer is more literal: if an external inequality is applied after conditioning on $\mathscr C=c$, its own sampling, weighting, moment, NED, and mixing hypotheses must first be verified under $\Pzero^c$, and any constants that are to be uniform in $c$ must depend only on the deterministic envelopes in the maintained assumptions. The Bernstein inequality of \citet{yuanspindler2025} is a useful irregular-spacing NED benchmark, but its published Theorem~1 concerns a scalar average under its own domain-expanding-infill geometry. The concentration argument needed here is instead proved directly in \Cref{lem:weighted_concentration} from the stronger setwise conditional mixing coefficient in \Cref{a:ned}, exponential local approximation, and the fixed-dimensional unit--time geometry. Thus, no unverified weighted extension of the Yuan--Spindler theorem is used below.
\end{proof}

For compactness in the appendix only, write
$$
Z_{it}=\Psi_{0,it}u_{it},
\qquad
\widehat Z_{it}=\widehat\Psi_{it}\widehat u_{it}.
$$
For the split-panel correction, write
$$
Z_{it}^{\mathrm{corr}}
=
u_{it}
\left[
3\Psi_{0,it}
-\Psi_{0,it}^{T}
-\Psi_{0,it}^{N}
\right],
$$
Equivalently, $Z_{it}^{\mathrm{corr}}$ is the bracketed Riesz weight multiplied by $u_{it}$. The feasible analogue combines the full- and split-panel residual scores as in \eqref{eq:corrected_spatial_variance}. These abbreviations are used only to keep the proofs readable.

\subsection{Conditional moments and local approximation}
\label{app:localization}

This subsection establishes the moment and localization results used in every later proof. The dynamic equation makes a current fitted value depend on lagged outcomes, and the covariates may depend on a long history of nearby unit-specific innovations. We first show that stability keeps the required moments bounded. We then replace each fitted value or score by a version that depends only on innovations from a finite time window and a finite spatial neighborhood. Conditional alpha mixing can be applied after this replacement because the innovation sets associated with sufficiently separated observations are then disjoint.

For $p\ge1$ and a random variable $X$, define its conditional $L_p$ norm by
$$
\|X\|_{p,\mathscr C}
=
\{\Ezero(|X|^p\mid\mathscr C)\}^{1/p}.
$$
Thus, $\|X\|_{p,\mathscr C}$ measures the size of $X$ under the conditional probability law obtained after the common shocks are held fixed. All bounds below are uniform over the regular common-shock realizations covered by the assumptions.

Recall from \Cref{a:ned} that $\mathscr E_{it}(h,r)$ contains unit-specific innovations from the previous $h$ time periods and from units within distance $r$ of unit $i$, with current covariate innovations included and current outcome innovations excluded. The approximation error has three sources: propagation through the dynamic equation, dependence of the covariates on remote past innovations, and dependence of the covariates on spatially distant innovations. For $h\ge3P$ and $r\ge0$, collect these three errors in
$$
\chi(h,r)
=
\underbrace{\rho^{\lfloor h/3\rfloor}}_{\text{dynamic propagation}}
+
\underbrace{\psi_T(\lfloor h/3\rfloor)}_{\text{remote time innovations}}
+
\underbrace{\psi_S(\lfloor r/3\rfloor)}_{\text{remote spatial innovations}}.
$$
Dividing both localization radii by three is only a proof device: it leaves separate margins for the outcome lags, the dynamic recursion, the covariate approximation, and the mixing separation. Any fixed division with sufficiently large margins would give the same asymptotic conclusion. By stability and the NED assumption, $\chi(h,r)\to0$ as $h,r\to\infty$.

The field $\mathscr E_{it}(h,r)$ is appropriate for approximating predetermined covariates because it excludes the current outcome innovation. A score such as $u_{it}\mathcal Y_{it}(\bDelta)$ contains the current innovation explicitly. For score approximations only, define the enlarged local information set
$$
\mathscr E_{it}^{+}(h,r)
=
\mathscr E_{it}(h,r)
\vee
\sigma\{u_{jt}:d_N(i,j)\le r\}.
$$
The added sigma-field contains current outcome innovations for unit $i$ and its spatial neighbors. This enlargement does not alter the timing assumption for the covariates; it is used only because the score itself contains $u_{it}$.

\begin{lemma}[Moments and local approximations]
\label{lem:localization}
Under \Cref{a:stab,a:ned,a:moments}, outcomes and their fixed number of lags have uniformly bounded conditional $8+\eta$ moments. For every $\ell=1,\ldots,P$, $y_{i,t-\ell}$ has an approximation measurable with respect to $\mathscr E_{it}(h,r)\vee\mathscr C$ whose conditional $8+\eta$ error is bounded by $C\chi(h,r)$. Consequently, for every fixed $\bDelta$,
$$
\left\|
\mathcal Y_{it}(\bDelta)
-
\mathcal Y_{it}^{[h,r]}(\bDelta)
\right\|_{8+\eta,\mathscr C}
\le
C\|\bDelta\|\chi(h,r),
$$
where $\mathcal Y_{it}^{[h,r]}(\bDelta)$ is measurable with respect to $\mathscr E_{it}(h,r)\vee\mathscr C$. The product
$$
u_{it}\mathcal Y_{it}(\bDelta)
$$
has an approximation measurable with respect to $\mathscr E_{it}^{+}(h,r)\vee\mathscr C$ with conditional $4+\eta/2$ error bounded by $C\|\bDelta\|\chi(h,r)$. Products used in the Gram-matrix and variance calculations satisfy the corresponding $2+\eta/4$ bounds.
\end{lemma}

\begin{proof}
Let
$$
\bm y_{it}
=
(y_{it},y_{i,t-1},\ldots,y_{i,t-P+1})'
$$
and let $\be_1$ be the first coordinate vector in $\R^P$. The dynamic model can be written as
$$
\bm y_{it}
=
\bC_{0,it}\bm y_{i,t-1}
+
\be_1
\left\{
\sum_{k=1}^{K}x_{it,k}\beta_{0,it,k}
+h_{0,it}
+u_{it}
\right\}.
$$
For $j\ge1$, define
$$
\bPhi_{i,t}(j)
=
\bC_{0,it}\bC_{0,i,t-1}\cdots\bC_{0,i,t-j+1},
\qquad
\bPhi_{i,t}(0)=\bm I_P.
$$
Iterating backward $m$ periods gives
\begin{align}
\bm y_{it}
={}&
\bPhi_{i,t}(m)\bm y_{i,t-m}
\nonumber\\
&+
\sum_{j=0}^{m-1}
\bPhi_{i,t}(j)\be_1
\left\{
\sum_{k=1}^{K}x_{i,t-j,k}\beta_{0,i,t-j,k}
+h_{0,i,t-j}
+u_{i,t-j}
\right\}.
\label{eq:companion_expansion}
\end{align}
By \eqref{eq:stability}, $\|\bPhi_{i,t}(j)\|_{\op}\le C\rho^j$. The coefficient entries are bounded, $P$ and $K$ are fixed, and \Cref{a:moments} bounds the moments of $x_{it,k}$ and $u_{it}$. Conditional Minkowski applied to \eqref{eq:companion_expansion}, followed by $m\to\infty$, yields
$$
\sup_{i,t}\|\bm y_{it}\|_{8+\eta,\mathscr C}\le C.
$$

Take $m=\lfloor h/3\rfloor$. Replace every retained covariate in \eqref{eq:companion_expansion} by its conditional expectation given the corresponding local innovation sigma-field. The omitted initial term is bounded by $C\rho^m$. The accumulated temporal and spatial covariate approximation errors are bounded by
$$
C\sum_{j=0}^{m-1}\rho^j
\{\psi_T(m)+\psi_S(r)\}
\le
C\{\psi_T(m)+\psi_S(r)\}.
$$
For $\ell\ge1$, every disturbance entering $y_{i,t-\ell}$ is dated no later than $t-1$, so the approximation is measurable with respect to $\mathscr E_{it}(h,r)\vee\mathscr C$. This proves the lagged-outcome statement.

The fitted-value map is a finite sum of bounded coefficient entries multiplied by the approximated lagged outcomes and covariates. Conditional Minkowski therefore gives the stated bound for $\mathcal Y_{it}(\bDelta)$. For products, write
$$
AB-A^{[h,r]}B^{[h,r]}
=
(A-A^{[h,r]})B
+
A^{[h,r]}(B-B^{[h,r]}).
$$
Conditional H\"older with the $8+\eta$ moment bound yields a $4+\eta/2$ product bound. Applying the same argument once more gives the $2+\eta/4$ bounds for quadratic products. The score approximation includes $u_{it}$ itself and is therefore measurable with respect to $\mathscr E_{it}^{+}(h,r)\vee\mathscr C$.
\end{proof}

For two locally measurable variables whose innovation neighborhoods are separated by distance $a$, the conditional mixing covariance inequality and \Cref{a:ned} imply
\begin{equation}
\label{eq:conditional_covariance_bound}
\left|
\operatorname{Cov}_0(X,Y\mid\mathscr C)
\right|
\le
C
\|X\|_{p,\mathscr C}
\|Y\|_{p,\mathscr C}
\bar\alpha(a)^{\eta/(8+\eta)}
\end{equation}
for the values of $p$ used below. Approximation errors add terms proportional to $\chi(h,r)$. Exponential decay and the polynomial neighborhood bound in \Cref{a:geometry} make the resulting covariance sums finite.

\subsection{Target directions and the population Riesz-score limit}
\label{app:population}

This subsection has three steps. First, it relates the size of the population Riesz representer to the direction $\bD$ of the parameter of interest. Second, it derives the variance rates and verifies that no single observation dominates the Riesz score. Third, it applies the conditional NED central limit theorem to prove \Cref{thm:population_clt}.

Let
$$
\Gamma_0(\bDelta,\bQ)
=
\Ezero\!\left[
\sum_{i,t}
\mathcal Y_{it}(\bDelta)
\mathcal Y_{it}(\bQ)
\middle|\mathscr C
\right].
$$
Because $\Gamma_0$ is uniformly positive definite on $\calT_0$, the Riesz equation has a unique solution. Taking $\bDelta=\bq_0$ in \eqref{eq:population_target_weight} gives
\begin{equation}
\label{eq:q_norm_identity}
\Gamma_0(\bq_0,\bq_0)
=
\ip{\calP_0\bD}{\bq_0}.
\end{equation}
The upper and lower Gram-matrix bounds therefore imply
\begin{equation}
\label{eq:q_target_norm_equivalence}
c\|\calP_0\bD\|
\le
\|\bq_0\|
\le
C\|\calP_0\bD\|.
\end{equation}

For use in both the score and recovery proofs, introduce one common set of local loading and factor coordinates. If $\bM_0=\bU_M\bSigma_M\bV_M'$ is a positive-rank coefficient matrix, take
$$
\bm\Lambda_{M,0}=\sqrt N\,\bU_M,
\qquad
\bm F_{M,0}=N^{-1/2}\bV_M\bSigma_M.
$$
Then $\bM_0=\bm\Lambda_{M,0}\bm F_{M,0}'$, $N^{-1}\bm\Lambda_{M,0}'\bm\Lambda_{M,0}=\bm I_{r_M}$, and the eigenvalues of $T^{-1}\bm F_{M,0}'\bm F_{M,0}=\bSigma_M^2/(NT)$ are bounded above and away from zero by \Cref{a:signal}. This gives the explicit map between the singular-value and loading--factor parameterizations used in the paper.

Let $z_{it,M}$ be the regressor multiplying matrix $M$:
$$
z_{it,A^{(\ell)}}=y_{i,t-\ell},
\qquad
z_{it,B^{(k)}}=x_{it,k},
\qquad
z_{it,H}=1.
$$
Stack the loading rows over the positive-rank matrices in $\bm\lambda_{0,i}$, stack their factor rows in $\bm f_{0,t}$, and set
$$
\mathbb Z_{it}
=
\operatorname{diag}\{z_{it,M}\bm I_{r_M}:r_M>0\}.
$$
After imposing the usual normalization restrictions that remove loading--factor rotations, a local direction is written as $\bm h=(\{\bm h_i^\lambda\}_i,\{\bm h_t^f\}_t)$ and its fitted-value derivative is
\begin{equation}
\label{eq:local_derivative}
\dot{\mathcal Y}_{it}(\bm h)
=
\frac{(\mathbb Z_{it}\bm f_{0,t})'\bm h_i^\lambda}{\sqrt T}
+
\frac{(\mathbb Z_{it}'\bm\lambda_{0,i})'\bm h_t^f}{\sqrt N}.
\end{equation}
Define the conditional local Gram operator by
$$
\mathbb G_0(\bm h,\bm q)
=
\Ezero\!\left[
\sum_{i,t}
\dot{\mathcal Y}_{it}(\bm h)
\dot{\mathcal Y}_{it}(\bm q)
\middle|\mathscr C
\right].
$$

\begin{lemma}[Rows of the inverse local Gram operator]
\label{lem:inverse_gram_rows}
Under \Cref{a:moments,a:signal,a:identification,a:gram}, $\mathbb G_0$ is invertible on the normalized local coordinate space. Uniformly for a fixed target unit $i$ and target time $t$,
\begin{equation}
\label{eq:inverse_gram_rows}
\begin{aligned}
&\|[\mathbb G_0^{-1}]_{\lambda_i,\lambda_i}\|_{\op}\le C,
&&\sum_{j\ne i}\|[\mathbb G_0^{-1}]_{\lambda_i,\lambda_j}\|_F^2\le C/N,
&&\sum_s\|[\mathbb G_0^{-1}]_{\lambda_i,f_s}\|_F^2\le C/N,\\
&\|[\mathbb G_0^{-1}]_{f_t,f_t}\|_{\op}\le C,
&&\sum_{s\ne t}\|[\mathbb G_0^{-1}]_{f_t,f_s}\|_F^2\le C/T,
&&\sum_j\|[\mathbb G_0^{-1}]_{f_t,\lambda_j}\|_F^2\le C/T.
\end{aligned}
\end{equation}
\end{lemma}

\begin{proof}
Partition $\mathbb G_0$ into loading and factor coordinates,
$$
\mathbb G_0
=
\begin{pmatrix}
G_{\lambda\lambda}&G_{\lambda f}\\
G_{f\lambda}&G_{ff}
\end{pmatrix}.
$$
The diagonal block associated with loading row $i$ is exactly the conditional time-series Gram matrix in \Cref{a:gram}; the diagonal block associated with factor row $t$ is the corresponding conditional cross-sectional Gram matrix. Hence these diagonal blocks have eigenvalues between $c$ and $C$. The local part of \Cref{a:identification} bounds the full operator away from zero and therefore bounds both Schur complements away from zero.

The $(i,t)$ loading--factor cross block contains the factor $1/\sqrt{NT}$ from \eqref{eq:local_derivative}. Conditional Cauchy--Schwarz, the moment bounds, and incoherence give its Frobenius norm at most $C/\sqrt{NT}$. Consequently, the squared norm of all cross blocks in a fixed loading row is at most $T\,C/(NT)=C/N$, and the analogous sum for a fixed factor row is at most $C/T$.

Write $A=G_{\lambda\lambda}$, $D=G_{ff}$, and $B=G_{\lambda f}$. The matrices $A$ and $D$ are block diagonal across loading rows and factor rows, respectively, with uniformly bounded inverses. The block inverse identities
$$
\begin{pmatrix}A&B\\B'&D\end{pmatrix}^{-1}_{11}
=(A-BD^{-1}B')^{-1}
=
A^{-1}
+
A^{-1}B(D-B'A^{-1}B)^{-1}B'A^{-1},
$$
and
$$
\begin{pmatrix}A&B\\B'&D\end{pmatrix}^{-1}_{12}
=-(A-BD^{-1}B')^{-1}BD^{-1},
$$
together with their lower-block analogues, give the stated rates. In particular, an off-diagonal loading--loading block of the inverse factors through one loading--factor row of $B$, the bounded inverse factor Schur complement, and one factor--loading column of $B'$. The squared Frobenius norm of a fixed loading row of this correction is therefore bounded by $C/N$ after summing over $j\ne i$. The factor--factor correction is symmetric and has squared row norm bounded by $C/T$. The same factorization gives the loading--factor and factor--loading row bounds. This proves \eqref{eq:inverse_gram_rows}.
\end{proof}

\begin{proposition}[Riesz-score regularity and variance rates]
\label{prop:target_score}
Under \Cref{a:stab,a:exog,a:geometry,a:ned,a:moments,a:signal,a:identification,a:gram,a:target}, let $\delta=\eta/4$. For every parameter of interest defined in \Cref{sec:model},
\begin{equation}
\label{eq:score_lyapunov}
\frac{
\sum_{i,t}
\Ezero(|Z_{it}|^{2+\delta}\mid\mathscr C)
}{
s_{NT}^{2+\delta}
}
\longrightarrow0,
\qquad
\frac{
\max_{i,t}\Ezero(Z_{it}^2\mid\mathscr C)
}{
s_{NT}^2
}
\longrightarrow0.
\end{equation}
Moreover, $Z_{it}$ is conditionally $L_{2+\delta}$-NED on the innovation field, with an approximation coefficient bounded by $C\chi(h,r)$ after normalization by its parameter-specific scale. The conditional variance satisfies:
\begin{enumerate}[label=(\alph*),leftmargin=2em]
\item for a regular individual coefficient or a fixed unit--time lag sum,
$$
s_{NT}^2\asymp N^{-1}+T^{-1};
$$
\item for a mean over $G_N\times R_T$,
$$
s_{NT}^2\asymp (|G_N||R_T|)^{-1};
$$
\item for a contrast between disjoint groups $G_{1,N}$ and $G_{2,N}$ averaged over $R_T$,
$$
s_{NT}^2
\asymp
\frac1{|R_T|}
\left(
\frac1{|G_{1,N}|}
+
\frac1{|G_{2,N}|}
\right).
$$
\end{enumerate}
\end{proposition}

\begin{proof}
Let $\bm\psi_t=(\Psi_{0,1t},\ldots,\Psi_{0,Nt})'$. Conditional sequential exogeneity makes the time-level Riesz sums martingale differences, so cross-time covariances vanish. Because $\bm\psi_t$ is measurable with respect to current covariates, past panel information, and the common shocks,
$$
\begin{aligned}
s_{NT}^2
&=
\sum_{t=1}^{T}
\Ezero\left[
\bm\psi_t'
\bm\Sigma_{u,t}
\bm\psi_t
\middle|\mathscr C
\right].
\end{aligned}
$$
The absolute row-sum bound in \Cref{a:moments} gives
$$
s_{NT}^2
\le
C\Ezero\left[
\sum_{i,t}\Psi_{0,it}^{2}
\middle|\mathscr C
\right],
$$
while the target-specific noncancellation restriction in \Cref{a:target} gives the reverse inequality. Therefore,
$$
s_{NT}^2
\asymp
\Ezero\left[
\sum_{i,t}\Psi_{0,it}^{2}
\middle|\mathscr C
\right]
=
\Gamma_0(\bq_0,\bq_0).
$$
The Gram-matrix bounds, \eqref{eq:q_norm_identity}, and \eqref{eq:q_target_norm_equivalence} therefore imply
$$
s_{NT}^2
\asymp
\|\calP_0\bD\|^2.
$$
For an individual entry, the explicit tangent projection
$$
\calP_{\calT_M}(D)
=
\bU_M\bU_M'D
+
D\bV_M\bV_M'
-
\bU_M\bU_M'D\bV_M\bV_M'
$$
and the leverage bounds give
$$
\|\calP_0\bD\|^2
\asymp
\frac1N+\frac1T.
$$
For a fixed unit--time lag sum, the direction contains only the fixed number $P$ of entry directions. The positive-definiteness requirement in \Cref{a:target} rules out cancellation among this fixed collection, so its variance has the same order $N^{-1}+T^{-1}$.

For a mean over $G_N\times R_T$, orthogonal projection is contractive and \Cref{a:target} gives
$$
\|\calP_0\bD\|^2
\asymp
\|\bD\|^2
=
\|\bm a_{G_N}\|_2^2
\|\bm b_{R_T}\|_2^2
=
\frac1{|G_N||R_T|}.
$$
For a contrast between disjoint groups,
$$
\|\bm a_{G_{1,N}}-\bm a_{G_{2,N}}\|_2^2
=
|G_{1,N}|^{-1}+|G_{2,N}|^{-1},
$$
which gives the stated contrast variance after multiplication by $\|\bm b_{R_T}\|_2^2=|R_T|^{-1}$.

The block-inverse calculation in \eqref{eq:inverse_gram_rows}, which follows from \Cref{a:signal,a:identification,a:gram}, implies that the time-series Riesz-score contributions associated with the target unit have conditional $2+\delta$ norms bounded by $CT^{-1}$ and the cross-sectional contributions associated with the target time are bounded by $CN^{-1}$. Hence
\begin{align*}
\sum_{i,t}
\Ezero(|Z_{it}|^{2+\delta}\mid\mathscr C)
&\le
C\{T\cdot T^{-(2+\delta)}+N\cdot N^{-(2+\delta)}\},\\
s_{NT}^{2+\delta}
&\asymp
(N^{-1}+T^{-1})^{1+\delta/2},
\end{align*}
and their ratio converges to zero. The largest variance contribution is bounded by $C(N^{-2}+T^{-2})$, which is negligible relative to $N^{-1}+T^{-1}$.

For a mean over $G_N\times R_T$, the target direction is $\bm a_{G_N}\bm b_{R_T}'$. The largest unit and time weights are $|G_N|^{-1}$ and $|R_T|^{-1}$. The same block-inverse bounds transfer these maximum weights to the corresponding Riesz-score contributions. Because the maximum target weight converges to zero whenever a support grows, the same Lyapunov calculation gives \eqref{eq:score_lyapunov}. Fixed-time group means and contrasts use $R_T=\{t_0\}$; the cross-sectional averaging weights still make every individual unit contribution negligible. Lag sums contain only the fixed number $P$ of autoregressive coefficient matrices, so the same bounds hold after multiplication by a constant.

Finally, \Cref{lem:localization} applied to $\mathcal Y_{it}(\bq_0)u_{it}$ gives the conditional NED approximation. The norm equivalence in \eqref{eq:q_target_norm_equivalence} provides the target-specific normalization.
\end{proof}

\begin{proof}[Proof of \Cref{thm:population_clt}]
Conditional sequential exogeneity and measurability of $\Psi_{0,it}$ with respect to current covariates, lagged outcomes, and $\mathscr C$ imply
$$
\Ezero(Z_{it}\mid x_{it},\mathscr F_{t-1},\mathscr C)
=
\Psi_{0,it}
\Ezero(u_{it}\mid x_{it},\mathscr F_{t-1},\mathscr C)
=
0.
$$
Hence $\Ezero(Z_{it}\mid\mathscr C)=0$.

Treat $(i,t)$ as a location in $\R^{d_s+1}$ with distance $d_N(i,k)+|t-s|$. Uniform separation in space and integer spacing in time satisfy the sampling requirement of the NED random-field central limit theorem in \citet{jenishprucha2012}. By \Cref{lem:conditional_limit_transfer}, that theorem may be applied under the regular conditional law given $\mathscr C$. By \Cref{lem:localization,prop:target_score}, the normalized score array is conditionally $L_{2+\delta}$-NED, has a summable approximation rate, satisfies the conditional Lyapunov condition, and has no dominating observation. Exponential mixing and \Cref{a:geometry} imply the required mixing summability:
$$
\sum_{a=1}^{\infty}
(1+a)^{d_s}
\bar\alpha(a)^{\delta/(2+\delta)}
<
\infty.
$$
Let $\bm\psi_t=(\Psi_{0,1t},\ldots,\Psi_{0,Nt})'$. Conditional sequential exogeneity removes cross-time covariance, so
$$
s_{NT}^2
=
\sum_t
\Ezero[
\bm\psi_t'\bm\Sigma_{u,t}\bm\psi_t
\mid\mathscr C].
$$
The row-sum bound in \Cref{a:moments} gives the upper comparison
$$
s_{NT}^2
\le
C\Ezero\!\left[
\sum_t\|\bm\psi_t\|_2^2
\middle|\mathscr C
\right],
$$
while the Riesz-direction noncancellation restriction in \Cref{a:target} gives the matching lower comparison. The Riesz equation and the Gram bounds then make this squared-weight sum comparable to $\|\calP_0\bD\|^2$. The conditional variance is therefore nondegenerate at the target-specific rate established in \Cref{prop:target_score}. The cited theorem yields \eqref{eq:conditional_clt}. Applying the same argument to every conditionally variance-normalized linear combination of a fixed target vector and then using Cram\'er--Wold proves the joint statement.
\end{proof}

\subsection{Full-panel concentration and fixed-rank recovery}
\label{app:recovery}

This subsection proves \Cref{thm:recovery}. The argument first controls the residual score and the empirical prediction criterion uniformly over the admissible low-rank differences. The least-squares basic inequality then gives recovery of all coefficient matrices. Strong singular values convert matrix recovery into singular-space and tangent-space recovery, after which local loading and factor equations yield max-norm and Riesz-representer consistency.

Define the empirical Gram form
$$
\widehat\Gamma(\bDelta,\bQ)
=
\sum_{i,t}
\mathcal Y_{it}(\bDelta)
\mathcal Y_{it}(\bQ).
$$

Let
$$
\calD_{\max}
=
\left\{
\bTheta-\bTheta_0:
\ \bTheta\in\bigcup_{\br\in\calR_{\max}}\calM(\br),
\ \|\bTheta\|_{\max}\le B
\right\}
$$
denote the set of admissible coefficient differences. For every $\bDelta\in\calD_{\max}$ and every coefficient matrix $M$, $\rank(\bDelta_M)\le \bar r_M+r_{M,0}$ and $\|\bDelta_M\|_{\max}\le2B$.


\begin{lemma}[Weighted concentration on the unit--time index]
\label{lem:weighted_concentration}
For this lemma only, write $n=NT$ and $d=d_s+1$. Let $X_a$, $a=(i,t)$, be a scalar array satisfying, conditionally on $\mathscr C$,
\[
\Ezero(X_a\mid\mathscr C)=0,
\qquad
\sup_a\|X_a\|_{2+\delta,\mathscr C}\le C
\]
for some fixed $\delta>0$. Suppose that for every integer $q\ge1$ there is a bounded local approximation $X_a^{(q)}$, measurable with respect to the unit-specific innovations in the ball of radius $q/3$ around $a$ together with $\mathscr C$, such that
\[
|X_a^{(q)}|\le M_n,
\qquad
\sup_a\|X_a-X_a^{(q)}\|_{2+\delta,\mathscr C}
\le C M_n e^{-cq},
\]
where $M_n\le n^C$. Let $\mathcal W_n$ be a finite deterministic collection of weights $w=(w_a)$ satisfying
\[
\sum_a|w_a|\le1,
\qquad
\sum_aw_a^2\le v_n^2,
\qquad
\max_a|w_a|\le\omega_n.
\]
Then, for every fixed $A<\infty$, one may choose a constant $K_A$ such that, with $q_n=\lceil K_A\log n\rceil$, for every $x_n\in(0,1]$ that is bounded below by a fixed negative power of $n$,
\begin{equation}
\label{eq:weighted_bernstein}
\Pzero\!\left(
\max_{w\in\mathcal W_n}
\left|\sum_aw_aX_a\right|>x_n
\middle|\mathscr C
\right)
\le
2|\mathcal W_n|
\exp\!\left[
-\frac{c x_n^2}
{v_n^2+M_n\omega_n\{\log n\}^{d}x_n}
\right]
+Cn^{-A}
\end{equation}
almost surely for all sufficiently large $n$.

The same blocking argument has the following matrix consequence. If $X=(X_{it})$ is an $N\times T$ centered array satisfying the same bounded local-approximation conditions, then
\begin{equation}
\label{eq:matrix_block_concentration}
\|X\|_{\op}
=
O_p\!\left(
\sqrt{(N+T)\log n}
+
M_n\{\log n\}^{(d_s+3)/2}
\right).
\end{equation}
The conclusion continues to hold for arrays of fixed-size matrix blocks, with only the constants changed.
\end{lemma}

\begin{proof}
The combined unit--time index has polynomial ball growth of dimension $d=d_s+1$: by \Cref{a:geometry},
\[
\sup_{i,t}
\left|
\{(j,s):d_N(i,j)+|t-s|\le r\}
\right|
\le
C(1+r)^{d_s+1}.
\]
Fix $q=q_n$. Using the spatial locations from \Cref{a:geometry} together with the integer time coordinate, partition the product space into blocks of side length proportional to $q$. A fixed number, depending only on $d$, of residue classes is enough to color the blocks so that two distinct blocks with the same color are separated by at least $q$. Each block contains at most $Cq^d$ observed unit--time points.

First replace $X_a$ by $X_a^{(q)}$. Since the weights have $\ell_1$ norm at most one, the replacement is uniform over $\mathcal W_n$. Conditional Markov's inequality and a union bound over the $n$ observations give, for any fixed polynomially small $x_n$,
\[
\Pzero\!\left(
\max_a|X_a-X_a^{(q)}|>x_n/4
\middle|\mathscr C
\right)
\le
Cn\left(\frac{M_ne^{-cq}}{x_n}\right)^{2+\delta}.
\]
Choosing $K_A$ sufficiently large makes this probability $O(n^{-A})$.

For one color, collect the localized variables in each block into a single finite-dimensional random vector. The corresponding block sigma-fields are separated by at least a constant multiple of $q$. The definition of $\alpha_{\mathscr C}$ in \Cref{a:ned} already takes the supremum over arbitrary separated index sets, so there is no block-cardinality multiplier in the mixing coefficient. Iterating the coupling theorem of \citet{bradley1983} therefore constructs independent copies of the localized block vectors, preserving each block distribution, on an extension of the conditional probability space. Exponential mixing and $q=K_A\log n$ make the probability that any coupled block differs by more than the amount used in the Bernstein comparison $O(n^{-A})$. The number of colors is fixed, so the same bound holds for all colors simultaneously.

On the coupling event, for a fixed $w\in\mathcal W_n$ the sum within a block is bounded by
\[
CM_nq^d\omega_n.
\]
Moreover, the conditional covariance inequality in \eqref{eq:conditional_covariance_bound}, exponential NED approximation, and polynomial shell growth imply an absolute covariance-row-sum bound for the original centered field. The localized covariance matrix differs from it by $O(ne^{-cq})$ in row sum. Hence, uniformly over the blocks,
\[
\sum_{\text{blocks}}
\operatorname{Var}_0\!\left(
\sum_{a\ \text{in block}}w_aX_a^{(q)}
\middle|\mathscr C
\right)
\le
C\sum_aw_a^2
\le Cv_n^2.
\]
Scalar Bernstein's inequality applied to the independent coupled block sums, followed by the fixed-color union bound and then the union bound over $\mathcal W_n$, gives \eqref{eq:weighted_bernstein} after constants are enlarged.

For the matrix statement, couple the entire localized block matrices rather than scalar block sums. A block matrix has at most $Cq^d$ nonzero entries, so its operator norm is at most $CM_nq^{d/2}$. The same covariance-row-sum calculation gives the rectangular matrix variance bounds
\[
\left\|
\sum_{\text{blocks}}
\Ezero(BB'\mid\mathscr C)
\right\|_{\op}
\le CT,
\qquad
\left\|
\sum_{\text{blocks}}
\Ezero(B'B\mid\mathscr C)
\right\|_{\op}
\le CN.
\]
The rectangular matrix Bernstein inequality of \citet{tropp2012}, applied separately to the fixed number of color classes, therefore gives
\[
\|X^{(q)}\|_{\op}
=
O_p\!\left[
\sqrt{(N+T)\log n}
+M_nq^{d/2}\log n
\right].
\]
The localization and coupling errors are $o_p(1)$ after increasing $K_A$, and $q\asymp\log n$, which proves \eqref{eq:matrix_block_concentration}. Fixed block size only multiplies the displayed bounds by a constant.
\end{proof}

\begin{proposition}[Uniform full-panel concentration]
\label{prop:uniform_concentration}
Under \Cref{a:stab,a:exog,a:geometry,a:ned,a:moments,a:signal,a:identification,a:gram,a:growth},
\begin{equation}
\label{eq:uniform_score_bound}
\sup_{\bDelta\in\calD_{\max}\setminus\{0\}}
\frac{
|\sum_{i,t}u_{it}\mathcal Y_{it}(\bDelta)|
}{
\|\bDelta\|
}
=
O_p\{\sqrt{NT\zeta_{NT}}\},
\end{equation}
and, uniformly over $\calD_{\max}$,
\begin{equation}
\label{eq:empirical_lower_bound}
\sum_{i,t}\mathcal Y_{it}(\bDelta)^2
\ge
\frac c2\|\bDelta\|^2-CNT\zeta_{NT}.
\end{equation}
On $\calT_0$, the empirical least-squares Gram form converges uniformly to its conditional population counterpart.
\end{proposition}

\begin{proof}
Let $z_{it,M}$ denote the regressor multiplying coefficient matrix $M$. By \Cref{lem:localization}, the arrays $u_{it}z_{it,M}$ and $z_{it,M}z_{it,M'}$ have the conditional moments and exponential local-approximation rates required below. Put $b=b_{NT}$ and let $\mathcal E_b$ be the event on which every $|u_{it}|$ and every $|z_{it,M}|$ is at most $b$. Stability and the conditional $8+\eta$ moment bounds imply
\[
\Pzero(\mathcal E_b^c\mid\mathscr C)
\le
CNT b^{-(8+\eta)}
\le
C\{\log(NT)\}^{-(8+\eta)}
\longrightarrow0.
\]
Thus all truncations used below agree simultaneously with the original panel on an event whose conditional probability approaches one.

\emph{Step 1: the score term.}
For each $M$, clip $u_{it}$ and $z_{it,M}$ at $b$, multiply them, and subtract the conditional mean given $\mathscr C$. Let $\bm S_M^{b}$ denote the resulting centered $N\times T$ matrix. Its entries are bounded by $Cb^2$, have uniformly bounded conditional $2+\delta$ moments for some fixed $\delta>0$, and have exponential local approximations by \Cref{lem:localization}. The matrix conclusion of \Cref{lem:weighted_concentration} gives
\begin{equation}
\label{eq:score_matrix_block_bound}
\max_M\|\bm S_M^{b}\|_{\op}
=
O_p\!\left[
\sqrt{(N+T)\log(NT)}
+b^2\{\log(NT)\}^{(d_s+3)/2}
\right].
\end{equation}
Conditional sequential exogeneity gives $\Ezero(u_{it}z_{it,M}\mid\mathscr C)=0$. Since clipping changes a score entry only on $\mathcal E_b^c$, conditional H\"older and the $8+\eta$ moments imply that the matrix of clipping means has Frobenius norm $o(1)$. Therefore \eqref{eq:score_matrix_block_bound} also holds for the original score matrix $\bm S_M$. Under \Cref{a:growth},
\[
\sqrt{(N+T)\log(NT)}
+b^2\{\log(NT)\}^{(d_s+3)/2}
=
O\!\left(\sqrt{NT\zeta_{NT}}\right).
\]
Duality between the operator norm and the sum of singular values, together with the fixed rank caps, gives, uniformly over $\calD_{\max}$,
\[
\left|
\sum_{i,t}u_{it}\mathcal Y_{it}(\bDelta)
\right|
\le
C\max_M\|\bm S_M\|_{\op}\,\|\bDelta\|,
\]
which proves \eqref{eq:uniform_score_bound}.

\emph{Step 2: uniform prediction curvature on the rank-capped set.}
Write $n=NT$, $m=N+T$, and $d=d_s+1$ within this step. Expanding the centered empirical prediction norm gives
\begin{equation}
\label{eq:quadratic_pair_expansion}
\sum_{i,t}
\left[
\mathcal Y_{it}(\bDelta)^2
-
\Ezero\{\mathcal Y_{it}(\bDelta)^2\mid\mathscr C\}
\right]
=
\sum_{M,M'}\sum_{i,t}
X_{it}^{MM'}\Delta_{M,it}\Delta_{M',it},
\end{equation}
where
\[
X_{it}^{MM'}
=
z_{it,M}z_{it,M'}
-
\Ezero(z_{it,M}z_{it,M'}\mid\mathscr C).
\]
On $\mathcal E_b$, replace $X_{it}^{MM'}$ by its clipped version; this is bounded by $Cb^2$ and satisfies the hypotheses of \Cref{lem:weighted_concentration}.

Let
\[
r_0^2=C_0n\zeta_{NT}
=C_0b^2m\{\log n\}^{d+1},
\]
where $C_0$ is a sufficiently large fixed constant. If $\|\bDelta\|<r_0$, the lower bound \eqref{eq:empirical_lower_bound} is immediate after the constant multiplying $n\zeta_{NT}$ is enlarged, because the empirical prediction norm is nonnegative. It remains to consider $\|\bDelta\|\ge r_0$.

Peel this set into dyadic norm ranges $r\le\|\bDelta\|<2r$. There are only $O(\log n)$ nonempty ranges because the coefficient bound implies $\|\bDelta\|\le C\sqrt n$. In each range choose an $\epsilon_n$-net, in Frobenius norm, of the normalized rank-capped collections $\bDelta/\|\bDelta\|$, with
\[
\epsilon_n=c_0/b^2.
\]
The ranks and number of coefficient matrices are fixed, so the usual singular-vector parameterization yields
\[
\log|\mathcal N_r|
\le
Cm\log(Cb)
\le Cm\log n.
\]
The net centers may be chosen from the normalized admissible set itself. Hence, for a center $\bm\delta$ from a range with lower endpoint $r$,
\[
\max_{i,t,M}|\delta_{M,it}|
\le\frac{2B}{r}.
\]
For every fixed pair $(M,M')$, the weights
\[
w_{it}^{MM'}=\delta_{M,it}\delta_{M',it}
\]
satisfy
\[
\sum_{i,t}|w_{it}^{MM'}|\le1,
\qquad
\sum_{i,t}(w_{it}^{MM'})^2
\le\omega_r,
\qquad
\omega_r\le C/r^2.
\]
Apply \eqref{eq:weighted_bernstein} simultaneously to the finite collection of all net points and all matrix pairs, with a fixed small threshold $x>0$. Since $M_n=Cb^2$ and $q_n^d\asymp\{\log n\}^d$, the exponent is bounded below by
\[
\frac{cr^2}{b^2\{\log n\}^{d}}
\ge
cC_0m\log n.
\]
Choosing $C_0$ large enough makes this dominate the net entropy, the finite number of matrix pairs, and the $O(\log n)$ norm ranges. Thus, with conditional probability approaching one, the absolute value of \eqref{eq:quadratic_pair_expansion} is at most $c\|\bDelta\|^2/4$ at every net point.

On $\mathcal E_b$, multiplication by the regressors maps the coefficient collection into fitted values with operator norm at most $Cb$. Consequently, both the empirical and conditional-population quadratic forms are $Cb^2$-Lipschitz on the unit Frobenius sphere. The choice $\epsilon_n=c_0/b^2$ therefore extends the preceding bound from each net to its entire dyadic range after $c_0$ is chosen sufficiently small. Conditional prediction identification then gives, uniformly for $\|\bDelta\|\ge r_0$,
\[
\sum_{i,t}\mathcal Y_{it}(\bDelta)^2
\ge
\frac c2\|\bDelta\|^2.
\]
Combining the large- and small-norm ranges proves \eqref{eq:empirical_lower_bound}.

\emph{Step 3: the empirical Gram operator on the true tangent space.}
Use the normalized loading and factor coordinates introduced before \Cref{lem:inverse_gram_rows}. The loading--loading diagonal block for unit $i$ is a $T$-average of fixed-dimensional products, and the factor--factor diagonal block for time $t$ is an $N$-average. Apply \eqref{eq:weighted_bernstein} entry by entry with weights $1/T$ or $1/N$ and take a union bound over the $N+T$ blocks. Under \Cref{a:growth},
\[
\max_i\|\widehat G_{\lambda\lambda,i}-G_{\lambda\lambda,i}\|_{\op}
+
\max_t\|\widehat G_{ff,t}-G_{ff,t}\|_{\op}
=o_p(1).
\]
The loading--factor part is a rectangular block matrix whose $(i,t)$ block is scaled by $(NT)^{-1/2}$. Apply the matrix conclusion \eqref{eq:matrix_block_concentration} to the centered, clipped fixed-size block array and then divide by $\sqrt{NT}$ to obtain
\[
\|\widehat G_{\lambda f}-G_{\lambda f}\|_{\op}
=
O_p\!\left[
\sqrt{\frac{(N+T)\log(NT)}{NT}}
+
\frac{b^2\{\log(NT)\}^{(d_s+3)/2}}{\sqrt{NT}}
\right]
=o_p(1).
\]
The transpose block has the same bound. Hence
\[
\|\widehat{\mathbb G}-\mathbb G_0\|_{\op}=o_p(1)
\]
on the normalized tangent coordinates, which is equivalent to the stated uniform convergence of the empirical least-squares Gram form on $\calT_0$.
\end{proof}

\begin{proof}[Proof of \Cref{thm:recovery}]
Let
$$
\widehat\bDelta
=
\widehat\bTheta(\br_0)-\bTheta_0.
$$
Optimality of $\widehat\bTheta(\br_0)$ and feasibility of $\bTheta_0$ imply
\begin{align*}
0
&\ge
\frac1{2NT}
\sum_{i,t}
\left[
\{u_{it}-\mathcal Y_{it}(\widehat\bDelta)\}^2-u_{it}^2
\right]\\
&=
\frac1{2NT}
\widehat\Gamma(\widehat\bDelta,\widehat\bDelta)
-
\frac1{NT}
\sum_{i,t}
u_{it}\mathcal Y_{it}(\widehat\bDelta).
\end{align*}
Multiplying by $NT$ and applying \Cref{prop:uniform_concentration,a:growth} gives
$$
\frac c4\|\widehat\bDelta\|^2
\le
C\sqrt{NT\zeta_{NT}}\|\widehat\bDelta\|
+
CNT\zeta_{NT}.
$$
If $x=\|\widehat\bDelta\|$, the inequality $c x^2/4\le a x+b$ implies
$$
x
\le
\frac{4a}{c}
+
2\sqrt{\frac{b}{c}}.
$$
Therefore,
\begin{equation}
\label{eq:global_frobenius_rate}
\|\widehat\bDelta\|
=
O_p\{\sqrt{NT\zeta_{NT}}\}.
\end{equation}
Substituting \eqref{eq:global_frobenius_rate} into the basic inequality yields
$$
\sum_{i,t}
\mathcal Y_{it}(\widehat\bDelta)^2
=
O_p(NT\zeta_{NT}).
$$

For a positive-rank coefficient matrix $M$,
$$
\|\widehat\bM-\bM_0\|_{\op}
\le
\|\widehat\bM-\bM_0\|_F
=
O_p\{\sqrt{NT\zeta_{NT}}\}.
$$
Because $\sigma_{r_M}(\bM_0)\ge c\sqrt{NT}$, the Wedin $\sin\Theta$ perturbation bound \citep{wedin1972} gives
$$
\|\widehat\bU_M\widehat\bU_M'-\bU_M\bU_M'\|_{\op}
+
\|\widehat\bV_M\widehat\bV_M'-\bV_M\bV_M'\|_{\op}
=
O_p(\zeta_{NT}^{1/2}).
$$
The formula for the tangent projection then implies
\begin{equation}
\label{eq:tangent_projection_rate}
\|\widehat\calP-\calP_0\|_{\op}
=
O_p(\zeta_{NT}^{1/2}).
\end{equation}
A supplied rank-zero matrix is exactly zero.

We next derive the pointwise rates using the loading--factor normalization, the matrix $\mathbb Z_{it}$, and the local derivative introduced before \Cref{lem:inverse_gram_rows}. Align the estimated factors and loadings with that normalization and set
$$
\bm h_i^\lambda
=
\sqrt T(\widehat{\bm\lambda}_i-\bm\lambda_{0,i}),
\qquad
\bm h_t^f
=
\sqrt N(\widehat{\bm f}_t-\bm f_{0,t}).
$$
the fitted-value change is exactly
\begin{equation}
\label{eq:local_fitted_change}
\mathcal Y_{it}(\widehat\bTheta)-\mathcal Y_{it}(\bTheta_0)
=
\dot{\mathcal Y}_{it}(\bm h)
+
R_{it}(\bm h),
\end{equation}
where $\dot{\mathcal Y}_{it}(\bm h)$ is defined in \eqref{eq:local_derivative} and
\begin{equation}
\label{eq:local_product}
R_{it}(\bm h)
=
\frac{(\bm h_i^\lambda)'\mathbb Z_{it}\bm h_t^f}{\sqrt{NT}}.
\end{equation}
For a direction $\bm q$, the derivative of \eqref{eq:local_fitted_change} at $\bm h$ is
$$
\dot{\mathcal Y}_{it}(\bm q)
+
C_{it}(\bm h,\bm q),
$$
where
\begin{equation}
\label{eq:local_cross_derivative}
C_{it}(\bm h,\bm q)
=
\frac{
(\bm h_i^\lambda)'\mathbb Z_{it}\bm q_t^f
+
(\bm q_i^\lambda)'\mathbb Z_{it}\bm h_t^f
}{
\sqrt{NT}
}.
\end{equation}

Define the normalized score blocks
\begin{align*}
\bm s_i^\lambda
&=
\frac1{\sqrt T}
\sum_{t=1}^{T}
(\mathbb Z_{it}\bm f_{0,t})u_{it},\\
\bm s_t^f
&=
\frac1{\sqrt N}
\sum_{i=1}^{N}
(\mathbb Z_{it}'\bm\lambda_{0,i})u_{it}.
\end{align*}
The empirical first-order conditions and \eqref{eq:local_fitted_change}--\eqref{eq:local_cross_derivative} imply
\begin{equation}
\label{eq:local_normal_equation}
\widehat{\mathbb G}(\widehat{\bm h},\bm q)
=
\bm s(\bm q)
+
\mathcal N(\widehat{\bm h},\bm q)
\qquad
\text{for every }\bm q,
\end{equation}
where
\begin{align}
\mathcal N(\bm h,\bm q)
={}&
\sum_{i,t}
u_{it}C_{it}(\bm h,\bm q)
-
\sum_{i,t}
\dot{\mathcal Y}_{it}(\bm h)C_{it}(\bm h,\bm q)
\nonumber\\
&-
\sum_{i,t}
R_{it}(\bm h)\dot{\mathcal Y}_{it}(\bm q)
-
\sum_{i,t}
R_{it}(\bm h)C_{it}(\bm h,\bm q).
\label{eq:local_nonlinear_term}
\end{align}
The tangent-space Gram part of \Cref{prop:uniform_concentration} gives
$$
\|\widehat{\mathbb G}-\mathbb G_0\|_{\op}=o_p(1).
$$
By \Cref{lem:inverse_gram_rows}, the inverse local Gram operator has bounded diagonal rows and square-summable loading--factor cross rows at the rates in \eqref{eq:inverse_gram_rows}. The bounds in \eqref{eq:inverse_gram_rows}, the conditional moment bounds for $\bm s_i^\lambda$ and $\bm s_t^f$, and \eqref{eq:panel_growth} imply
$$
\|\mathcal N(\widehat{\bm h},\cdot)\|
=
o_p(1)
$$
for every fixed target loading or factor row. Inverting \eqref{eq:local_normal_equation} therefore gives
\begin{equation}
\label{eq:pointwise_factor_expansion}
\bm h_i^\lambda
=
[\mathbb G_0^{-1}\bm s]_i^\lambda+o_p(1),
\qquad
\bm h_t^f
=
[\mathbb G_0^{-1}\bm s]_t^f+o_p(1).
\end{equation}
This proves root-$T$ loading and root-$N$ factor recovery for fixed target rows. For the uniform version, truncate every coordinate of $\bm s_i^\lambda$ and $\bm s_t^f$ at the same level $b_{NT}^2$ used in \Cref{prop:uniform_concentration}. Applying \Cref{lem:weighted_concentration} to the localized row scores with the equal time or unit weights, followed by a union bound over the fixed-dimensional coordinates of all $N+T$ rows, gives
$$
\max_i\|\bm s_i^\lambda\|
+
\max_t\|\bm s_t^f\|
=
O_p\{b_{NT}\sqrt{\log(NT)}\}.
$$
The same-side and cross-side inverse-row bounds in \eqref{eq:inverse_gram_rows} imply the identical order for the rows of $\mathbb G_0^{-1}\bm s$. On the ball with this radius, \eqref{eq:local_product}--\eqref{eq:local_nonlinear_term}, conditional moments, and fixed ranks give
$$
\max_i\|\mathcal N(\widehat{\bm h},\cdot)_{\lambda_i}\|
+
\max_t\|\mathcal N(\widehat{\bm h},\cdot)_{f_t}\|
\le
C b_{NT}^2\log(NT)
\left(
N^{-1/2}+T^{-1/2}
\right)
+
o_p(1).
$$
The right side is $o_p(1)$ by \Cref{a:growth}. Applying the row bounds of $\mathbb G_0^{-1}$ once more to \eqref{eq:local_normal_equation} therefore yields
$$
\max_i\|\bm h_i^\lambda\|
=O_p\{b_{NT}\sqrt{\log(NT)}\},
\qquad
\max_t\|\bm h_t^f\|
=O_p\{b_{NT}\sqrt{\log(NT)}\}.
$$
Because the true loading and factor rows are bounded by incoherence and the signal normalization, expanding each fitted entry yields
\begin{equation}
\label{eq:max_norm_recovery}
\|\widehat\bTheta-\bTheta_0\|_{\max}
=
O_p\!\left[
b_{NT}\sqrt{\frac{\log(NT)}{T}}
+
b_{NT}\sqrt{\frac{\log(NT)}{N}}
+
\frac{b_{NT}^2\log(NT)}{\sqrt{NT}}
\right]
=O_p(\zeta_{NT}^{1/2})
=o_p(1).
\end{equation}
This proves the max-norm statement in \Cref{thm:recovery} and, together with the strict margin in \Cref{a:signal}, makes the coefficient bound inactive with probability approaching one.

Let $\bm h^*$ denote the local coordinate of the population Riesz representer. It is the unique solution of
\begin{equation}
\label{eq:local_target_weight}
\mathbb G_0(\bm q,\bm h^*)
=
\dot\phi(\bm q)
\qquad
\text{for every }\bm q,
\end{equation}
where $\dot\phi$ is the derivative of the parameter. The empirical Riesz coordinate satisfies the corresponding equation with $\widehat{\mathbb G}$ and the estimated tangent space. For a mean or contrast, \Cref{a:target,a:gram} and \eqref{eq:tangent_projection_rate} give
$$
\|\widehat\bq-\bq_0\|
\le
C
\left[
\|\widehat{\mathbb G}-\mathbb G_0\|_{\op}
+
\|\widehat\calP-\calP_0\|_{\op}
\right]
\|\calP_0\bD\|
=
o_p(\|\calP_0\bD\|).
$$
For an individual entry, the global tangent bound must be evaluated at the target row. The explicit tangent projection and \eqref{eq:pointwise_factor_expansion} give
$$
\|(\widehat\calP-\calP_0)\bD\|
=
o_p\left\{
\left(
\frac1N+\frac1T
\right)^{1/2}
\right\}
=
o_p(\|\calP_0\bD\|).
$$
The same inverse-Gram argument then yields the Riesz rate. Finally, conditional identification gives
$$
\sum_{i,t}
(\widehat\Psi_{it}-\Psi_{0,it})^2
\le
C\|\widehat\bq-\bq_0\|^2
=
o_p(\|\calP_0\bD\|^2).
$$
\paragraph{Weighted residual identity.}
Let
$$
L_{NT}(\bTheta)
=
\frac{1}{2NT}\sum_{i,t}\{y_{it}-\mathcal Y_{it}(\bTheta)\}^2.
$$
On the event that the coefficient bound is inactive, every $\bDelta\in\widehat\calT$ is the derivative at zero of a differentiable path $a\mapsto\bTheta(a)$ in the selected fixed-rank set with $\bTheta(0)=\widehat\bTheta$. First-order stationarity gives
$$
0
=
\left.\frac{d}{da}L_{NT}\{\bTheta(a)\}\right|_{a=0}
=
-\frac{1}{NT}\sum_{i,t}\widehat u_{it}\mathcal Y_{it}(\bDelta).
$$
Since $\widehat\bq\in\widehat\calT$, choosing $\bDelta=\widehat\bq$ proves \eqref{eq:weighted_residual_zero}. The max-norm conclusion of \Cref{thm:recovery} and the strict interior margin in \Cref{a:signal} set the coefficient-bound multipliers equal to zero with probability approaching one. Numerical objective, first-order, and Riesz equation residuals enter additively and are negligible under the stated tolerances.
\end{proof}

\subsection{Auxiliary local-deletion comparison}
\label{app:weights}

The empirical Riesz weights and $u_{it}$ are calculated from the same panel, so their product cannot be centered by treating the weights as fixed. This subsection introduces an auxiliary estimator that removes a small space--time neighborhood around $(i,t)$. The auxiliary estimator is used only in the proof. Conditional NED and alpha mixing separate it from the current score, while stability shows that deleting the neighborhood changes the full-panel estimator by a negligible amount.

Choose proof-only radii
$$
q_T=\lceil c_q\log(NT)\rceil,
\qquad
r_N=\lceil c_r\log(NT)\rceil,
$$
with fixed constants $c_q,c_r$ sufficiently large. Exponential dynamic, NED, and mixing decay, together with \Cref{a:geometry,a:growth}, gives
\begin{equation}
\label{eq:local_deletion_tail}
\begin{aligned}
&b_{NT}^2
\left\{
\frac{2q_T+1}{T}
+
\frac{v_N(r_N)}{N}
\right\}
\longrightarrow0,
\\
&\sqrt{NT}
\left[
\rho^{q_T-P}
+
\psi_T(\lfloor q_T/3\rfloor)
+
\psi_S(\lfloor r_N/3\rfloor)
+
\sum_{a>\lfloor\min(q_T,r_N)/3\rfloor}
(1+a)^{d_s}\bar\alpha(a)^{\eta/(8+\eta)}
\right]
\longrightarrow0.
\end{aligned}
\end{equation}
These radii are used only in the proof and are not estimator tuning parameters.

For observation $a=(i,t)$, define the proof-only neighborhood
$$
\mathcal B_a
=
\{(k,s):|s-t|\le q_T,\ d_N(i,k)\le r_N\}.
$$
Let $\widehat\bTheta^{(-a)}$ be the supplied-rank estimator after deleting $\mathcal B_a$, and let $\widehat\bq^{(-a)}$ be its Riesz representer. These quantities are not computed by the estimator.

\begin{proposition}[Local-deletion stability and comparison]
\label{prop:leave_neighborhood}
Under the assumptions of \Cref{thm:recovery} and \Cref{prop:uniform_concentration},
\begin{equation}
\label{eq:h_deletion_average}
\frac1{NT}\sum_a
\|\widehat{\bm h}-\widehat{\bm h}^{(-a)}\|^2
=o_p(1),
\end{equation}
and
\begin{equation}
\label{eq:q_deletion_average}
\frac1{NT}
\sum_a
\frac{
\|\widehat\bq-\widehat\bq^{(-a)}\|^2
}{
\|\calP_0\bD\|^2
}
=o_p(1).
\end{equation}
\end{proposition}

\begin{proof}
Work in the aligned local coordinates used in the recovery proof. Let $\widehat{\mathbb G}^{(-a)}$ and $\widehat{\bm s}^{(-a)}$ denote the Gram matrix and score after deleting $\mathcal B_a$. Deleting the neighborhood removes at most $2q_T+1$ observations from a loading-row time series and at most $v_N(r_N)$ observations from a factor-row cross-section. Therefore,
\begin{equation}
\label{eq:deleted_information_bound}
\|\widehat{\mathbb G}-\widehat{\mathbb G}^{(-a)}\|_{\op}
\le
C b_{NT}^2
\left\{
\frac{2q_T+1}{T}
+
\frac{v_N(r_N)}{N}
\right\}
=
o_p(1).
\end{equation}
The conditional second moment of the deleted score satisfies
\begin{equation}
\label{eq:deleted_score_bound}
\frac1{NT}
\sum_a
\|
\widehat{\bm s}-\widehat{\bm s}^{(-a)}
\|^2
\le
C(2q_T+1)v_N(r_N)
\left(
\frac1N+\frac1T
\right)
+
o_p(1).
\end{equation}
Subtracting the full and deleted normal equations and using uniform nonsingularity gives
\begin{align}
\|\widehat{\bm h}-\widehat{\bm h}^{(-a)}\|
\le C\Big[
&
\|\widehat{\bm s}-\widehat{\bm s}^{(-a)}\|
+
\|\widehat{\mathbb G}-\widehat{\mathbb G}^{(-a)}\|_{\op}
\|\widehat{\bm h}\|
\nonumber\\
&+
\|\mathcal N(\widehat{\bm h},\cdot)
-\mathcal N^{(-a)}(\widehat{\bm h}^{(-a)},\cdot)\|
\Big].
\label{eq:deleted_estimator_resolvent}
\end{align}
The nonlinear difference is of smaller order by recovery and \eqref{eq:panel_growth}. Averaging the square of \eqref{eq:deleted_estimator_resolvent} and applying \eqref{eq:deleted_information_bound}--\eqref{eq:deleted_score_bound} gives \eqref{eq:h_deletion_average}.

The Riesz equations satisfy, in aligned coordinates,
\begin{align*}
\widehat{\bm h}^*
-
\widehat{\bm h}^{*(-a)}
=
(\widehat{\mathbb G}^{(-a)})^{-1}
\Big[
&
(\widehat{\dot\phi}-\widehat{\dot\phi}^{(-a)})
\\
&-
(\widehat{\mathbb G}-\widehat{\mathbb G}^{(-a)})
\widehat{\bm h}^*
\Big].
\end{align*}
The first difference is controlled by the tangent-space stability just obtained, and the second by \eqref{eq:deleted_information_bound}. Averaging the squared norm proves \eqref{eq:q_deletion_average}.

The same deletion construction will also be used below as a decorrelation device. After replacing an anchored regressor or score contribution by its local innovation approximation, that contribution is measurable inside $\mathcal B_a$, whereas the corresponding deleted coordinate is measurable outside $\mathcal B_a$. The conditional NED covariance inequality therefore bounds their conditional covariance by the dependence-tail quantity in \eqref{eq:local_deletion_tail} times their conditional moments. For centered sums, splitting pairs according to the distance between their local innovation neighborhoods, then using exponential mixing and polynomial neighborhood growth, makes the covariance bound summable. This is the specific local-deletion property invoked in \Cref{prop:bilinear_decorrelation}; no feasible-score decomposition is needed here.
\end{proof}

\subsection{Full-panel target expansion and split-panel correction}
\label{app:bias}

This subsection proves \Cref{thm:target_expansion,thm:twoway}. The local coordinates and score equation were constructed in \Cref{app:recovery}. We now separate their first-order solution from the remaining loading--factor products. Substitution into the parameter of interest produces the Riesz score, a loading-side bias of order $T^{-1}$, a factor-side bias of order $N^{-1}$, and smaller terms. The final part shows algebraically why the two split-panel comparisons remove the two leading biases.

Let $\widehat{\bm h}$ be the local coordinate of the full-panel estimate defined in \Cref{app:recovery}, and define the first-order coordinate $\bm h^{(1)}$ by
\begin{equation}
\label{eq:first_order_coordinate}
\mathbb G_0(\bm h^{(1)},\bm q)
=
\bm s(\bm q)
\qquad
\text{for every }\bm q.
\end{equation}
The target has the exact local expansion
\begin{equation}
\label{eq:local_target_expansion}
\phi\{\bTheta(\bm h)\}
-
\phi(\bTheta_0)
=
\dot\phi(\bm h)
+
\mathcal Q(\bm h),
\end{equation}
where
\begin{equation}
\label{eq:target_quadratic}
\mathcal Q(\bm h)
=
\sum_{M:r_M>0}
\frac{
\ip{\bD_M}{\bm H_M^\lambda(\bm H_M^f)'}
}{
\sqrt{NT}
}.
\end{equation}
Here $\bm H_M^\lambda$ and $\bm H_M^f$ collect the loading and factor coordinates for matrix $M$. By \eqref{eq:local_target_weight} and \eqref{eq:first_order_coordinate},
\begin{equation}
\label{eq:first_order_influence_identity}
\dot\phi(\bm h^{(1)})
=
\mathbb G_0(\bm h^{(1)},\bm h^*)
=
\bm s(\bm h^*)
=
\sum_{i,t}\Psi_{0,it}u_{it}.
\end{equation}

\begin{lemma}[Exact local target decomposition]
\label{lem:exact_target_decomposition}
For an exact interior minimizer,
\begin{align}
\widehat\phi-\phi(\bTheta_0)
={}&
\sum_{i,t}\Psi_{0,it}u_{it}
+
\{\mathbb G_0-\widehat{\mathbb G}\}
(\widehat{\bm h},\bm h^*)
\nonumber\\
&+
\mathcal N(\widehat{\bm h},\bm h^*)
+
\mathcal Q(\widehat{\bm h}).
\label{eq:exact_local_target_decomposition}
\end{align}
For a computed approximate stationary point, an additional numerical remainder is bounded by the reported first-order and Riesz-equation tolerances.
\end{lemma}

\begin{proof}
Subtract the first-order equation \eqref{eq:first_order_coordinate} from the exact local normal equation \eqref{eq:local_normal_equation}. For every direction $\bm q$,
$$
\mathbb G_0(\widehat{\bm h}-\bm h^{(1)},\bm q)
=
\{\mathbb G_0-\widehat{\mathbb G}\}
(\widehat{\bm h},\bm q)
+
\mathcal N(\widehat{\bm h},\bm q).
$$
Set $\bm q=\bm h^*$ and use the local Riesz equation \eqref{eq:local_target_weight} to obtain
$$
\dot\phi(\widehat{\bm h}-\bm h^{(1)})
=
\{\mathbb G_0-\widehat{\mathbb G}\}
(\widehat{\bm h},\bm h^*)
+
\mathcal N(\widehat{\bm h},\bm h^*).
$$
Combining this identity with the exact target expansion \eqref{eq:local_target_expansion} and the first-order influence identity \eqref{eq:first_order_influence_identity} gives \eqref{eq:exact_local_target_decomposition}.
\end{proof}

Define the leading second-order expression
\begin{equation}
\label{eq:leading_second_order_term}
\mathcal B_{NT}
=
\{\mathbb G_0-\widehat{\mathbb G}\}
(\bm h^{(1)},\bm h^*)
+
\mathcal N(\bm h^{(1)},\bm h^*)
+
\mathcal Q(\bm h^{(1)}).
\end{equation}

\begin{proposition}[Decorrelation of the bilinear remainders]
\label{prop:bilinear_decorrelation}
Under the assumptions of \Cref{thm:target_expansion}, in the aligned local coordinates of \Cref{app:recovery},
\begin{align}
&\{\mathbb G_0-\widehat{\mathbb G}\}
(\widehat{\bm h},\bm h^*)
-
\{\mathbb G_0-\widehat{\mathbb G}\}
(\bm h^{(1)},\bm h^*)
\nonumber\\
&\hspace{30mm}
=o_p(s_{NT})+o_p(N^{-1}+T^{-1}),
\label{eq:gram_bilinear_replacement}\\
&\mathcal N(\widehat{\bm h},\bm h^*)
-
\mathcal N(\bm h^{(1)},\bm h^*)
=
o_p(s_{NT})+o_p(N^{-1}+T^{-1}),
\label{eq:nonlinear_bilinear_replacement}\\
&\mathcal Q(\widehat{\bm h})
-
\mathcal Q(\bm h^{(1)})
=
o_p(s_{NT})+o_p(N^{-1}+T^{-1}).
\label{eq:curvature_bilinear_replacement}
\end{align}
For the mixed loading--factor part of \eqref{eq:leading_second_order_term},
\begin{equation}
\label{eq:mixed_bilinear_decorrelation}
\Ezero(\mathcal B_{NT}^{\lambda f}\mid\mathscr C)
=
o(N^{-1}+T^{-1}),
\qquad
\operatorname{Var}_0(
\mathcal B_{NT}^{\lambda f}
\mid\mathscr C)
=
o(s_{NT}^2).
\end{equation}
Consequently,
$$
\mathcal B_{NT}^{\lambda f}
=
o_p(s_{NT})
+
o_p(N^{-1}+T^{-1}).
$$
\end{proposition}

\begin{proof}
Let $\delta_{NT}$ denote the bracketed dependence-tail expression in
\eqref{eq:local_deletion_tail}. By construction of the proof-only radii,
$\sqrt{NT}\,\delta_{NT}\to0$. Define
$$
\epsilon_{NT}
=
b_{NT}^2\log(NT)
\left(
N^{-1/2}+T^{-1/2}
\right)
+
\sqrt{NT}\,\delta_{NT}.
$$
The panel-growth condition and \eqref{eq:local_deletion_tail} imply
$\epsilon_{NT}\to0$.

We first prove the three replacement statements. Expand the empirical Gram
difference and the nonlinear term block by block in loading and factor
coordinates. Each summand is a product of a local regressor or score
contribution and one row of $\widehat{\bm h}-\bm h^{(1)}$, multiplied by a
row of the fixed target coordinate $\bm h^*$. For the summand anchored at
$a=(i,t)$, add and subtract the corresponding coordinate from the estimator
that deletes $\mathcal B_a$. The average squared difference between the
full and deleted estimator is $o_p(1)$ by
\eqref{eq:h_deletion_average}; the uniform row bounds in
\Cref{app:recovery} and Cauchy--Schwarz therefore make the contribution of
this replacement bounded by the first term in $\epsilon_{NT}$ times
$s_{NT}+N^{-1}+T^{-1}$.

After deletion, localize the regressor and score contribution to radii
$\lfloor q_T/3\rfloor$ and $\lfloor r_N/3\rfloor$. The localized
contribution attached to $a$ is measurable inside $\mathcal B_a$, while
the localized deleted coordinate is measurable outside that neighborhood.
The conditional NED covariance inequality used in the proof of
\Cref{prop:leave_neighborhood} therefore bounds the conditional mean by
$C\delta_{NT}$ times the relevant conditional second moments. The centered
covariances are split according to the distance between the two local
innovation neighborhoods. Exponential mixing, polynomial spatial
neighborhood growth, the same-side and cross-side inverse-Gram row bounds
in \eqref{eq:inverse_gram_rows}, and
$\|\bm h^*\|\asymp\|\calP_0\bD\|$ then make the sum of the centered
covariances $O(\epsilon_{NT}^2s_{NT}^2)$. The localization errors satisfy
the same order. This proves
\eqref{eq:gram_bilinear_replacement} and
\eqref{eq:nonlinear_bilinear_replacement}. Finally,
$\mathcal Q$ is bilinear in one loading and one factor coordinate.
Expanding
$\mathcal Q(\widehat{\bm h})-\mathcal Q(\bm h^{(1)})$ and using the
uniform row recovery rates with the preceding deletion bound gives
\eqref{eq:curvature_bilinear_replacement}.

It remains to verify that the mixed loading--factor expectation is smaller
than the two one-sided biases. The primitive normalized score blocks satisfy
a sharper cross-side covariance bound than their same-side variances. For a
loading score at unit $i$ and a factor score at time $t$, sequential
exogeneity eliminates every cross-time term in their product. Indeed, if
$s<t$, the time-$t$ disturbance has conditional mean zero given the
time-$s$ disturbance and the intervening past; if $s>t$, the time-$s$
disturbance has conditional mean zero given the time-$t$ disturbance and
the intervening past. Hence only $s=t$ remains, and
\begin{align}
&\left\|
\Ezero\!\left[
\bm s_i^\lambda(\bm s_t^f)'
\middle|\mathscr C
\right]
\right\|
\nonumber\\
&\quad\le
\frac{C}{\sqrt{NT}}
\sum_j
\left|
\Ezero(u_{it}u_{jt}\mid
x_t,\mathscr F_{t-1},\mathscr C)
\right|
\le
\frac{C}{\sqrt{NT}},
\label{eq:primitive_cross_score_covariance}
\end{align}
where the last inequality uses the conditional covariance row-sum bound and
the moment bounds for the regressors and factor/loadings.

Write $\mathbb G_0^{-1}$ in loading--loading, loading--factor,
factor--loading, and factor--factor blocks and substitute
$\bm h^{(1)}=\mathbb G_0^{-1}\bm s$. The diagonal loading and factor blocks
are bounded, while every off-diagonal same-side row and every cross-side
row has squared norm $O(N^{-1})$ or $O(T^{-1})$ by
\eqref{eq:inverse_gram_rows}. Combining these row bounds with
\eqref{eq:primitive_cross_score_covariance} and conditional covariance
summability gives, uniformly in $i,t$,
$$
\left\|
\Ezero\!\left[
h_i^{(1),\lambda}(h_t^{(1),f})'
\middle|\mathscr C
\right]
\right\|
\le
C\left[
(NT)^{-1/2}
+
N^{-1/2}
+
T^{-1/2}
\right].
$$
Every mixed term in \eqref{eq:leading_second_order_term} has one additional
factor $(NT)^{-1/2}$. The entry, mean, contrast, and fixed lag-sum
directions in \Cref{sec:model} have uniformly bounded total absolute
weight. Therefore,
$$
\left|
\Ezero(\mathcal B_{NT}^{\lambda f}\mid\mathscr C)
\right|
\le
C\left[
\frac1{NT}
+
\frac1{N\sqrt T}
+
\frac1{\sqrt N\,T}
\right]
=
o(N^{-1}+T^{-1}).
$$

For the centered mixed term, substitute the four inverse-Gram block
expansions just used, localize each score product, and again delete the
space--time neighborhood around one anchor. The covariance of two
localized products is bounded by the summable conditional mixing envelope;
the squared coefficients are summable at the rates in
\eqref{eq:inverse_gram_rows}. The resulting bound is
$$
\operatorname{Var}_0\!\left(
\mathcal B_{NT}^{\lambda f}
\middle|\mathscr C
\right)
\le
C\epsilon_{NT}^2s_{NT}^2
=
o(s_{NT}^2).
$$
This proves \eqref{eq:mixed_bilinear_decorrelation} and the proposition.
\end{proof}

\begin{proposition}[Orders and local form of the second-order terms]
\label{prop:second_order_bias}
Collect the terms in \eqref{eq:leading_second_order_term} according to whether the two first-order score coordinates are both loading coordinates, both factor coordinates, or one of each:
$$
\mathcal B_{NT}
=
\mathcal B_{NT}^{\lambda\lambda}
+
\mathcal B_{NT}^{ff}
+
\mathcal B_{NT}^{\lambda f}.
$$
Uniformly over the parameters of interest in Assumption~\ref{a:target},
\begin{align*}
\Ezero(\mathcal B_{NT}^{\lambda\lambda}\mid\mathscr C)
&=B_T/T,
& |B_T|&\le C,\\
\Ezero(\mathcal B_{NT}^{ff}\mid\mathscr C)
&=B_N/N,
& |B_N|&\le C,\\
\Ezero(\mathcal B_{NT}^{\lambda f}\mid\mathscr C)
&=o(N^{-1}+T^{-1}),
&
\operatorname{Var}_0(\mathcal B_{NT}^{\lambda f}\mid\mathscr C)
&=o(s_{NT}^2).
\end{align*}
The centered loading--loading and factor--factor terms have conditional variances $o(s_{NT}^2)$. Hence the mixed term is $o_p(s_{NT})+o_p(N^{-1}+T^{-1})$.

The two leading constants are linear in the target direction. To make the local contributions explicit, for the coordinate direction $\bm E_{M,it}$ that has a one in entry $(i,t)$ of coefficient matrix $M$ and zeros elsewhere, define
$$
b_{M,it}^{T}
=
T\,\Ezero\!\left\{
\mathcal B_{NT}^{\lambda\lambda}(\bm E_{M,it})
\middle|\mathscr C
\right\},
\qquad
b_{M,it}^{N}
=
N\,\Ezero\!\left\{
\mathcal B_{NT}^{ff}(\bm E_{M,it})
\middle|\mathscr C
\right\}.
$$
Linearity of $\bm h^*$ and of the target curvature in $\bD$ then gives the exact representation
\begin{equation}
\label{eq:bias_local_representation}
B_T
=
\sum_M\sum_{i,t}D_{M,it}b_{M,it}^{T},
\qquad
B_N
=
\sum_M\sum_{i,t}D_{M,it}b_{M,it}^{N}.
\end{equation}
The inverse-Gram row bounds and the conditional moment assumptions imply $\sup_{M,i,t}\{|b_{M,it}^{T}|+|b_{M,it}^{N}|\}\le C$ for the target classes in \Cref{sec:model}. For a target restricted to a set of rows or columns, the corresponding leading bias constant is obtained by restricting the same sums in \eqref{eq:bias_local_representation}.
\end{proposition}

\begin{proof}
Every first-order loading coordinate is a row of $\mathbb G_0^{-1}$ applied to a score normalized by $T^{-1/2}$; every first-order factor coordinate is the analogous object normalized by $N^{-1/2}$. The inverse-row bounds in \eqref{eq:inverse_gram_rows} and the conditional covariance summability in \eqref{eq:conditional_covariance_bound} imply
$$
\sup_i\Ezero\|h_i^{(1),\lambda}\|^2\le C,
\qquad
\sup_t\Ezero\|h_t^{(1),f}\|^2\le C.
$$
Every term with two loading coordinates carries the explicit factor $T^{-1}$ from the local parameterization; every term with two factor coordinates carries $N^{-1}$. Conditional H\"older and the $8+\eta$ moment bounds therefore give
$$
\left|\Ezero(\mathcal B_{NT}^{\lambda\lambda}\mid\mathscr C)\right|\le C/T,
\qquad
\left|\Ezero(\mathcal B_{NT}^{ff}\mid\mathscr C)\right|\le C/N,
$$
which defines bounded $B_T$ and $B_N$.

The mixed loading--factor term is handled by \Cref{prop:bilinear_decorrelation}. The key point is not the explicit prefactor $(NT)^{-1/2}$ by itself, but the additional cross-side decorrelation in \eqref{eq:primitive_cross_score_covariance} together with the same-side and cross-side inverse-Gram row decay. These yield an expectation $o(N^{-1}+T^{-1})$ and a centered variance $o(s_{NT}^2)$.

For a centered same-side term, replace its first-order scores by local approximations. The covariance of two products is bounded by a summable function of the distance between their space--time neighborhoods. Summing these covariance bounds and using the Riesz-weight rates in \Cref{prop:target_score} gives conditional variance $o(s_{NT}^2)$.

It remains to record the structure needed by the split-panel proof. The local Riesz coordinate $\bm h^*$ is linear in $\bD$, and every same-side conditional expectation in \eqref{eq:leading_second_order_term} is therefore a bounded linear functional of $\bD$. Evaluating these two functionals on the coordinate directions $\bm E_{M,it}$ gives the coefficients $b_{M,it}^{T}$ and $b_{M,it}^{N}$ defined above and proves \eqref{eq:bias_local_representation}. Because the functionals are finite-degree combinations of bounded conditional moments and blocks of $\mathbb G_0^{-1}$, the inverse-Gram bounds, conditional $8+\eta$ moments, and fixed ranks make the coefficients uniformly bounded. Restricting the target direction to selected rows or columns simply removes the corresponding terms from the exact linear sums.
\end{proof}

\begin{proof}[Proof of \Cref{thm:target_expansion}]
Lemma~\ref{lem:exact_target_decomposition} gives the exact expansion. The rate-level replacements \eqref{eq:gram_bilinear_replacement}--\eqref{eq:curvature_bilinear_replacement} in \Cref{prop:bilinear_decorrelation} show, for the specific bilinear objects appearing in that decomposition, that replacing $\widehat{\bm h}$ by $\bm h^{(1)}$ changes the last three terms by
$$
o_p(s_{NT})+o_p(N^{-1}+T^{-1}).
$$
Propositions~\ref{prop:bilinear_decorrelation} and \ref{prop:second_order_bias} then yield
$$
\widehat\phi-\phi(\bTheta_0)
=
\sum_{i,t}\Psi_{0,it}u_{it}
+
\frac{B_T}{T}
+
\frac{B_N}{N}
+
r_{NT},
$$
with the remainder stated in the theorem. For an individual coefficient and a fixed unit--time lag sum, the variance rates in \Cref{prop:target_score} make the two bias terms $o_p(s_{NT})$. For means and contrasts, the support-growth conditions stated in \Cref{thm:target_expansion} give the same conclusion. The conditional population limit and Slutsky yield the stated normal limits.
\end{proof}

\begin{proposition}[Balanced split-panel preservation]
\label{prop:split_preservation}
Consider a two-way mean or contrast covered by \Cref{thm:twoway}. Under the assumptions of \Cref{thm:target_expansion}, the independent balanced time and unit partitions in \Cref{subsec:split_correction} have the following properties, conditionally on $\mathscr C$, with probability approaching one over the random partitions.
\begin{enumerate}[label=(\roman*),leftmargin=2em]
\item The restricted target directions reassemble the full target exactly:
$$
\sum_{h=1}^{2}\bD^{T,h}=\bD,
\qquad
\sum_{g=1}^{2}\bD^{N,g}=\bD.
$$
\item Every positive rank and every lower Gram-matrix eigenvalue bound is preserved in the corresponding split panel.
\item Let $\pi_{T,h}=T_h/T$ and $\pi_{N,g}=N_g/N$. The leading constants for the restricted time targets satisfy
$$
B_T^{T,h}=\pi_{T,h}B_T+o_p(1),
\qquad
B_N^{T,h}=\pi_{T,h}B_N+o_p(1),
$$
and the restricted unit targets satisfy
$$
B_T^{N,g}=\pi_{N,g}B_T+o_p(1),
\qquad
B_N^{N,g}=\pi_{N,g}B_N+o_p(1).
$$
\item If $\Psi_{0,it}^{T}$ and $\Psi_{0,it}^{N}$ are the population split Riesz weights defined in \Cref{subsec:variance}, then
\begin{align}
&\sum_{h=1}^{2}\sum_{t\in\mathcal I_{T,h}}\sum_i
(\Psi_{0,it}^{T}-\Psi_{0,it})^2
+
\sum_{g=1}^{2}\sum_{i\in\mathcal I_{N,g}}\sum_t
(\Psi_{0,it}^{N}-\Psi_{0,it})^2
\nonumber\\
&\hspace{35mm}
=o_p\!\left(
\Ezero\!\left[\sum_{i,t}\Psi_{0,it}^2\middle|\mathscr C\right]
\right).
\label{eq:split_riesz_preservation}
\end{align}
Consequently,
\begin{equation}
\label{eq:corrected_variance_equivalence}
\frac{(s_{NT}^{\mathrm{corr}})^2}{s_{NT}^2}
\longrightarrow_p1.
\end{equation}
\end{enumerate}
\end{proposition}

\begin{proof}
We condition throughout on the realized common shocks and on the full panel, and take probability with respect to the independent split assignments. The conclusions are then integrated over the panel and the splits.

\emph{Step 1: the split targets reassemble the original parameter.}
By construction, each entry of $\bD$ is retained in exactly one time-restricted direction and exactly one unit-restricted direction. Therefore,
$$
\sum_h\ip{\bD^{T,h}}{\bTheta_0}
=
\ip{\bD}{\bTheta_0},
\qquad
\sum_g\ip{\bD^{N,g}}{\bTheta_0}
=
\ip{\bD}{\bTheta_0}.
$$
This argument applies without change to a proper time subset $R_T$, a proper unit subset $G_N$, and a signed contrast, because the original target weights and signs are retained rather than renormalized inside a split.

\emph{Step 2: ranks and Gram matrices are preserved.}
For a positive-rank matrix $\bM_0=\bU_M\bSigma_M\bV_M'$, let $v_{M,t}'$ be row $t$ of $\bV_M$. Conditional on $\bV_M$, sampling a balanced random time half without replacement gives
$$
\operatorname E_{\mathrm{split}}
\left(
\sum_{t\in\mathcal I_{T,h}}v_{M,t}v_{M,t}'
\right)
=
\pi_{T,h}\bm I.
$$
Incoherence bounds every summand by $C/T$. The finite-population variance formula therefore gives, entry by entry,
$$
\operatorname{Var}_{\mathrm{split}}
\left(
\sum_{t\in\mathcal I_{T,h}}v_{M,t,a}v_{M,t,b}
\right)
\le C/T.
$$
Since $r_M$ is fixed,
$$
\bV_{M,h}'\bV_{M,h}
=
\pi_{T,h}\bm I+o_p(1).
$$
Hence the restricted matrix retains rank $r_M$ and its nonzero singular values remain of order $\sqrt{NT_h}$. The same calculation with the rows of $\bU_M$ establishes rank preservation in a random unit half.

The loading and factor Gram matrices are finite-dimensional averages of terms with uniformly bounded second moments. Applying the same sampling-without-replacement variance calculation to each entry gives
$$
\max_h\|G_{\lambda,i}^{T,h}-G_{\lambda,i}\|_{\op}=o_p(1),
\qquad
\max_g\|G_{f,t}^{N,g}-G_{f,t}\|_{\op}=o_p(1),
$$
with the corresponding statements for the other Gram blocks. The lower eigenvalue bounds in \Cref{a:gram,a:identification} are therefore preserved.

\emph{Step 3: the leading bias constants are preserved.}
By \eqref{eq:bias_local_representation}, the loading- and factor-side constants are weighted sums of the explicitly defined local coefficients $b_{M,it}^{T}$ and $b_{M,it}^{N}$. These coefficients are finite-degree functions of: (a) the relevant blocks of the inverse local Gram operator; (b) the aligned loading and factor rows; and (c) a fixed-dimensional vector of conditional moments of the regressors and disturbances that enters the two same-side expectations in \Cref{prop:second_order_bias}. On the set where the Gram eigenvalues are bounded below by $c$, matrix inversion is Lipschitz. Consequently, if $b_{M,it}^{j,T,h}$ denotes the corresponding coefficient computed from time split $h$, there is a random envelope $\xi_{M,it}^{j}$ with $\sup_{M,i,t}\Ezero[(\xi_{M,it}^{j})^2\mid\mathscr C]\le C$ such that
\begin{equation}
\label{eq:bias_local_lipschitz}
\left|
b_{M,it}^{j,T,h}-b_{M,it}^{j}
\right|
\le
C\Delta_{T,h}\xi_{M,it}^{j},
\qquad
j\in\{T,N\},
\end{equation}
where $\Delta_{T,h}$ is the maximum of the operator-norm differences between the split and full aligned Gram blocks and the Euclidean differences between the corresponding fixed-dimensional conditional-moment vectors. Step~2 and the sampling-without-replacement variance bound give $\max_h\Delta_{T,h}=o_p(1)$. The unit-split coefficients satisfy the analogous bound with $\Delta_{N,g}=o_p(1)$. Thus the phrase ``preservation of the local bias contributions'' is a direct Lipschitz consequence of the objects entering their definition rather than an appeal to continuity alone.

For the time split, write
$$
a_t^{T}=\sum_M\sum_iD_{M,it}b_{M,it}^{T},
\qquad
a_t^{N}=\sum_M\sum_iD_{M,it}b_{M,it}^{N}.
$$
Then $B_T=\sum_ta_t^{T}$ and $B_N=\sum_ta_t^{N}$. For a two-way target, the largest time weight tends to zero; the uniform bounds in \Cref{prop:second_order_bias} therefore imply
$$
\sum_t(a_t^{T})^2+\sum_t(a_t^{N})^2=o_p(1).
$$
The finite-population variance formula gives, for $j\in\{T,N\}$,
$$
\operatorname E_{\mathrm{split}}
\left[
\left(
\sum_{t\in\mathcal I_{T,h}}a_t^{j}
-\pi_{T,h}\sum_ta_t^{j}
\right)^2
\middle|\mathscr C,\text{panel}
\right]
\le C\sum_t(a_t^{j})^2=o_p(1).
$$
The Lipschitz bound \eqref{eq:bias_local_lipschitz} and the bounded total absolute target weight also give
$$
B_j^{T,h}
=
\sum_{t\in\mathcal I_{T,h}}\sum_{M,i}
D_{M,it}b_{M,it}^{j,T,h}
=
\sum_{t\in\mathcal I_{T,h}}a_t^j+o_p(1)
=
\pi_{T,h}B_j+o_p(1),
\qquad j\in\{T,N\}.
$$
This proves the two time-split constant relations. For the unit split, set $c_i^{j}=\sum_M\sum_tD_{M,it}b_{M,it}^{j}$ and repeat the argument. Stratification within every group entering a contrast preserves the signed group weights, and the vanishing maximum unit weight gives $\sum_i(c_i^{j})^2=o_p(1)$. The unit-split analogue of \eqref{eq:bias_local_lipschitz} then yields $B_j^{N,g}=\pi_{N,g}B_j+o_p(1)$.

For a time split, the loading-side bias of subtarget $h$ is $B_T^{T,h}/T_h$. Since $\pi_{T,h}/T_h=1/T$, summing over the two halves gives $2B_T/T+o_p(T^{-1})$. Its factor-side biases sum to $B_N/N+o_p(N^{-1})$. The unit-split calculation gives $B_T/T+2B_N/N+o_p(N^{-1}+T^{-1})$.

\emph{Step 4: split Riesz weights and the corrected-score scale.}
Let $\Gamma_{T,h}$ be the population Gram form on time split $h$, and let $\bq_0^{T,h}$ solve its Riesz equation for $\bD^{T,h}$. After aligning the split and full singular spaces, let $\bq_0^{[T,h]}$ denote the restriction of the full-panel Riesz coordinates to the loading rows and factor rows used by time split $h$. Define $\Gamma_{N,g}$, $\bq_0^{N,g}$, and $\bq_0^{[N,g]}$ analogously.

We now make the factor-of-$\pi_{T,h}$ bookkeeping explicit. Let $\dot\phi(\bDelta)=\ip{\bD}{\bDelta}$ and $\dot\phi_{T,h}(\bDelta)=\ip{\bD^{T,h}}{\bDelta}$. The same sampling-without-replacement calculation as in Step~2, now applied to every fixed-dimensional Gram block and to the target derivative, yields
\begin{align}
\sup_{\|\bDelta\|=1}
\left|
\pi_{T,h}^{-1}
\Gamma_{T,h}(\bq_0^{[T,h]},\bDelta)
-
\Gamma_0(\bq_0,\bDelta)
\right|
&=
o_p(\|\calP_0\bD\|),
\label{eq:split_gram_scaled}\\
\sup_{\|\bDelta\|=1}
\left|
\pi_{T,h}^{-1}
\dot\phi_{T,h}(\bDelta)
-
\dot\phi(\bDelta)
\right|
&=
o_p(\|\calP_0\bD\|).
\label{eq:split_target_scaled}
\end{align}
For means and contrasts, the second display follows from the vanishing maximum time weight; for a full mean it reduces to the ordinary finite-population variance formula. Since the full Riesz equation gives $\Gamma_0(\bq_0,\bDelta)=\dot\phi(\bDelta)$, subtracting \eqref{eq:split_target_scaled} from \eqref{eq:split_gram_scaled} and multiplying by $\pi_{T,h}$ gives
$$
\sup_{\|\bDelta\|=1}
\left|
\Gamma_{T,h}(\bq_0^{[T,h]},\bDelta)
-
\ip{\bD^{T,h}}{\bDelta}
\right|
=
o_p(\|\calP_0\bD\|).
$$
The same argument with $\pi_{N,g}$ proves the unit-split approximate Riesz equation. Set $\bm d_{T,h}=\bq_0^{T,h}-\bq_0^{[T,h]}$. Subtracting the two Riesz equations gives
$$
\Gamma_{T,h}(\bm d_{T,h},\bDelta)
=
-\left\{
\Gamma_{T,h}(\bq_0^{[T,h]},\bDelta)
-
\ip{\bD^{T,h}}{\bDelta}
\right\}
$$
for every split tangent direction $\bDelta$. Taking $\bDelta=\bm d_{T,h}$ and using the split-panel Gram lower bound yields
$$
\Gamma_{T,h}(\bm d_{T,h},\bm d_{T,h})
\le
C
\sup_{\|\bDelta\|=1}
\left|
\Gamma_{T,h}(\bq_0^{[T,h]},\bDelta)
-
\ip{\bD^{T,h}}{\bDelta}
\right|^2
=
o_p(\|\calP_0\bD\|^2).
$$
Since this Gram norm is the squared prediction norm of the split Riesz-weight difference, summing over $h$ gives the first term in \eqref{eq:split_riesz_preservation}. Repeating the same argument with $\bm d_{N,g}=\bq_0^{N,g}-\bq_0^{[N,g]}$ gives the second term.

Let
$$
e_{it}^{T}=\Psi_{0,it}^{T}-\Psi_{0,it},
\qquad
 e_{it}^{N}=\Psi_{0,it}^{N}-\Psi_{0,it}.
$$
The corrected population weight equals
$$
3\Psi_{0,it}
-\Psi_{0,it}^{T}
-\Psi_{0,it}^{N}
=
\Psi_{0,it}-e_{it}^{T}-e_{it}^{N}.
$$
The conditional covariance row-sum bound in \Cref{a:moments} turns \eqref{eq:split_riesz_preservation} into an $o_p(s_{NT}^2)$ bound for the variance of the score difference. The target-specific lower variance bound in \Cref{a:target} gives $s_{NT}^2\asymp\Ezero[\sum_{i,t}\Psi_{0,it}^2\mid\mathscr C]$. Expanding the variance of the last display and applying conditional Cauchy--Schwarz proves \eqref{eq:corrected_variance_equivalence}.
\end{proof}

For the split-panel estimator, define the population subtargets
$$
\phi_{T,h}(\bTheta_0)
=
\ip{\bD^{T,h}}{\bTheta_0},
\qquad
\phi_{N,g}(\bTheta_0)
=
\ip{\bD^{N,g}}{\bTheta_0}.
$$
Applying \eqref{eq:target_expansion} to each split fit gives
\begin{align*}
\widehat\phi_{T,h}-\phi_{T,h}(\bTheta_0)
&=
S_{T,h}
+
\frac{B_T^{T,h}}{T_h}
+
\frac{B_N^{T,h}}{N}
+
 r_{T,h},\\
\widehat\phi_{N,g}-\phi_{N,g}(\bTheta_0)
&=
S_{N,g}
+
\frac{B_T^{N,g}}{T}
+
\frac{B_N^{N,g}}{N_g}
+
 r_{N,g}.
\end{align*}
Here the $S$ terms are the corresponding population Riesz scores and the remainders have the same target-specific order as in \Cref{thm:target_expansion}. Summing the two time equations and using \Cref{prop:split_preservation} gives
$$
\overline\phi_T-\phi(\bTheta_0)
=
\overline S_T
+
\frac{2B_T}{T}
+
\frac{B_N}{N}
+
o_p(s_{NT})+o_p(N^{-1}+T^{-1}).
$$
Similarly,
$$
\overline\phi_N-\phi(\bTheta_0)
=
\overline S_N
+
\frac{B_T}{T}
+
\frac{2B_N}{N}
+
o_p(s_{NT})+o_p(N^{-1}+T^{-1}).
$$
Therefore,
\begin{align*}
\widehat\phi^{\mathrm{corr}}-\phi(\bTheta_0)
={}&
3\left(
\frac{B_T}T+\frac{B_N}N
\right)
-
\left(
\frac{2B_T}T+\frac{B_N}N
\right)
\nonumber\\
&-
\left(
\frac{B_T}T+\frac{2B_N}N
\right)
+
\sum_{i,t}Z_{it}^{\mathrm{corr}}
+
o_p(s_{NT})
\nonumber\\
={}&
\sum_{i,t}Z_{it}^{\mathrm{corr}}
+
o_p(s_{NT}^{\mathrm{corr}}),
\end{align*}
where the final equality uses \eqref{eq:corrected_variance_equivalence}. The vector consisting of the full-panel and split-panel scores is a finite collection of conditionally NED random-field sums. By \Cref{lem:conditional_limit_transfer,prop:target_score}, the conditional random-field central limit theorem applies jointly. The conclusions hold for almost every independent balanced partition sequence.

\subsection{Spatial HAC variance estimation}
\label{app:variance}

This subsection completes the proof of feasible inference. It first shows that cross-time covariance is zero because the time-level Riesz sums are martingale differences. It then shows that the spatial weighting rule captures all nonnegligible same-time covariance. Finally, it replaces the population Riesz scores by the fitted Riesz weights and residuals.

\begin{proposition}[Consistency of the spatial variance estimator]
\label{prop:spatial_consistency}
Under the assumptions of \Cref{thm:twoway}, the distance cutoff $h_N=\lceil c_h\log(NT)\rceil$ with sufficiently large $c_h$ makes the omitted conditional spatial covariance negligible. Together with the recovery conclusions in \Cref{thm:recovery}, the target expansion in \Cref{thm:target_expansion}, and split preservation in \Cref{prop:split_preservation}, this implies consistency of \eqref{eq:spatial_variance} and \eqref{eq:corrected_spatial_variance}.
\end{proposition}

\begin{proof}
The conditional covariance inequality for NED fields bounds the covariance between two population Riesz-score contributions at spatial distance $a$ by
$$
C\{\psi_S(\lfloor a/3\rfloor)+\bar\alpha(\lfloor a/3\rfloor)^{\eta/(8+\eta)}\}
$$
times their conditional $2+\eta/4$ norms. Spatial shell growth is polynomial by \Cref{a:geometry}, whereas both dependence terms decay exponentially. Summing over distances larger than $h_N$ gives $o(s_{NT}^2)$ for sufficiently large $c_h$.

For retained pairs, the fourth-moment bounds from the localization lemma and $v_N(h_N)\le C\{\log(NT)\}^{d_s}$ make the centered spatial quadratic form $o_p(s_{NT}^2)$ under \Cref{a:growth}. Finally, the Riesz and residual recovery rates imply that replacing the population Riesz score by the feasible residual Riesz score changes the quadratic form by $o_p(s_{NT}^2)$ through Cauchy--Schwarz and $\|\bm W_N\|_{\op}\le Cv_N(h_N)$.
\end{proof}

For either the uncorrected or corrected estimator, write $Z_{it}^{*}$ for the corresponding population Riesz score and $\widehat Z_{it}^{*}$ for its feasible version. The star is only a proof abbreviation that allows the same calculation to cover both estimators. Define
$$
\widetilde s_{\mathrm{sp}}^{2,*}
=
\sum_t\sum_{i,k}
w_{ik,N}Z_{it}^{*}Z_{kt}^{*}.
$$

\begin{proof}[Proof of \Cref{thm:twoway}]
The split-panel calculation above gives the corrected influence expansion. It remains to establish consistency of the feasible spatial variance.

Conditional sequential exogeneity gives
$$
\Ezero\left(
\sum_i Z_{it}^{*}
\middle|
\mathscr F_{t-1},\mathscr C
\right)
=
0.
$$
Hence, for $s<t$,
\begin{align*}
\Ezero\left[
\left(\sum_i Z_{it}^{*}\right)
\left(\sum_k Z_{ks}^{*}\right)
\middle|\mathscr C
\right]
&=
\Ezero\left[
\left(\sum_k Z_{ks}^{*}\right)
\Ezero\left(
\sum_i Z_{it}^{*}
\middle|
\mathscr F_{t-1},\mathscr C
\right)
\middle|\mathscr C
\right]\\
&=0.
\end{align*}
Therefore,
$$
(s_{NT}^{*})^2
=
\sum_t\sum_{i,k}
\Ezero(Z_{it}^{*}Z_{kt}^{*}\mid\mathscr C).
$$
The covariance-tail argument in the proof of \Cref{prop:spatial_consistency} implies
$$
\left|
\Ezero(\widetilde s_{\mathrm{sp}}^{2,*}\mid\mathscr C)
-
(s_{NT}^{*})^2
\right|
=
o\{(s_{NT}^{*})^2\}.
$$

For the centered statistic, expand
\begin{align*}
\operatorname{Var}_0(
\widetilde s_{\mathrm{sp}}^{2,*}
\mid\mathscr C)
=
\sum_t
\sum_{i,k,j,\ell}
w_{ik,N}w_{j\ell,N}
\operatorname{Cov}_0(
Z_{it}^{*}Z_{kt}^{*},
Z_{jt}^{*}Z_{\ell t}^{*}
\mid\mathscr C).
\end{align*}
Every unit has at most $v_N(h_N)$ partners with nonzero weight. Splitting the quadruples according to their minimum spatial separation and applying \eqref{eq:conditional_covariance_bound} gives
$$
\operatorname{Var}_0(
\widetilde s_{\mathrm{sp}}^{2,*}
\mid\mathscr C)
\le
C v_N(h_N)^3
\sum_{i,t}
\Ezero(|Z_{it}^{*}|^4\mid\mathscr C)
+
o\{(s_{NT}^{*})^4\}.
$$
Because $h_N$ is logarithmic, $v_N(h_N)^3$ contributes only a fixed power of $\log(NT)$. The fourth moments of local and two-way Riesz scores contain the additional polynomial factors $T^{-1}$ or $(NT)^{-1}$ established in \Cref{prop:target_score}; these polynomial gains dominate every fixed logarithmic power under \Cref{a:growth}. Thus, the exponent $d_s+6$ in the main growth condition governs the leading recovery and bias terms, while the larger logarithmic power appearing in this intermediate variance bound is asymptotically harmless. The fourth-moment argument in the proof of \Cref{prop:spatial_consistency} therefore makes the right side $o\{(s_{NT}^{*})^4\}$. Conditional Chebyshev yields
$$
\frac{
\widetilde s_{\mathrm{sp}}^{2,*}
}{
(s_{NT}^{*})^2
}
\to_p1.
$$

Let $e_{it}^{*}=\widehat Z_{it}^{*}-Z_{it}^{*}$. In vector notation,
$$
\widehat s_{\mathrm{sp}}^{2,*}
-
\widetilde s_{\mathrm{sp}}^{2,*}
=
\sum_t
\left[
2(Z_t^{*})'\bm W_N e_t^{*}
+
(e_t^{*})'\bm W_Ne_t^{*}
\right].
$$
Cauchy--Schwarz and positive semidefiniteness give
\begin{align*}
\left|
\widehat s_{\mathrm{sp}}^{2,*}
-
\widetilde s_{\mathrm{sp}}^{2,*}
\right|
\le
\|\bm W_N\|_{\op}
\left[
2
\left\{
\sum_{i,t}(Z_{it}^{*})^2
\sum_{i,t}(e_{it}^{*})^2
\right\}^{1/2}
+
\sum_{i,t}(e_{it}^{*})^2
\right].
\end{align*}
It remains to verify the feasible-score bound for the corrected statistic, which contains residuals from three distinct fitted objects. Write
$$
\widehat Z_{it}^{\mathrm{corr}}
=
3\widehat\Psi_{it}\widehat u_{it}
-
\widehat\Psi_{it}^{T}\widehat u_{it}^{T}
-
\widehat\Psi_{it}^{N}\widehat u_{it}^{N},
\qquad
Z_{it}^{\mathrm{corr}}
=
u_{it}
\left(
3\Psi_{0,it}-\Psi_{0,it}^{T}-\Psi_{0,it}^{N}
\right).
$$
For a time split $h$,
$$
\widehat u_{it}^{T,h}-u_{it}
=
-\mathcal Y_{it}
(\widehat\bTheta_{T,h}-\bTheta_0),
\qquad
t\in\mathcal I_{T,h},
$$
and the analogous identity holds for each unit split. By \Cref{prop:split_preservation}, every split preserves the supplied ranks and the Gram lower bounds. Applying the recovery argument of \Cref{thm:recovery} to the $N\times T_h$ and $N_g\times T$ panels therefore gives, uniformly over the two halves,
$$
\sum_{t\in\mathcal I_{T,h}}\sum_i
(\widehat u_{it}^{T,h}-u_{it})^2
=
O_p\{NT_h\zeta_{N,T_h}\},
\qquad
\sum_t\sum_{i\in\mathcal I_{N,g}}
(\widehat u_{it}^{N,g}-u_{it})^2
=
O_p\{N_gT\zeta_{N_g,T}\}.
$$
Balanced halves imply $\zeta_{N,T_h}\asymp\zeta_{NT}$ and $\zeta_{N_g,T}\asymp\zeta_{NT}$. The split Riesz recovery in \eqref{eq:split_riesz_preservation}, together with the full-panel Riesz rate, then yields
$$
\|\bm W_N\|_{\op}
\sum_{i,t}
(\widehat Z_{it}^{\mathrm{corr}}-Z_{it}^{\mathrm{corr}})^2
=
o_p\{(s_{NT}^{\mathrm{corr}})^2\}.
$$
Substituting this bound for $\sum(e_{it}^{*})^2$ in the preceding Cauchy--Schwarz display makes
$\widehat s_{\mathrm{corr}}^2-\widetilde s_{\mathrm{sp}}^{2,\mathrm{corr}}
=o_p\{(s_{NT}^{\mathrm{corr}})^2\}$. For the uncorrected statistic the same calculation uses only the full-panel recovery rates. Studentization follows from the corresponding influence expansion, the population limit, and conditional Slutsky.
\end{proof}

\subsection{Rank selection}
\label{app:rank}

This subsection proves \Cref{cor:rank_pilot_rate,thm:rank_consistency}. No information criterion or singular-value cutoff is used.

For the pilot proof, define
\begin{equation}
\label{eq:pilot_difference_class}
\calD_{\max}^{+}
=
\left\{
\bTheta-\bTheta_0:
\|\bTheta\|_{\max}\le B,
\ \rank(\bM)\le\bar r_M+1\ \text{for every }M
\right\}.
\end{equation}
For $\bDelta\in\calD_{\max}^{+}$, rank subadditivity gives $\rank(\bDelta_M)\le\bar r_M+1+r_{M,0}\le2\bar r_M+1$. These bounds and the number of coefficient matrices are fixed.

\begin{lemma}[Regularity of the spectral scale weights]
\label{lem:rank_scale_weights}
Under \Cref{a:stab,a:geometry,a:ned,a:moments,a:pilot_identification,a:growth}, the finite collection of ratios $w_M/w_{A,1}$ is bounded above and away from zero in probability.
\end{lemma}

\begin{proof}
Let $z_{it,M}$ denote the regressor multiplying coefficient matrix $M$: $z_{it,A^{(\ell)}}=y_{i,t-\ell}$, $z_{it,B^{(k)}}=x_{it,k}$, and $z_{it,H}=1$. The conditional moment and localization results in \Cref{lem:localization,lem:weighted_concentration} imply, for each of the finitely many nonconstant regressors,
\[
\frac1{NT}\sum_{i,t}z_{it,M}^2
-
\frac1{NT}\sum_{i,t}\Ezero(z_{it,M}^2\mid\mathscr C)
=o_p(1).
\]
The conditional averages on the right are uniformly bounded above by the maintained moment conditions.

For the lower bound, fix an observation $(i,t)$ and a nonconstant regressor block $M$. Choose a scalar $d\ne0$ small enough that adding $d$ to entry $(i,t)$ of the true coefficient matrix remains inside the box. Such a one-entry change raises matrix rank by at most one, so the perturbed coefficient collection belongs to the pilot class. Assumption~\ref{a:pilot_identification} applied to this perturbation gives
\[
d^2\Ezero(z_{it,M}^2\mid\mathscr C)\ge c_+d^2.
\]
Thus every relevant conditional second moment is bounded below by $c_+$. Since $P\ge1$, the reference weight $w_{A,1}$ is available; $w_H=1$ by definition. The finite-collection argument and the sample-moment approximation above prove that all $w_M/w_{A,1}$ are bounded above and away from zero in probability.
\end{proof}

\begin{proposition}[Uniform concentration on the spectral-pilot class]
\label{prop:uniform_concentration_pilot}
Under \Cref{a:stab,a:exog,a:geometry,a:ned,a:moments,a:signal,a:identification,a:gram,a:growth} and the additional pilot identification condition,
\begin{equation}
\label{eq:pilot_score_bound}
\sup_{\bDelta\in\calD_{\max}^{+}\setminus\{0\}}
\frac{|\sum_{i,t}u_{it}\mathcal Y_{it}(\bDelta)|}{\|\bDelta\|}
=O_p\{\sqrt{NT\zeta_{NT}}\},
\end{equation}
and, uniformly over $\calD_{\max}^{+}$,
\begin{equation}
\label{eq:pilot_curvature_bound}
\sum_{i,t}\mathcal Y_{it}(\bDelta)^2
\ge
\frac{c_+}{2}\|\bDelta\|^2-CNT\zeta_{NT}
\end{equation}
for the fixed $c_+>0$ in Assumption~\ref{a:pilot_identification}.
\end{proposition}

\begin{proof}
For a coefficient matrix $M$, let $S_M$ be the $N\times T$ score matrix with entries $u_{it}z_{it,M}$. The matrix-score part of \Cref{prop:uniform_concentration} is obtained before restricting the rank of $\bDelta$ and therefore gives
\[
\max_M\|S_M\|_{\op}=O_p\{\sqrt{NT\zeta_{NT}}\}.
\]
For every $\bDelta\in\calD_{\max}^{+}$, rank subadditivity gives $\rank(\bDelta_M)\le2\bar r_M+1$. Hence the standard operator--Frobenius norm inequality for a matrix of fixed rank implies
\[
\left|\sum_{i,t}u_{it}\mathcal Y_{it}(\bDelta)\right|
\le
\sum_M\|S_M\|_{\op}\sqrt{\rank(\bDelta_M)}\|\bDelta_M\|_F
\le C\max_M\|S_M\|_{\op}\|\bDelta\|,
\]
where $C$ depends only on the fixed pilot ranks and number of coefficient matrices. This proves \eqref{eq:pilot_score_bound}.

For the prediction term, write $n=NT$, $b=b_{NT}$, and choose $r_n^2=C_0n\zeta_{NT}$ with $C_0$ fixed and sufficiently large. On the truncation event used in the proof of \Cref{prop:uniform_concentration}, the centered regressor products are bounded by $Cb^2$ and satisfy the same exponential local-approximation bounds. Peel the set $\{\bDelta\in\calD_{\max}^{+}:\|\bDelta\|\ge r_n\}$ into dyadic Frobenius-norm shells. On each normalized shell, the singular-vector parameterization for matrices of ranks at most $2\bar r_M+1$ gives an $\epsilon$-net with
\[
\log |\mathcal N_\epsilon|
\le
C_{\mathrm{pil}}(N+T)\log(C/\epsilon),
\]
where $C_{\mathrm{pil}}$ depends only on the fixed pilot ranks. Taking $\epsilon=c_0/b^2$, applying \Cref{lem:weighted_concentration} simultaneously to the finitely many block pairs and net points, and choosing $C_0$ large enough to dominate the net entropy and the $O(\log n)$ shells yields
\[
\sup_{\substack{\bDelta\in\calD_{\max}^{+}\\\|\bDelta\|\ge r_n}}
\frac{\left|\sum_{i,t}\{\mathcal Y_{it}(\bDelta)^2-
\Ezero[\mathcal Y_{it}(\bDelta)^2\mid\mathscr C]\}\right|}{\|\bDelta\|^2}
\le
\frac{c_+}{2}
\]
with conditional probability approaching one. The Lipschitz extension from the net uses the same $Cb^2$ bound as in \Cref{prop:uniform_concentration}. Assumption~\ref{a:pilot_identification} then gives
\[
\sum_{i,t}\mathcal Y_{it}(\bDelta)^2
\ge
\frac{c_+}{2}\|\bDelta\|^2
\qquad
(\|\bDelta\|\ge r_n).
\]
If $\|\bDelta\|<r_n$, the empirical prediction norm is nonnegative, and subtracting a sufficiently large multiple of $n\zeta_{NT}$ makes \eqref{eq:pilot_curvature_bound} automatic. Combining the two regions proves the result. Enlarging each fixed rank cap by one changes only $C_{\mathrm{pil}}$; all $N,T$ rates and the rectangular growth condition are unchanged.
\end{proof}

\begin{proof}[Proof of \Cref{cor:rank_pilot_rate}]
Let $\widehat\bDelta^{\mathrm{pil}}=\widehat\bTheta^{\mathrm{pil}}-\bTheta_0$ and let $\delta_{NT}$ denote the objective gap in \eqref{eq:rank_pilot_objective_gap}. The truth is feasible for the spectral pilot. Expanding the least-squares objective around $\bTheta_0$ and using approximate optimality gives
\begin{equation}
\label{eq:pilot_basic_inequality}
\frac12\sum_{i,t}\mathcal Y_{it}(\widehat\bDelta^{\mathrm{pil}})^2
\le
\sum_{i,t}u_{it}\mathcal Y_{it}(\widehat\bDelta^{\mathrm{pil}})
+NT\delta_{NT}.
\end{equation}
By \Cref{prop:uniform_concentration_pilot}, with probability approaching one,
\[
\frac{c_+}{4}\|\widehat\bDelta^{\mathrm{pil}}\|^2
\le
C\sqrt{NT\zeta_{NT}}\|\widehat\bDelta^{\mathrm{pil}}\|
+CNT\zeta_{NT}+NT\delta_{NT}.
\]
Since $\delta_{NT}=o_p(\zeta_{NT})$, solving this quadratic inequality yields $\|\widehat\bDelta^{\mathrm{pil}}\|=O_p\{\sqrt{NT\zeta_{NT}}\}$. Each coefficient-matrix operator norm is bounded by the joint Frobenius norm, proving \eqref{eq:pilot_operator_rate}.
\end{proof}

\begin{lemma}[Normalized pilot spectra]
\label{lem:rank_spectral_separation}
Under the assumptions of \Cref{thm:rank_consistency}, uniformly over the fixed coefficient-matrix collection and the finitely many indices used by the selector,
\[
\widehat\lambda_{M,j}=\Theta_p(1)\quad(j\le r_{M,0}),
\qquad
\widehat\lambda_{M,j}=O_p(\zeta_{NT})\quad(j>r_{M,0}).
\]
For positive-rank indices the first statement includes a lower bound away from zero in probability.
\end{lemma}

\begin{proof}
By Weyl's inequality and \eqref{eq:pilot_operator_rate}, $|\sigma_j(\widehat\bM^{\mathrm{pil}})-\sigma_j(\bM_0)|=O_p\{\sqrt{NT\zeta_{NT}}\}$. If $j\le r_{M,0}$, \Cref{a:signal} places $\sigma_j(\bM_0)$ between fixed positive multiples of $\sqrt{NT}$; if $j>r_{M,0}$, the population singular value is zero. Squaring, dividing by $NT$, and applying \Cref{lem:rank_scale_weights} gives the stated orders.
\end{proof}

\begin{lemma}[Rectangular ridge separation]
\label{lem:rank_ridge_rate}
Under \Cref{a:growth}, $a_{NT}=1/\log(NT)$ satisfies $a_{NT}\to0$ and $\zeta_{NT}/a_{NT}\to0$, without an aspect-ratio restriction.
\end{lemma}

\begin{proof}
Let $G_{NT}=(NT)^{4/(8+\eta)}\{\log(NT)\}^{d_s+6}(1/N+1/T)$, which converges to zero by \Cref{a:growth}. Since
\[
\frac{\zeta_{NT}}{a_{NT}}
=
\frac{G_{NT}}{(NT)^{2/(8+\eta)}\log(NT)},
\]
the ratio converges to zero because the denominator diverges. Also $a_{NT}\to0$ as $NT\to\infty$. The calculation retains $1/N+1/T$ and therefore does not use $N\asymp T$.
\end{proof}

\begin{proof}[Proof of \Cref{thm:rank_consistency}]
Fix a coefficient matrix $M$ and write $r=r_{M,0}$. First suppose $1\le r\le\bar r_M$. For every $1\le j<r$, both adjacent normalized singular values are bounded above and away from zero by \Cref{lem:rank_spectral_separation}; hence $R_{M,j}$ is bounded away from zero in probability. The anchored ratio $R_{M,0}$ is also bounded away from zero because $\widehat\lambda_{M,1}$ is bounded away from zero in probability. At the true rank,
\[
R_{M,r}
=
\frac{O_p(\zeta_{NT})+a_{NT}}{\Theta_p(1)+a_{NT}}
\longrightarrow_p0
\]
by \Cref{lem:rank_ridge_rate}. For $j>r$, both adjacent spectral values are $o_p(a_{NT})$, so $R_{M,j}\to_p1$. When $r=\bar r_M$, the numerator at the true rank uses the permitted $(\bar r_M+1)$st pilot singular value; the zero limit is therefore a consequence of spectral separation rather than rank truncation.

If $r=0$, every pilot spectral value is $O_p(\zeta_{NT})$. Therefore $R_{M,0}=\{a_{NT}+o_p(a_{NT})\}/(1+a_{NT})\to_p0$, while $R_{M,j}\to_p1$ for every $j\ge1$. Thus the true rank is the unique asymptotic minimizer in all cases, and $\Pzero(\widehat r_M=r_{M,0}\mid\mathscr C)\to_p1$. Because the number of matrices and all caps are fixed,
\[
\Pzero(\widehat\br\ne\br_0\mid\mathscr C)
\le
\sum_M\Pzero(\widehat r_M\ne r_{M,0}\mid\mathscr C)
\longrightarrow_p0.
\]
On $\{\widehat\br=\br_0\}$, the final post-selection estimator is exactly the supplied-rank estimator analyzed in \Cref{thm:recovery,thm:target_expansion,thm:twoway}. For any fixed target, the difference between the distribution functions of the selected- and supplied-rank statistics is bounded by $\Pzero(\widehat\br\ne\br_0\mid\mathscr C)$ plus the supplied-rank approximation error. Hence all stated recovery and inference results transfer.
\end{proof}

\section{Additional Monte Carlo Results}
\label{app:simulation}
\subsection{Rank-stress data-generating processes}
\label{app:rank-stress-dgps}

We study the true-rank vectors $(1,1,1)$, $(2,1,1)$, and
$(1,0,2)$.  All rank-stress coefficient matrices are first generated over
the complete burn-in and observed horizon of length $L=T+50$, after which
the final $T$ columns are retained.  The coefficients multiplying the added loading--factor terms below are one.
All added loading and factor draws are mutually independent and independent
of the baseline DGP primitives.

For the rank-one design, we use the main-DGP matrices
\begin{align*}
 A_{it}^{(1),\mathrm{raw}}&=\lambda_{a,i}f_{a,t}, &
 B_{it}^{(1)}&=\lambda_{b,i}f_{b,t}, &
 H_{it}^{(1),\mathrm{raw}}&=\lambda_{h,i}g_{h,t},
\end{align*}
where $f_{a,t}=0.5+0.1g_{a,t}$,
$f_{b,t}=0.6+0.15g_{b,t}$, and
\[
 g_{m,t}=0.5g_{m,t-1}+\sqrt{0.75}\,v_{m,t},\qquad
 v_{m,t}\stackrel{\mathrm{iid}}{\sim}U[-\sqrt{3},\sqrt{3}],
 \qquad g_{m,-50}=0,
\]
for $m\in\{a,b,h\}$.  Also
$\lambda_{h,i}\stackrel{\mathrm{iid}}{\sim}U[-\sqrt{3},\sqrt{3}]$,
while $\lambda_{a,i}$ and $\lambda_{b,i}$ are the loadings from the
corresponding main DGP.  No rank-stress rescaling is applied when the
requested rank is one.  We set
\[
 c_a=\min\left\{1,
 \frac{0.85}{\max_{i,t\leq L}|A_{it}^{(1),\mathrm{raw}}|}\right\},
 \qquad A^{(1)}=c_aA^{(1),\mathrm{raw}},
 \qquad H_0=c_\xi c_H H^{(1),\mathrm{raw}}.
\]
The resulting true-rank vector is $(1,1,1)$ almost surely.

For the $(2,1,1)$ design, only the autoregressive-coefficient matrix is
expanded.  Independently draw
\[
 \widetilde\lambda_{A,i}\stackrel{\mathrm{iid}}{\sim}
 U[-\sqrt{3},\sqrt{3}],\qquad
 \widetilde f_{A,t}\stackrel{\mathrm{iid}}{\sim}
 U[-\sqrt{3},\sqrt{3}].
\]
Thus $\widetilde f_{A,t}$ is iid over time and does not follow an AR
recursion.  Define
\[
 E_{A,\mathrm{raw}}
 =L_{A,\max}(0.5+3\cdot0.1),\qquad
 s_A=\frac{E_{A,\mathrm{raw}}}{E_{A,\mathrm{raw}}+3},
\]
where $L_{A,\max}=1+0.1\sqrt{3}$ in DGPs 1--3 and
$L_{A,\max}=1.1+0.08\sqrt{3}$ in DGP 4.  The rank-two raw matrix is
\[
 A_{it}^{(2),\mathrm{raw}}
 =s_A\left(\lambda_{a,i}f_{a,t}
 +\widetilde\lambda_{A,i}\widetilde f_{A,t}\right).
\]
Both terms have coefficient one before the same deterministic common
rescaling.  The added term has support envelope three.  We then recompute
\[
 c_a=\min\left\{1,
 \frac{0.85}{\max_{i,t\leq L}|A_{it}^{(2),\mathrm{raw}}|}\right\},
 \qquad A^{(2)}=c_aA^{(2),\mathrm{raw}}.
\]
There is no entrywise clipping.  The rank-one $B$ and $H$ constructions
are unchanged, giving true ranks $(2,1,1)$ almost surely.

For the $(1,0,2)$ design, $A$ is the baseline rank-one matrix (including
its full-horizon stability scale), and
\[
 B_{it}\equiv0.
\]
Independently draw
\[
 \widetilde\lambda_{H,i}\stackrel{\mathrm{iid}}{\sim}
 U[-\sqrt{3},\sqrt{3}],\qquad
 \widetilde g_{H,t}\stackrel{\mathrm{iid}}{\sim}
 U[-\sqrt{3},\sqrt{3}],
\]
where the second factor is iid over time.  With
\[
 E_H=3\sqrt{3},\qquad s_H=\frac{3\sqrt{3}}{3\sqrt{3}+3},
\]
the raw rank-two interactive effect is
\[
 H_{it}^{(2),\mathrm{raw}}
 =s_H\left(\lambda_{h,i}g_{h,t}
 +\widetilde\lambda_{H,i}\widetilde g_{H,t}\right).
\]
Its population observed-entry variance is
\[
 V_{H,2}(T)=s_H^2\left\{
 \frac1T\sum_{k=51}^{50+T}(1-0.5^{2k})+1\right\},
\]
and the population-moment scale is
\[
 c_H=\left\{\frac{0.30}{0.70}\frac{1}{V_{H,2}(T)}\right\}^{1/2}.
\]
For this zero-slope design we impose the normalization $c_\xi=1$ and set
\texttt{r2\_scale\_identified=false}.  Accordingly, the pooled
$R^2$ is the value induced by this normalization rather than a calibrated
target.  The final interactive effect is
$H_0=c_HH^{(2),\mathrm{raw}}$, and the true-rank vector is $(1,0,2)$
almost surely.

For the rank-stress designs, the conditioning field is augmented by the additional bounded unit loadings and time-factor paths used to construct the second low-rank term. Conditional on this enlarged field, the rank-stress coefficient matrices are fixed and the remaining idiosyncratic innovations retain the same dependence structure as in the corresponding baseline DGP.

The population $c_H$ and common-random-number $c_\xi$ calibrations are
performed separately for each DGP, $(N,T)$ cell, and true-rank vector.
The calibration uses master seed 20260807 and 50 independent calibration
draws, none of which is reused as a reported Monte Carlo replication.

\subsection{Calibration and coefficient summaries}

\input{tables/mc/app_mc_calibration}

\subsection{Entrywise and fixed-time targets}

\input{tables/mc/app_mc_entry}

\input{tables/mc/app_mc_fixed_time}

\subsection{Two-way averages and bias correction}

\input{tables/mc/app_mc_bias}

\input{tables/mc/app_mc_halves}

\subsection{Rank-selection, spatial-weight, and computation diagnostics}

\placeholdertable{Blockwise ridge-ratio diagnostics}{tab:app_mc_rank_sensitivity}{Revision-10 shell reporting: DGP; $N$; $T$; coefficient block; true rank; selected rank; normalized pilot singular values through $\bar r_M+1$; $R_{M,0},\ldots,R_{M,\bar r_M}$; minimum ratio; second-smallest ratio; ratio gap; pilot numerical validity; pilot boundary activity; pilot objective, multistart-stability, and stationarity diagnostics; and final post-refit validity. Candidate counts, candidate coverage, information criteria, IC gaps, penalty or threshold multipliers, and IC completion are excluded. No numerical entries are populated in RR3.1.}


\input{tables/mc/app_mc_bandwidth}

\input{tables/mc/app_mc_runtime}

\input{tables/mc/app_mc_failures}

\input{tables/mc/app_mc_extreme}

\input{figures/mc/app_mc_bias_decomposition}

\input{figures/mc/app_mc_interval_length}

\input{figures/mc/app_mc_runtime}

\input{figures/mc/app_mc_tails}

%================================================
\section{Empirical Appendix}
\label{app:empirical}

\subsection{Data construction}

\textit{Placeholder.} Document sources, sample filters, seasonal adjustment, transformations, lag construction, geographic matching, missing-value treatment, and the definitions of all predetermined controls for the unemployment and housing panels.

\placeholdertable{Empirical sample construction}{tab:app_emp_samples}{Placeholder for raw and final sample sizes, time ranges, exclusions, coordinate matches, and control availability.}

\subsection{Rank selection and numerical diagnostics}

\placeholdertable{Empirical ridge-ratio rank selection and numerical diagnostics}{tab:app_emp_rank}{Placeholder for reportable and pilot caps, spectral pilot singular values, fitted-value scale weights, normalized block spectra, all ridge ratios, selected ranks, minimum and second-smallest ratio margins, pilot objective and stationarity diagnostics, boundary activity, final post-refit validity, and neighboring supplied-rank target sensitivity.}

\placeholdertable{Full- and split-panel numerical diagnostics}{tab:app_emp_numerics}{Placeholder for objective paths, first-order residuals, Riesz solver convergence, Gram-matrix eigenvalues, envelope checks, and split-panel stability.}

\subsection{Unemployment robustness}

\placeholdertable{Unemployment specifications and controls}{tab:app_emp_unemp_controls}{Placeholder for baseline and control-augmented specifications, alternative sample windows, group definitions, selected ranks, estimates, and spatial standard errors.}

\placeholdertable{Unemployment target details}{tab:app_emp_unemp_targets}{Placeholder for full means, fixed-time means, high and low group means, contrasts, selected entries, and the two-way correction terms.}

\placeholderfigure{Unemployment coefficient paths and maps}{fig:app_emp_unemp_figures}{Placeholder for time paths, entrywise estimates, geographic patterns, and common-shock diagnostics.}

\subsection{Housing robustness}

\placeholdertable{Housing specifications and controls}{tab:app_emp_house_controls}{Placeholder for lag choices, control-augmented specifications, alternative samples, selected ranks, estimates, and spatial standard errors.}

\placeholdertable{Housing target details}{tab:app_emp_house_targets}{Placeholder for lag means, lag-sum persistence, price-tier means, contrasts, selected entries, and the two-way correction terms.}

\placeholderfigure{Housing coefficient paths and maps}{fig:app_emp_house_figures}{Placeholder for time paths, price-tier contrasts, entrywise estimates, and geographic patterns.}

\subsection{Variance and dependence diagnostics}

\placeholdertable{Empirical spatial-weight and cluster sensitivity}{tab:app_emp_dependence}{Placeholder for alternative geographic distance cutoffs, known-cluster calculations, residual spatial covariance, and conditional variance diagnostics.}

\placeholdertable{Residual adequacy and homogeneous benchmarks}{tab:app_emp_fit}{Placeholder for fit measures, residual serial correlation, residual spatial correlation, remaining singular-value shares, and comparisons with homogeneous dynamic-panel benchmarks.}

\end{document}
